\subsection{Quartic CM fields at \texorpdfstring{$n=2$}{n=2}}
\label{ssec:quartic}

By \Cref{thm:closure}(ii) the only form of (P2) on a minimal route to the
conjecture is the one for the split families of the CM fields of degree at
least four, and by \Cref{prop:descent}(i) the Weil classes of any family of
such a field in a higher dimension give those of its families at $n=2$. For a
quartic field these are families of abelian eightfolds, the smallest
families of a field of degree at least four. This subsection determines how
far the available mechanisms reach there.

Divisor classes reach the Weil classes exactly on a Noether--Lefschetz locus,
which has codimension two when $\delta$ is trivial (\Cref{prop:quarticnl},
\Cref{cor:quarticbase}). For biquadratic $F$, Markman's theorem reaches them
on a further locus of dimension four (\Cref{prop:quarticscalar}). For
$F=\QQ(\zeta_{5}),\QQ(\zeta_{8}),\QQ(\zeta_{12})$ and trivial $\delta$,
Schoen's cycles on Prym varieties of cyclic covers reach them on a locus of
dimension at most six, at whose very general points the Hodge group is all of
$G_{F}$ (\Cref{rem:quarticprym}). The Casimir class of the quadratic form on
the weight two part $R_{F}$ below, the natural class of degree four attached
to it, is algebraic exactly when the Weil classes are
(\Cref{prop:quarticcasimir}), and correspondences with abelian varieties whose
simple factors have dimension at most seven cannot help at a member with full
Hodge group (\Cref{prop:quarticnogo}). For a general shape the polarised
criterion asks $\dim\Ext^{2}(E,E)=112$ of one complex
(\Cref{thm:quarticp2prime}). No Orlov product of sheaves on a principally
polarised fourfold with real multiplication meets the criterion through a
dimension count, whatever the quartic field (\Cref{thm:quarticobstruction}),
and no Orlov product of complexes does so unless both factors have nonzero self-Ext
groups in some degree at most $-3$ (\Cref{rem:quarticnegext}). Those products
are in any case candidates only for a twisted form of the criterion, the form
of Markman's reduction (\Cref{rem:quarticsecantshape}). Objects with pure
spinor characters give no Weil class at all (\Cref{prop:quarticexclusions}).
Together these results fix the loci that
divisors, scalar extension and correspondences reach inside $\cD_{F}$, and
the numbers that one complex has to attain on the rest.

Throughout, $F$ is a quartic CM field and $F_{0}$ its real quadratic subfield,
with real places $\tau_{1},\tau_{2}$, and $(V,H)$ is of $(F,2)$-Weil type with
discriminant $\delta$; so $m=n=2$, the members are abelian eightfolds, and
$\cD_{F}$ has dimension $mn^{2}=8$. By \Cref{prop:cmweil}(i),
$H^{1}(A,\CC)=\bigoplus_{\sigma}V_{\sigma}$ over the four embeddings $\sigma$ of
$F$, with $\dim V_{\sigma}^{1,0}=\dim V_{\sigma}^{0,1}=2$. Put
\[
  R_{F}=\textstyle\bigwedge^{2}_{F}V\;\subseteq\;\bigwedge^{2}_{\QQ}V=H^{2}(A,\QQ),
\]
the $\QQ$-subspace with complexification
$\bigoplus_{\sigma}\bigwedge^{2}V_{\sigma}$; for $m=1$ it is the space $R$ of
\Cref{prop:divfails}. It is an $F$-vector space of dimension $6$, of
$\QQ$-dimension $24$, with Hodge numbers $(4,16,4)$. Fix a generator $\Omega$
of the $F$-line $W_{F}=\bigwedge^{4}_{F}V$, so that
$W_{F}=\{a\Omega:a\in F\}$, and define a symmetric $F$-bilinear form $q$ on
$R_{F}$ by $x\wedge_{F}y=q(x,y)\,\Omega$. Let $h$ be the hermitian form on
$R_{F}$ induced by $H$, $h(x_{1}\wedge x_{2},y_{1}\wedge y_{2})
=\det\bigl(H(x_{i},y_{j})\bigr)$, and $h_{4}$ the one induced on
$\bigwedge^{4}_{F}V$; in a basis $e_{1},\dots,e_{4}$ with
$\Omega=e_{1}\wedge e_{2}\wedge e_{3}\wedge e_{4}$, $h_{4}(\Omega,\Omega)$ is
$\det H$, so its class in $F_{0}^{\times}/\Nm_{F/F_{0}}(F^{\times})$ is
$\delta$. With dual bases $(\omega_{k})$, $(\omega_{k}^{*})$ of $F$ for the
trace form, as in \Cref{thm:cmbasepoint}, put
$x\star y=\sum_{k}(\omega_{k}x)\cup(\omega_{k}^{*}y)$ for $x,y\in R_{F}$. Then
\begin{equation}\label{eq:quarticstar}
  x\star y=q(x,y)\,\Omega ,
\end{equation}
because the complexification of $(\omega_{k}x)\cup(\omega_{k}^{*}y)$ is
$\sum_{\sigma,\sigma'}\sigma(\omega_{k})\sigma'(\omega_{k}^{*})\,
x_{\sigma}\wedge y_{\sigma'}$, the sum over $k$ of the coefficient is $1$ for
$\sigma=\sigma'$ and $0$ otherwise, as in the proof of
\Cref{thm:cmbasepoint}(ii), and $\sum_{\sigma}x_{\sigma}\wedge y_{\sigma}$ is
the complexification of $q(x,y)\Omega$. Finally let $\theta_{j}\in H^{1,1}(A)$
be the component of the polarisation that pairs $V_{\sigma}$ with
$V_{\bbar\sigma}$ for the two embeddings $\sigma,\bbar\sigma$ over $\tau_{j}$.
Then $\eta_{t}=\sum_{j}\tau_{j}(t)\,\theta_{j}$ for the classes $\eta_{t}$ of
\Cref{prop:cmweil}(iii), and $\CC[\theta_{1},\theta_{2}]$ is the
complexification of the $\QQ$-algebra generated by the $\eta_{t}$.

\begin{proposition}[Hodge classes, and a quaternion algebra acting on $R_{F}$]
\label{prop:quarticquat}
Let $A$ be a member of an $(F,2,\delta)$-family.
\begin{enumerate}[label=\textup{(\roman*)},leftmargin=2.6em,itemsep=2pt]
\item If $\Hg(A)=G_{F}$, the space of Hodge classes of $A$ has dimension $2$
  in degree two, spanned by the $\eta_{t}$, and $7$ in degree four, spanned by
  the three products of two $\eta_{t}$ and by $W_{F}(A)$, which meets the span
  of those products in $0$.
\item The map $\phi\colon R_{F}\to R_{F}$ defined by $q(x,y)=h(x,\phi y)$
  satisfies $\phi(ax)=\bbar a\,\phi(x)$ for $a\in F$ and
  $\phi^{2}=h_{4}(\Omega,\Omega)^{-1}$, and it commutes with $\SU(V,H)$. So
  $\phi$ and the action of $F$ are Hodge endomorphisms of $R_{F}$ at every
  member, and the Hodge classes in $R_{F}$ form a module over the algebra
  $\cQ_{\delta}=F\oplus F\phi$.
\item $\cQ_{\delta}$ is the cyclic quaternion algebra $(F/F_{0},\delta)$ over
  $F_{0}$, which is split if and only if $\delta$ is trivial. If
  $\Hg(A)=G_{F}$, then $\cQ_{\delta}$ is the whole algebra of Hodge
  endomorphisms of $R_{F}$.
\item If $\delta$ is trivial, then $R_{F}=W\otimes_{F_{0}}F$ for a rational
  sub-Hodge structure $W$, stable under $F_{0}$, of dimension $12$ and with
  Hodge numbers $(2,8,2)$.
\end{enumerate}
\end{proposition}

\begin{proof}
(i) As in the proof of \Cref{prop:cmweil}(iii), over $\CC$ the group $G_{F}$
is $\SL(U_{1})\times\SL(U_{2})$ with $\dim U_{j}=4$, and
$H^{1}(A,\CC)=\bigoplus_{j}(U_{j}\oplus U_{j}^{\vee})$. The invariants of
$\SL(U)$ in $\bigwedge^{a}U\otimes\bigwedge^{b}U^{\vee}$ form a line if $a=b$
or $\{a,b\}=\{0,4\}$, and are zero otherwise. In degree two this leaves the
bidegree $(1,1)$ in one factor: two lines, spanned by $\theta_{1}$ and
$\theta_{2}$. In degree four it leaves $(2,2)$ in one factor or $(1,1)$ in
both, three lines spanned by $\theta_{1}^{2}$, $\theta_{1}\theta_{2}$ and
$\theta_{2}^{2}$, which are nonzero because $\dim U_{j}=4$; and $(4,0)$ or
$(0,4)$ in one factor, four lines spanning $W_{F}\otimes\CC$. These lie in
different summands, so $W_{F}(A)$ meets the span of the products in $0$.

(ii) Choose an $H$-orthogonal basis $e_{1},\dots,e_{4}$ of $V$ and put
$a_{i}=H(e_{i},e_{i})\in F_{0}$. Replacing $\Omega$ by $\lambda\Omega$
replaces $\phi$ by $\bbar\lambda^{-1}\phi$ and $h_{4}(\Omega,\Omega)$ by
$\Nm(\lambda)h_{4}(\Omega,\Omega)$, so we may take
$\Omega=e_{1}\wedge e_{2}\wedge e_{3}\wedge e_{4}$, with
$h_{4}(\Omega,\Omega)=a_{1}a_{2}a_{3}a_{4}$. For $J=\{j<j'\}$ put
$e_{J}=e_{j}\wedge e_{j'}$ and $a_{J}=a_{j}a_{j'}$, and let $J^{c}$ be the
complement, with $e_{J^{c}}\wedge e_{J}=\epsilon_{J}\Omega$. The $e_{J}$ form
an $h$-orthogonal $F$-basis with $h(e_{J},e_{J})=a_{J}$, and
$q(e_{K},e_{J})$ is $\epsilon_{J}$ for $K=J^{c}$ and $0$ otherwise, so
$\phi(e_{J})=\epsilon_{J}a_{J^{c}}^{-1}e_{J^{c}}$. Two-forms commute, so
$\epsilon_{J^{c}}=\epsilon_{J}$ and
$\phi^{2}(e_{J})=(a_{J}a_{J^{c}})^{-1}e_{J}=h_{4}(\Omega,\Omega)^{-1}e_{J}$.
Semilinearity holds because $q$ is linear and $h$ antilinear in the second
variable. An element of $\SU(V,H)$ preserves $h$ and fixes $\Omega$, hence
preserves $q$ and commutes with $\phi$, and it commutes with $F$. The Hodge
group of every member lies in $G_{F}$, because the circle defining the Hodge
structure preserves $E$, commutes with $F$ and acts on each $V_{\sigma}$ with
determinant $z^{2}z^{-2}=1$; so $\phi$ and $F$ commute with it, and they are
Hodge endomorphisms.

(iii) The relations $\phi a=\bbar a\phi$ and $\phi^{2}=h_{4}(\Omega,\Omega)^{-1}$
present $F\oplus F\phi$ as the cyclic algebra
$(F/F_{0},h_{4}(\Omega,\Omega)^{-1})$, which is $(F/F_{0},\delta)$ because a
square in $F_{0}$ is a norm; and a cyclic algebra $(F/F_{0},c)$ is split
exactly when $c$ is a norm from $F$. It has $\QQ$-dimension $8$, since $\phi$
is not $F$-linear. If $\Hg(A)=G_{F}$, the Hodge endomorphisms of $R_{F}$ are
the commutant of $G_{F}$. Over $\CC$,
$R_{F}\otimes\CC=\bigoplus_{j}\bigl(\bigwedge^{2}U_{j}\oplus
\bigwedge^{2}U_{j}^{\vee}\bigr)$, where
$\bigwedge^{2}U_{j}\cong\bigwedge^{2}U_{j}^{\vee}$ is irreducible for
$\SL(U_{j})$ and trivial for the other factor, so the commutant is
$M_{2}(\CC)\times M_{2}(\CC)$, of dimension $8$.

(iv) If $\delta$ is trivial, $h_{4}(\Omega,\Omega)=\Nm(b)$ for some $b\in F$,
and $\psi=b\phi$ satisfies $\psi^{2}=b\bbar b\,\phi^{2}=1$. Its fixed space
$W$ is an $F_{0}$-subspace and a sub-Hodge structure, $\psi$ being a Hodge
endomorphism, and an element $\beta\in F$ with $\bbar\beta=-\beta\ne0$ maps it
isomorphically onto the $(-1)$-eigenspace, since $\psi\beta=-\beta\psi$. So
$R_{F}=W\oplus\beta W=W\otimes_{F_{0}}F$, and $W$ has half the dimension and
half the Hodge numbers of $R_{F}$.
\end{proof}

\begin{proposition}[What divisor classes reach]\label{prop:quarticnl}
Let $A$ be a member of an $(F,2,\delta)$-family.
\begin{enumerate}[label=\textup{(\roman*)},leftmargin=2.6em,itemsep=2pt]
\item The following are equivalent: \textup{(a)} $W_{F}(A)$ lies in the
  subalgebra of $H^{*}(A,\QQ)$ generated by $\NS(A)_{\QQ}$; \textup{(b)}
  $W_{F}(A)$ meets that subalgebra in a nonzero class; \textup{(c)} $R_{F}$
  contains a nonzero rational Hodge class. When they hold, $R_{F}$ contains a
  rational Hodge class $t$ with $q(t,t)\ne0$, and
  $W_{F}(A)=\{(ft)\star t:f\in F\}$, a space of polynomials with rational
  coefficients in divisor classes.
\item For a nonzero $\cQ_{\delta}$-submodule $M\subseteq R_{F}$ let
  $\cD(M)\subseteq\cD_{F}$ be the set of points at which $M$ consists of Hodge
  classes, and let $\mathrm{NL}(R_{F})\subseteq\cD_{F}$ be the set of points at
  which \textup{(c)} holds. Then $\mathrm{NL}(R_{F})=\bigcup_{M}\cD(M)$, a
  countable union. Each nonempty $\cD(M)$ is a closed complex submanifold,
  isomorphic to a product of two bounded symmetric domains of type IV and
  dimension $4-d$ each, where $d=\dim_{F}M$, and $d=1$ occurs only if $\delta$ is
  trivial. So if $\delta$ is trivial, $\mathrm{NL}(R_{F})$ is the union of the
  nonempty $\cD(M)$ with $d=1$, each of dimension exactly $6$; if $\delta$ is
  not trivial, every $\cD(M)$ has dimension at most $4$.
\item $\mathrm{NL}(R_{F})$ contains the base point $s_{0}$ of
  \Cref{thm:cmbasepoint}, and for every $\delta$ it contains the locus of
  products $B_{1}\times B_{2}$ of two abelian fourfolds of $(F,1)$-Weil type,
  which has dimension $4$. The split locus, of members isogenous to
  $X\otimes_{\cO_{F_{0}}}\cO_{F}$ with $X$ an abelian fourfold with
  $F_{0}\subseteq\End^{0}(X)$, has dimension $6$, lies in the family of
  trivial discriminant, and is contained in $\mathrm{NL}(R_{F})$.
\end{enumerate}
\end{proposition}

\begin{proof}
(i) Since $W_{F}(A)\ne0$, (a) implies (b). Let $F^{\times}$ act on $H^{*}(A,\QQ)$ by
$\alpha\mapsto\iota(\alpha)^{*}$, which on
$\bigotimes_{\sigma}\bigwedge^{a_{\sigma}}V_{\sigma}$ is the character
$\prod_{\sigma}\sigma^{a_{\sigma}}$. The isotypic parts for the characters
$\sigma^{2}$ in degree two and $\sigma^{4}$ in degree four are $R_{F}$ and
$W_{F}$. The projections $p_{R}$ and $p_{W}$ onto them are rational, the two
sets of characters being Galois orbits, and they are morphisms of Hodge
structures, being $\QQ$-linear combinations of the $\iota(\alpha)^{*}$. The
characters occurring in $H^{2}$ are the products $\rho\rho'$, and the only way
to write $\sigma^{4}$ as a product of two of them is $\sigma^{2}\cdot\sigma^{2}$,
so $p_{W}(D\cup D')=q(p_{R}D,p_{R}D')\,\Omega$ for $D,D'\in H^{2}(A,\QQ)$. If a
nonzero $w\in W_{F}$ equals $\sum_{i}D_{i}\cup D_{i}'$ with
$D_{i},D_{i}'\in\NS(A)_{\QQ}$, then
$w=p_{W}(w)=\sum_{i}q(p_{R}D_{i},p_{R}D_{i}')\,\Omega$, so some $p_{R}D_{i}$ is
a nonzero rational Hodge class in $R_{F}$, and (b) implies (c). Conversely,
let $t\in R_{F}$ be a nonzero rational Hodge class. The part of type $(1,1)$
of $\bigwedge^{2}V_{\sigma}$ is $V_{\sigma}^{1,0}\wedge V_{\sigma}^{0,1}$, and
$H$ is definite of opposite signs on the two factors by \Cref{prop:cmweil}(i),
so $h$ is negative definite on it at each real place, and $h(t,t)\in F_{0}$ is
totally negative. By \Cref{prop:quarticquat}(ii),
$q(t,\phi t)=h(t,\phi^{2}t)=h_{4}(\Omega,\Omega)^{-1}h(t,t)\ne0$, so one of the
rational Hodge classes $t$, $\phi t$, $t+\phi t$, call it $t'$, has
$q(t',t')\ne0$. Each $\omega_{k}ft'$ and each $\omega_{k}^{*}t'$ is a rational
Hodge class in $H^{2}$, hence a divisor class, and by \eqref{eq:quarticstar}
$(ft')\star t'=f\,q(t',t')\,\Omega$, which runs over $W_{F}$ as $f$ runs over
$F$. So (c) implies (a), with the stated formula.

(ii) The Hodge classes in $R_{F}$ form a $\cQ_{\delta}$-module, so
$\mathrm{NL}(R_{F})$ is the union of the $\cD(M)$, and there are countably
many $M$. Fix a real place $\tau$, let $\sigma$ be an embedding over it, and
put $R_{\tau}=R_{F}\otimes_{F_{0},\tau}\RR$, a complex vector space of
dimension $6$ through $F\otimes_{F_{0},\tau}\RR=\CC$, identified with
$\bigwedge^{2}V_{\sigma}$. Since $\det H$ has sign $(-1)^{n}=1$ at $\tau$, the
map $\phi$ induces an antilinear $\phi_{\tau}$ on $R_{\tau}$ with
$\phi_{\tau}^{2}=\lambda^{2}$, $\lambda>0$. Its eigenspace
$W_{\tau}=\{x:\phi_{\tau}x=\lambda x\}$ is a real form of $R_{\tau}$ and a real
sub-Hodge structure with $W_{\tau}\otimes\CC\cong\bigwedge^{2}V_{\sigma}$, of
Hodge numbers $(1,4,1)$, on which $q=\lambda h$ is real of signature $(2,4)$,
positive on the plane $\Pi_{\tau}$ of real classes of type $(2,0)+(0,2)$ and
negative on those of type $(1,1)$. The group $\SU(V_{\tau},H_{\tau})$ acts on
$W_{\tau}$ preserving $q$, through the isogeny $\SU(2,2)\to\mathrm{SO}(2,4)^{0}$,
and the map from $\cD_{2,2}$ to the space of oriented positive planes in
$W_{\tau}$ that sends a point to $\Pi_{\tau}$ is the equivariant isomorphism of
the symmetric domains of this group onto one component,
$\cD^{\mathrm{IV}}(W_{\tau})$, of dimension $4$. For a quadratic space $P$ of
signature $(2,k)$ write likewise $\cD^{\mathrm{IV}}(P)$, of dimension $k$. Let
$M$ have $F$-dimension $d$. Then $M_{\tau}=M\otimes_{F_{0},\tau}\RR$ is stable
under $\phi$ and $F$, so $M_{\tau}=P_{\tau}\otimes\CC$ with
$P_{\tau}=M_{\tau}\cap W_{\tau}$ of real dimension $d$, and $M$ consists of
Hodge classes exactly when every $P_{\tau}$ is of type $(1,1)$, that is,
orthogonal to $\Pi_{\tau}$. So
$\cD(M)=\prod_{\tau}\cD^{\mathrm{IV}}(P_{\tau}^{\perp})$, which is empty unless
every $P_{\tau}$ is negative definite and then has dimension $2(4-d)$. If
$d=1$, then $M=Fx$ with $\phi x=bx$ for some $b\in F$, and
$\phi^{2}x=\bbar bb\,x$ gives $h_{4}(\Omega,\Omega)^{-1}=\Nm(b)$, so $\delta$ is
trivial. If $\delta$ is trivial, $\cQ_{\delta}\cong M_{2}(F_{0})$, every nonzero
module contains one with $d=1$, and $\cD(M)\subseteq\cD(M')$ for
$M'\subseteq M$.

(iii) The classes $\delta_{i}(f)$ of \Cref{thm:cmbasepoint}(ii) have
complexification in $\bigoplus_{\sigma}\bigwedge^{2}V_{\sigma}$, so they lie in
$R_{F}$, and they are Hodge classes at $s_{0}$. For a product,
$V=V_{1}\oplus V_{2}$ with $H$-orthogonal $F$-planes of signature $(1,1)$ at
each place, and $\bigwedge^{2}_{F}V_{1}\subseteq R_{F}$ is $W_{F}(B_{1})$, of
type $(1,1)$ by \Cref{prop:cmweil}(i). Every $\delta$ is the product of the
discriminants of two such planes, for instance of $\mathrm{diag}(1,-\delta)$
and $\mathrm{diag}(1,-1)$, and each factor varies in a family of dimension
$m=2$. On the split locus $H^{1}(A,\QQ)=U\otimes_{F_{0}}F$, with
$U=H^{1}(X,\QQ)$ and the complex structure of $X$ on the first factor. The
polarisation class $e$ of $X$ lies in $\bigwedge^{2}_{F_{0}}U$, the hermitian
form of $A$ is $c$ times the sesquilinear extension of the $F_{0}$-bilinear
form $e$, with $c\in F$ totally imaginary, and so $\det H=c^{4}\,\mathrm{Pf}(e)^{2}$
is a square in $F_{0}$ and $\delta$ is trivial. The class
$e\otimes1\in(\bigwedge^{2}_{F_{0}}U)\otimes_{F_{0}}F=R_{F}$ is fixed by
$\Hg(A)=\Hg(X)\subseteq\Res_{F_{0}/\QQ}\Sp(U,e)$, so it is a nonzero Hodge
class. The polarised abelian fourfolds with real multiplication by $F_{0}$
form families of dimension $[F_{0}:\QQ]\cdot3=6$.
\end{proof}

\begin{corollary}[A base locus of codimension two]\label{cor:quarticbase}
In every $(F,2,\delta)$-family, $\Sigma_{F}\supseteq\mathrm{NL}(R_{F})$, and at
every point of $\mathrm{NL}(R_{F})$ the Weil classes are the polynomials in
divisor classes of \Cref{prop:quarticnl}(i). If $\delta$ is trivial,
$\mathrm{NL}(R_{F})$ is a countable union of closed complex submanifolds of
$\cD_{F}$ of codimension exactly two, and it contains the split locus; if
$\delta$ is not trivial, they have codimension at least four, and the product
loci have codimension four. In both cases $\mathrm{NL}(R_{F})$ contains the
base point $s_{0}$ of \Cref{thm:cmbasepoint} and is stable under
$G_{F}(\QQ)$. At a point outside $\mathrm{NL}(R_{F})$ no nonzero Weil class is
a polynomial in divisor classes.
\end{corollary}

\begin{proof}
Everything except the stability is \Cref{prop:quarticnl}. An element
$g\in G_{F}(\QQ)$ commutes with $F$ and with $\phi$, so it carries
$\cQ_{\delta}$-submodules to $\cQ_{\delta}$-submodules, and
$g\,\cD(M)=\cD(gM)$.
\end{proof}

\Cref{thm:cmbasepoint} gives one CM point, whose $G_{F}(\QQ)$-orbit is dense
but countable. When $\delta$ is trivial, \Cref{cor:quarticbase} replaces it by
a locus of codimension two on which the Weil classes are balanced products of
divisor classes, as they are at $s_{0}$; nothing is claimed about the rest of
the Hodge ring at its points. It does not decide $\mathbf{W}(F,2,\delta)$:
$\mathrm{NL}(R_{F})$ is a countable union of proper closed submanifolds, hence
meagre, and \Cref{thm:cmpropagation}(iii) allows $\Sigma_{F}$ to be meagre.
The last sentence of the corollary is the analogue of \Cref{prop:divfails}:
no argument through divisor classes reaches a member outside
$\mathrm{NL}(R_{F})$. The codimension is two rather than
$h^{2,0}(R_{F})=4$ (compare \Cref{rem:nlcodim}) because Hodge classes in
$R_{F}$ come in $\cQ_{\delta}$-modules, and a module of $F$-dimension one costs
one condition at each real place. \Cref{fig:quarticloci} shows where the Weil
classes of such a family are known to be algebraic.

\begin{figure}[tbp]
\centering
  \includegraphics[width=0.94\textwidth]{fig_quarticloci}
\caption{Where the Weil classes of an $(F,2,\delta)$-family are known to be
algebraic, $F$ a quartic CM field; each panel is the period domain
$\cD_{F}$, of dimension $8$. The grass regions are components of the
Noether--Lefschetz locus $\mathrm{NL}(R_{F})$, on which $W_{F}$ consists of
polynomials in divisor classes (\Cref{prop:quarticnl},
\Cref{cor:quarticbase}): of dimension exactly $6$ when $\delta$ is trivial,
where the quaternion algebra $\cQ_{\delta}$ of \Cref{prop:quarticquat} is
split, and at most $4$ otherwise. The ochre curve is the locus of scalar
extensions of Weil fourfolds, of dimension $4$ when $F$ is biquadratic, on
which Markman's theorem applies (\Cref{prop:quarticscalar}); it is not
contained in $\mathrm{NL}(R_{F})$ but meets it (ring) where the fourfold is a
product of Weil surfaces. The indigo points are CM points; the base point
$s_{0}$ of \Cref{thm:cmbasepoint} lies on $\mathrm{NL}(R_{F})$. For
$F=\QQ(\zeta_{5}),\QQ(\zeta_{8}),\QQ(\zeta_{12})$ and trivial $\delta$ the Prym
locus of \Cref{rem:quarticprym}, of dimension at most $6$, is not drawn. On
the rest of the domain, the very general member included, the algebraicity
of the Weil classes is the open propagation statement.}
\label{fig:quarticloci}
\end{figure}

\begin{proposition}[Scalar extension from an imaginary quadratic subfield]
\label{prop:quarticscalar}
Let $F=KF_{0}$ with $K$ imaginary quadratic; these are the biquadratic
quartic CM fields, and the only ones containing an imaginary quadratic field.
Let $C$ be an abelian fourfold of $(K,2,\delta)$-Weil type and
$B_{F}=C\otimes_{\cO_{K}}\cO_{F}$, which is of $(F,2,\iota(\delta))$-Weil type by
\Cref{prop:descent}(ii); write $(V_{C},H_{C})$ for the hermitian space of $C$.
For $f\in F$ let $g_{f}\in\Hom(B_{F},C)\otimes\QQ$ be the element inducing
$v\mapsto v\otimes f$ from $H^{1}(C,\QQ)$ to
$H^{1}(B_{F},\QQ)=H^{1}(C,\QQ)\otimes_{K}F$.
\begin{enumerate}[label=\textup{(\roman*)},leftmargin=2.6em,itemsep=2pt]
\item $W_{F}(B_{F})$ lies in the $\QQ$-span of the classes $g_{f}^{*}w$, $f\in F$,
  $w\in W_{K}(C)$, which has dimension $10$. Hence $W_{F}(B_{F})$ is
  algebraic, by Markman's theorem $\mathbf{W}(K,2,\delta)$ \cite{Mar25a}
  (\Cref{prop:separate}(ii)).
\item The members $B_{F}$ form a locus of dimension $4$ in the
  $(F,2,\iota(\delta))$-family. At those with $\Hg(C)=\SU(V_{C},H_{C})$ the
  space $R_{F}$ contains no nonzero Hodge class, so
  this locus is not contained in $\mathrm{NL}(R_{F})$; it meets
  $\mathrm{NL}(R_{F})$, for instance where $C=C_{1}\times C_{2}$ with $C_{1}$,
  $C_{2}$ abelian surfaces of Weil type, and then $B_{F}$ is a product of two
  abelian fourfolds of $(F,1)$-Weil type.
\end{enumerate}
\end{proposition}

\begin{proof}
(i) The map $v\mapsto v\otimes f$ commutes with the complex structures, that of
$B_{F}$ being the one of $C$ tensored with $F$, so $g_{f}$ exists. Over $\CC$,
for each embedding $\kappa$ of $K$ let $V_{\kappa}$ be the $\kappa$-eigenspace
of $K$ on $H^{1}(C,\CC)$ and $U_{\kappa}=F\otimes_{K,\kappa}\CC\cong\CC^{2}$, the
sum of the lines $L_{\sigma}$, $\sigma|\kappa$, on which $F$ acts through
$\sigma$. The $\kappa$-part of $H^{1}(B_{F},\CC)$ is $V_{\kappa}\otimes
U_{\kappa}$ and its $\sigma$-part is $V_{\kappa}\otimes L_{\sigma}$, so the
$\sigma$-component $\bigwedge^{4}V_{\kappa}\otimes L_{\sigma}^{\otimes4}$ of
$W_{F}(B_{F})\otimes\CC$ lies in $\bigwedge^{4}V_{\kappa}\otimes
\mathrm{Sym}^{4}U_{\kappa}\subseteq\bigwedge^{4}(V_{\kappa}\otimes U_{\kappa})$.
The map $g_{f}^{*}$ sends a generator $\alpha_{\kappa}$ of $\bigwedge^{4}V_{\kappa}$
to $\alpha_{\kappa}\otimes f_{\kappa}^{\otimes4}$, with $f_{\kappa}$ the image
of $f$ in $U_{\kappa}$. The fourth powers of the elements of the Zariski
dense subset $F$ of $U_{\kappa}$ span $\mathrm{Sym}^{4}U_{\kappa}$, so the
complexified span of the $g_{f}^{*}w$ is $\bigoplus_{\kappa}\bigwedge^{4}V_{\kappa}
\otimes\mathrm{Sym}^{4}U_{\kappa}$, of dimension $2\cdot5$, and it contains
$W_{F}(B_{F})\otimes\CC$. A rational subspace whose complexification lies in
that of a rational subspace lies in it, and pull-backs of algebraic classes
are algebraic.

(ii) The map $C\mapsto B_{F}$ embeds the four-dimensional domain of the
$K$-family into $\cD_{F}$. As rational Hodge structures
$H^{1}(B_{F},\QQ)\cong H^{1}(C,\QQ)^{2}$, so $\Hg(B_{F})=\Hg(C)$, and
$R_{F}\otimes\CC=\bigoplus_{\sigma}\bigwedge^{2}V_{\kappa}\otimes
L_{\sigma}^{\otimes2}$, with $\bigwedge^{2}V_{\kappa}$ irreducible and
nontrivial for $\SL(V_{\kappa})$; so when $\Hg(C)=\SU(V_{C},H_{C})$, $R_{F}$
has no invariants. For $C=C_{1}\times C_{2}$ one has
$B_{F}=(C_{1}\otimes_{\cO_{K}}\cO_{F})\times(C_{2}\otimes_{\cO_{K}}\cO_{F})$, and
\Cref{prop:quarticnl}(iii) applies.
\end{proof}

\Cref{prop:quarticscalar} gives no new case of the Hodge conjecture: $B_{F}$ is
isogenous to $C\times C$, the classes are Hodge classes there, and when
$\Hg(C)=\SU(V_{C},H_{C})$ \Cref{prop:powers} already gives the conjecture for
$C\times C$ from \cite{Mar25a}.

\begin{remark}[Prym varieties of cyclic covers]\label{rem:quarticprym}
For three cyclotomic fields published work gives a further locus, at whose
very general points the Hodge group is the whole of $G_{F}$. Let $N\ge3$,
let $\widetilde C\to C$ be an \'etale cyclic cover of degree $N$ of a smooth
projective curve of genus
$g\ge2$, and let $P\subseteq J(\widetilde C)$ be the abelian subvariety on
whose $H^{1}$ the covering group acts through characters of exact order $N$.
Then $F=\QQ(\zeta_{N})$ acts on $P$, and every embedding of $F$ occurs in
$H^{1,0}(P)$ with multiplicity $g-1$ \cite[Lemma~1.6a]{Sch88},
\cite[Lemma~5.1]{PZ25}, so $P$ is of $(F,g-1)$-Weil type
(\Cref{rem:cmmember}). The hermitian form is split, because the part of
$H_{1}(\widetilde C,\ZZ)$ on which the covering group acts through
characters of exact order $N$ has a hyperbolic $\ZZ[\zeta_{N}]$-basis
\cite[Proposition~2.2]{Loo97}; so $\delta=[(-1)^{g-1}]$, which is trivial for
$g=3$. Schoen proved that the Weil classes of $P$ are algebraic
\cite[Theorem~2.0 with $r=0$, Corollary~3.1]{Sch88}, and Patel and Zhang give
a proof by unramified geometric class field theory
\cite[Theorems~1.1, 1.2 and~5.3, Remark~3.7]{PZ25}. By
\cite[Theorem~2.4]{Loo97}, for every $g\ge2$, the monodromy group of the
family of the $P$ over a finite cover of the moduli of such covers contains a
subgroup of finite index of an arithmetic subgroup of $G_{F}$. Every factor
$\SU(g-1,g-1)$ of $G_{F}(\RR)$ is noncompact, so by Borel's density theorem
\cite{Bor60} its Zariski closure contains $G_{F}$; the Hodge group of every
member lies in $G_{F}$, so by Andr\'e's theorem \cite[Theorem~1]{And92} a very
general $P$ has $\Hg(P)=G_{F}$, and by \Cref{prop:cmweil}(iii) the Hodge
conjecture holds for it in every degree. All of this is a consequence of
\cite{Sch88,Loo97,And92}, and we make no claim of priority.

For $\varphi(N)=4$, that is for $F=\QQ(\zeta_{5})=\QQ(\zeta_{10})$,
$\QQ(\zeta_{8})$ and $\QQ(\zeta_{12})$, and $g=3$, the $P$ are abelian
eightfolds of $(F,2,1)$-Weil type. With their $G_{F}(\QQ)$-translates, which
lie in $\Sigma_{F}$ by \Cref{thm:cmpropagation}(ii), they fill a countable
union of analytic subvarieties of $\cD_{F}$, the \emph{Prym locus} of $F$, of
positive dimension, since the monodromy is infinite, and of dimension at most
$3g-3=6$ \cite[pp.~25--26]{Sch88}. At a very general point of it $\Hg=G_{F}$, so
$R_{F}$ has no nonzero Hodge class there, each $\bigwedge^{2}V_{\sigma}$ being
a nontrivial irreducible representation of $\SL(V_{\sigma})$, and the Hodge
group is larger than that of a scalar extension of a Weil fourfold: such
points lie outside $\mathrm{NL}(R_{F})$, outside the locus of
\Cref{prop:quarticscalar}, and are not CM points. They are the only members
with $\Hg=G_{F}$ at which we know the Weil classes of a quartic CM field to be
algebraic. They do not decide $\mathbf{W}(F,2,1)$: the count
$3g-3=6<8=\tfrac12\varphi(N)(g-1)^{2}=\dim\cD_{F}$ of \cite[pp.~25--26]{Sch88}
makes the union meagre, as $\mathrm{NL}(R_{F})$ is, and
\Cref{thm:cmpropagation}(iii) allows $\Sigma_{F}$ to be meagre. Nor does
\Cref{prop:quarticnogo} exclude them: the cycles come from $J(\widetilde C)$,
which has $P$ as a factor. For imaginary quadratic fields the count is
$3n\ge n^{2}$ for $n\le3$, and there the Prym varieties do reach whole
families: Schoen proved the algebraicity for the Weil fourfolds and for a
family of Weil sixfolds of $\QQ(\sqrt{-3})$ \cite[Theorem~3.2]{Sch88},
\cite{Sch98}, and Koike for a family of Weil sixfolds of $\QQ(\sqrt{-1})$
\cite[Theorem~2.1, Corollary~2.1]{Koi04}.
\end{remark}

\begin{proposition}[The Casimir class]\label{prop:quarticcasimir}
Let $A$ be a member of an $(F,2,\delta)$-family and $c\in F^{\times}$, and let
\[
  C_{q}\in R_{F}\otimes_{\QQ}R_{F}\subseteq H^{2}(A,\QQ)\otimes H^{2}(A,\QQ)
  \subseteq H^{4}(A\times A,\QQ)
\]
be the Casimir element $\sum_{j}u_{j}\otimes u_{j}^{*}$ of the nondegenerate
symmetric form $\Tr_{F/\QQ}(c^{-1}q)$ on $R_{F}$, for a basis $(u_{j})$ and its
dual basis $(u_{j}^{*})$. Then $C_{q}$ is a Hodge class,
$\Delta^{*}C_{q}=6c\,\Omega$ for the diagonal $\Delta\colon A\to A\times A$,
and $C_{q}$ lies in the $\QQ$-span of the classes $f^{*}w$, $w\in W_{F}(A)$,
$f(x,y)=\iota(a)x+\iota(b)y$ with $a,b\in\cO_{F}$. Hence $W_{F}(A)$ is
algebraic if and only if $C_{q}$ is.
\end{proposition}

\begin{proof}
For an $F$-basis $(x_{I})$ of $R_{F}$ and its $q$-dual basis $(x_{I}^{q})$, the
elements $\omega_{k}x_{I}$ and $c\,\omega_{k}^{*}x_{I}^{q}$ form dual bases for
$\Tr_{F/\QQ}(c^{-1}q)$. So by \eqref{eq:quarticstar}
$\Delta^{*}C_{q}=\sum_{I}x_{I}\star(c\,x_{I}^{q})=\sum_{I}c\,\Omega=6c\,\Omega$,
and if $C_{q}$ is algebraic so is $W_{F}(A)$, by \Cref{prop:cmweil}(ii).
Conversely, $\SU(V,H)$ preserves $q$, so $C_{q}$ is invariant under the
diagonal action of $G_{F}$, which contains $\Hg(A\times A)$, and $C_{q}$ is a
Hodge class. Since $q$ pairs $\bigwedge^{2}V_{\sigma}$ only with itself, the
complexification of $C_{q}$ lies in $\bigoplus_{\sigma}\bigwedge^{2}V_{\sigma}
\otimes\bigwedge^{2}V_{\sigma}$, inside the part
$\bigoplus_{\sigma}\bigwedge^{4}(V_{\sigma}\otimes\CC^{2})$ of
$H^{4}(A\times A,\CC)$ on which $F^{\times}$, acting diagonally, has the
characters $\sigma^{4}$. The group $G_{F}$ acts on $V_{\sigma}$ through
$\SL(V_{\sigma})$, whose invariants in $\bigwedge^{4}(V_{\sigma}\otimes\CC^{2})$
are $\bigwedge^{4}V_{\sigma}\otimes\mathrm{Sym}^{4}\CC^{2}$ by Cauchy's
formula, so the $G_{F}$-invariants in that part form the rational space
$W_{F}(A)\otimes_{F}\mathrm{Sym}^{4}_{F}(F^{2})$, of dimension $20$. The map
$f^{*}$ sends $w$ to $w\otimes(a,b)^{\otimes4}$, and the fourth powers of the
elements of $\cO_{F}^{2}$ span $\mathrm{Sym}^{4}_{F}(F^{2})$; so the classes
$f^{*}w$ span that space, and $C_{q}$ is algebraic if $W_{F}(A)$ is.
\end{proof}

\begin{remark}[What the exceptional isogeny gives]\label{rem:quarticisogeny}
By the proof of \Cref{prop:quarticnl}(ii), at each real place $R_{F}$ becomes
$W_{\tau}\otimes\CC$ with $W_{\tau}$ of K3 type, $\SU(V_{\tau},H_{\tau})$ acting
through $\SU(2,2)\to\mathrm{SO}(2,4)$, and when $\delta$ is trivial the
$W_{\tau}$ are the real components of the rational structure $W$ of
\Cref{prop:quarticquat}(iv). A route of Kuga--Satake type would start from
this weight two structure and would have to return to degree four, on $A$ or
on $A\times A$; by \Cref{prop:quarticcasimir} the natural class it meets
there, the Casimir class of $q$, is algebraic exactly when the Weil classes
are. So this mechanism restates $\mathbf{W}(F,2,\delta)$ at each member and
does not reduce it.
\end{remark}

\begin{proposition}[No correspondence from a small abelian variety]
\label{prop:quarticnogo}
Let $A$ be a member with $\Hg(A)=G_{F}$, and let $X$ be an abelian variety
such that $\Hg(X)$ has no normal subgroup isogenous to $G_{F}$; by the proof,
this holds when every simple factor of $X$ has dimension at most $7$, in
particular when $\dim X\le7$, and it holds when $X$ is of CM type. Then
$\Hg(A\times X)=\Hg(A)\times\Hg(X)$, and
$\Hdg^{*}(A^{a}\times X^{b})=\Hdg^{*}(A^{a})\otimes\Hdg^{*}(X^{b})$ for all
$a,b\ge0$. So the class of an algebraic cycle $Z$ on $A^{a}\times X^{b}$ is
$\sum_{i}\beta_{i}\otimes\xi_{i}$ with $\beta_{i}\in\Hdg^{*}(A^{a})$ and
linearly independent $\xi_{i}\in\Hdg^{*}(X^{b})$; the correspondence $Z$
carries $H^{*}(X^{b},\QQ)$ into the span of the $\beta_{i}$; and if the Hodge
conjecture holds for $X^{b}$, every $\beta_{i}$ is algebraic.
\end{proposition}

\begin{proof}
The group $G_{F}=\Res_{F_{0}/\QQ}\SU(V,H)$ is almost simple over $\QQ$. By
Goursat's lemma the kernel $N$ of the projection $\Hg(A\times X)\to\Hg(X)$
maps onto a normal subgroup of $\Hg(A)=G_{F}$. If that subgroup is $G_{F}$,
then $\Hg(A\times X)$ contains $G_{F}\times1$ and is the product. Otherwise
$N$ is finite, $\Hg(X)$ has a normal subgroup $H$ isogenous to $G_{F}$,
contrary to the hypothesis, and the cocharacter $\mu$ defining the Hodge
structure of $X$ projects through $H\to G_{F}$ to that of $A$. This situation
also cannot occur when every simple factor of $X$ has dimension at most $7$,
or when $X$ is of CM type, as follows. Over $\CC$, $H$ is $\SL_{4}\times\SL_{4}$,
and $\mu$ projects in each factor to
$\mathrm{diag}(\frac12,\frac12,-\frac12,-\frac12)$ modulo the centre. Since
$\mu$ acts on $H^{1}(X,\CC)$ with two weights differing by one, the weights of
its projection spread by at most one on every irreducible constituent for
$H$; for the representation of $\SL_{4}$ of highest weight
$a\varpi_{1}+b\varpi_{2}+c\varpi_{3}$ the spread is $a+2b+c$, and spreads add
over the two factors. So every nontrivial constituent is the standard
representation of one factor or its dual. These four are permuted
transitively by $\Gal(\bbar\QQ/\QQ)$, as the embeddings of $F$ are; so a simple
factor $Y$ of $X$ on whose $H^{1}$ the group $H$ acts nontrivially has
$\dim_{\QQ}H^{1}(Y,\QQ)\ge16$, that is $\dim Y\ge8$. A torus has no such $H$.
Next, $\Hg(A^{a}\times X^{b})$ is $\Hg(A)\times\Hg(X)$ acting factorwise on
$H^{*}(A^{a})\otimes H^{*}(X^{b})$, and the invariants of a product of groups
in a tensor product of representations of the factors are the tensor product
of the invariants. Finally, the intersection pairing is perfect on
$\Hdg^{*}(X^{b})$, the invariants being orthogonal to the other isotypic
parts, so there are $\xi_{j}^{*}\in\Hdg^{*}(X^{b})$ with
$\int\xi_{i}\xi_{j}^{*}=\delta_{ij}$; they are algebraic when the conjecture
holds for $X^{b}$, and then $\beta_{j}=pr_{A^{a}*}(Z\cdot pr_{X^{b}}^{*}\xi_{j}^{*})$
is algebraic.
\end{proof}

So a correspondence with such an $X$ can produce a Weil class of $A$ only from
an algebraic cycle on some $A^{a}\times X^{b}$ that already yields one on
$A^{a}$. At a member with full Hodge group this excludes every abelian
variety of dimension at most seven, among them the fourfolds and sixfolds of
Weil type of the imaginary quadratic fields on which Markman's theorems are
proved, and every CM abelian variety. The same argument, with real points in
place of dimensions, excludes rational Hodge structures $Y$ of K3 type: by
Zarhin's theorem \cite{Zar83} the Hodge group of such a $Y$ is the restriction
of scalars of an orthogonal or unitary group over the endomorphism field of
its transcendental part, noncompact at exactly one real place, so its real Lie
algebra has at most one simple factor isomorphic to $\mathfrak{su}(2,2)$,
while that of $G_{F}$ has two; hence $\Hg(A\times Y)=\Hg(A)\times\Hg(Y)$.
This last statement rests on Zarhin's theorem, which we quote without proof.

\begin{theorem}[The polarised criterion for a quartic field at $n=2$]
\label{thm:quarticp2prime}
Let $A$ be a member of an $(F,2,\delta)$-family, let
$\omega=\sum_{\sigma}u_{\sigma}\alpha_{\sigma}$, with $\alpha_{\sigma}$ spanning
$\bigwedge^{4}V_{\sigma}$ and every $u_{\sigma}\ne0$, as holds for every
nonzero rational Weil class, let $p\in\CC[\theta_{1},\theta_{2}]$, put
$\gamma=\omega+p$, and let $r(\gamma)$ be the rank of the contraction
$HT^{2}(A)\to H^{*}(A,\CC)$, $\xi\mapsto\xi\lrcorner\gamma$. Write
$T_{\sigma}=(V_{\sigma}^{1,0})^{\vee}\subseteq T_{0}A$.
\begin{enumerate}[label=\textup{(\roman*)},leftmargin=2.6em,itemsep=2pt]
\item $\dim HT^{2}(A)=28+64+28=120$, and
  \[
    \Ann_{HT^{2}}(\omega)=\textstyle\bigoplus_{\sigma}
    \Hom\bigl(V_{\sigma}^{1,0},V_{\sigma}^{0,1}\bigr)\;\oplus\;
    \bigoplus_{\{\sigma,\sigma'\},\,\sigma\ne\sigma'}T_{\sigma}\otimes T_{\sigma'},
  \]
  of dimension $0+16+24$: the $F$-linear first order deformations of $A$ and
  the bivectors mixing two embeddings. So $r(\omega)=80$. Unlike the case
  $m=1$ of \Cref{thm:p2forces}(i), no nonzero element of $H^{0,2}(A)$ kills
  $\omega$.
\item $r(\gamma)\le112$ for every $p$. For a general $p$, $r(\gamma)=112$, and
  whenever $r(\gamma)=112$ the annihilator of $\gamma$ is the tangent space
  $T_{[A]}\cD_{F}$ of the family, of dimension $8$, inside $H^{1}(T_{A})$. The
  value $112$ is also taken at a general element, with nonzero component in
  $W_{F}\otimes\CC$, of the complexified Hodge ring of a member with
  $\Hg(A)=G_{F}$.
\item $r(\gamma)\ge80$ for every $p$, with equality exactly when $p$ is a
  constant.
\item The rank $r(\gamma)$ takes exactly the fifteen values $80$, $84$, $87$,
  $88$, $91$, $92$, $94$, $95$, $96$, $104$, $107$, $108$, $110$, $111$,
  $112$; \Cref{thm:quarticrank} computes it for every $p$. On the monomials
  $\theta_{1}^{i}\theta_{2}^{j}$ with a general coefficient it takes the
  six values $80$, $84$, $88$, $104$, $108$, $112$.
\item Let $E$ be a perfect complex on $A$ with $\ch(E)=N\omega_{1}+P$, where
  $N\ne0$, $\omega_{1}$ is a nonzero rational Weil class and $P$ is a
  polynomial with rational coefficients in the classes $\eta_{t}$. Then
  $\dim\Ext^{2}(E,E)\ge r(\ch E)\ge80$. If $\dim\Ext^{2}(E,E)=r(\ch E)$, then
  $\operatorname{ev}_{E}$ is onto with kernel $\Ann_{HT^{2}}(\ch E)$ and
  $\sigma_{E}$ is injective, so that $\mathbf{W}(F,2,\delta)$ follows by
  \Cref{thm:cmpropagation}(v).
\end{enumerate}
\end{theorem}

\begin{proof}
(i) $HT^{2}(A)=\bigwedge^{2}H^{0,1}\oplus H^{0,1}\otimes T_{0}A\oplus
\bigwedge^{2}T_{0}A$, with $H^{0,1}$ and $T_{0}A$ of dimension $8$. As in the
proof of \Cref{thm:p2forces}(i) the three summands map into different Hodge
components, so the annihilator is the sum of the three annihilators; and the
torus $(\CC^{\times})^{4}$ acting on each $V_{\sigma}$ by scalars separates the
images further. A block $V_{\rho}^{0,1}\wedge V_{\rho'}^{0,1}$ of
$\bigwedge^{2}H^{0,1}$ kills $\alpha_{\sigma}$ exactly when
$\sigma\in\{\rho,\rho'\}$, and otherwise maps injectively into a weight space
that no other pair of a block and an embedding reaches; there are four
embeddings, so some $\sigma$ lies outside $\{\rho,\rho'\}$, and no nonzero
element of $\bigwedge^{2}H^{0,1}$ kills $\omega$. A block
$\Hom(V_{\rho}^{1,0},V_{\rho'}^{0,1})$ of $H^{1}(T_{A})$ kills $\alpha_{\sigma}$
for $\sigma\ne\rho$; it kills $\alpha_{\rho}$ when $\rho'=\rho$, because
$V_{\rho}^{0,1}$ is already exhausted in $\alpha_{\rho}$, and maps it
injectively into its own weight space when $\rho'\ne\rho$. A bivector in
$T_{\rho}\otimes T_{\rho'}$ with $\rho\ne\rho'$ kills every $\alpha_{\sigma}$,
whose two factors of type $(1,0)$ both lie in $V_{\sigma}^{1,0}$, and
$\bigwedge^{2}T_{\sigma}$ maps $\alpha_{\sigma}$ to a nonzero multiple of a
generator of $\bigwedge^{2}V_{\sigma}^{0,1}$. The dimensions are $0$,
$4\cdot4$ and $6\cdot4$.

(ii) The classes $\theta_{j}$ and the rational Weil classes stay of Hodge type
along $\cD_{F}$, so each is killed by $T_{[A]}\cD_{F}\subseteq H^{1}(T_{A})$, as
in \Cref{prop:p2primelocus}(i); hence $r(\gamma)\le120-8$, and when equality
holds the annihilator, of dimension $8$, is $T_{[A]}\cD_{F}$. Take a basis of
$H^{1}(A,\CC)$ adapted to the $V_{\sigma}$ and to the Hodge decomposition, in
which the polarisation pairs the basis of $V_{\sigma}^{1,0}$ with that of
$V_{\bbar\sigma}^{0,1}$. The torus acting diagonally on it contains a subtorus
that multiplies each $\theta_{j}$ by a scalar and each $\alpha_{\sigma}$ by an
arbitrary scalar, and contraction is equivariant; so $r(\omega+p)$ depends on
$\omega$ only through a rescaling of the variables of $p$, and each statement
of the theorem holds for every $\omega$ once it holds for one. The rank is
lower semicontinuous in $p$, and it equals $112$ at
$p=\theta_{1}^{2}\theta_{2}^{2}$ by \Cref{thm:quarticrank}(iii) below, so it
equals $112$ for a general $p$. Every element of the complexified Hodge ring
of a member with $\Hg(A)=G_{F}$ is invariant under $G_{F}$, so its components
stay of Hodge type along $\cD_{F}$ and it is killed by $T_{[A]}\cD_{F}$; its
rank is therefore at most $112$. The ring contains $\omega+\theta_{1}^{2}
\theta_{2}^{2}$, since it contains $W_{F}\otimes\CC$ and the classes
$\theta_{1}$, $\theta_{2}$, whose span is defined over $\QQ$, and the rank
is lower semicontinuous; so it is $112$ at a general element of the ring, whose
component in $W_{F}\otimes\CC$ is not zero.

(iii) and (iv) are \Cref{thm:quarticrank}(iii) and (iv) below.

(v) By the commutative square $\sigma_{E}\circ\operatorname{ev}_{E}=(\,\cdot\,)
\lrcorner\ch(E)$ of \Cref{setup:criterion}, which holds on every abelian
variety, $\dim\Ext^{2}(E,E)\ge\operatorname{rk}\operatorname{ev}_{E}\ge
r(\ch E)$, and equality forces $\operatorname{ev}_{E}$ to be onto with kernel
the annihilator and $\sigma_{E}$ to be injective, as in
\Cref{thm:p2primenumber}(iv); and $r(\ch E)\ge80$ by (iii), since the
complexification of $\ch(E)$ is a $\gamma$ as above with $\omega=N\omega_{1}$.
The last implication is \Cref{thm:cmpropagation}(v), whose character is
exactly of this form.
\end{proof}

At $m=1$ and $n=2$ the corresponding numbers are $28$, $12$ and $24$
(\Cref{thm:p2primenumber}); in both cases the general value is the dimension
of $HT^{2}$ minus that of the family. For the pure shape the argument of
\Cref{thm:p2forces}(ii) and (iv) transfers with the annihilator of
\Cref{thm:quarticp2prime}(i): an injective $\sigma_{E}$ would force
$A_{\theta}A_{\theta'}=0$ for $\theta\in T_{\rho}$, $\theta'\in T_{\rho'}$,
$\rho\ne\rho'$, and a smooth component of the support of codimension at least
two along which $E$ is locally free of positive rank would be tangent to three
of the four $T_{\rho}$. The pure shape is excluded, for every CM
field, by \Cref{thm:cmpurefalse}.

The rank $r(\gamma)$ can be computed for every $p$. In the basis of the
proof of \Cref{thm:quarticp2prime}(ii) write $x_{\sigma,1},x_{\sigma,2}$ for
the basis of $V_{\sigma}^{1,0}$, $y_{\sigma,1},y_{\sigma,2}$ for that of
$V_{\sigma}^{0,1}$, and $\partial_{\sigma,1},\partial_{\sigma,2}$ for the dual
basis of $T_{\sigma}$, so that $\theta_{j}=\sum_{\sigma|\tau_{j}}\sum_{k}
x_{\sigma,k}\,y_{\bbar\sigma,k}$ and $\alpha_{\sigma}=x_{\sigma,1}x_{\sigma,2}
y_{\sigma,1}y_{\sigma,2}$, and use the torus of that proof to make every
$u_{\sigma}$ equal to $1$. Expand
$p=\sum_{i,j=0}^{4}e_{ij}\,\theta_{1}^{i}\theta_{2}^{j}/(i!\,j!)$, so that
$e_{ij}$ is a derivative of $p$ at $0$, and put:
\begin{itemize}[leftmargin=1.6em,itemsep=1pt]
\item $\mu(p)=0$ if $\partial_{1}\partial_{2}p=0$ and $\mu(p)=1$ otherwise,
  the derivatives taken in $\CC[\theta_{1},\theta_{2}]$;
\item $\rho_{1}(p)$ the dimension of the span of $\partial_{1}p$ and
  $\partial_{1}^{2}p$ modulo $\theta_{1}^{3}$, which is the rank of the
  $2\times15$ matrix whose columns are the pairs $(e_{i,j},e_{i+1,j})$,
  $1\le i\le3$, $0\le j\le4$; and $\rho_{2}(p)$ likewise, with the two indices
  exchanged;
\item for $t=1,2$, with $\sigma,\bbar\sigma$ the embeddings over $\tau_{t}$,
  $W_{t}\subseteq HT^{2}(A)$ the span of the eight elements
  $y_{\sigma,1}\wedge y_{\sigma,2}$, $y_{\bbar\sigma,1}\wedge
  y_{\bbar\sigma,2}$, $\partial_{\sigma,1}\wedge\partial_{\sigma,2}$,
  $\partial_{\bbar\sigma,1}\wedge\partial_{\bbar\sigma,2}$ and
  $y_{\sigma,k}\otimes\partial_{\bbar\sigma,l}$,
  $y_{\bbar\sigma,k}\otimes\partial_{\sigma,l}$ with $k\ne l$, and $R_{t}(p)$
  the rank of $\xi\mapsto\xi\lrcorner\gamma$ on $W_{t}$;
\item $h_{1}=e_{1}e_{3}-e_{2}^{2}$, $h_{2}=e_{2}e_{4}-e_{3}^{2}$,
  $h_{3}=e_{1}e_{4}-e_{2}e_{3}$ on $\CC^{4}$, and $V_{\pm}\subset\CC^{4}$ the
  set where $h_{3}=\pm e_{3}$, $2h_{2}=\pm e_{4}$ and $2h_{1}+1=\pm e_{2}$.
\end{itemize}

\begin{theorem}[The rank of the criterion at a quartic field]
\label{thm:quarticrank}
In this notation:
\begin{enumerate}[label=\textup{(\roman*)},leftmargin=2.6em,itemsep=2pt]
\item $r(\gamma)=64+16\mu(p)+4\rho_{1}(p)+4\rho_{2}(p)+R_{1}(p)+R_{2}(p)$.
\item $R_{t}(p)\in\{7,8\}$. If $R_{1}(p)=7$, then
  $(e_{10},e_{20},e_{30},e_{40})$ lies on $V_{+}\cup V_{-}$ and
  $\rho_{1}(p)=2$; conversely, if $\mu(p)=0$ and
  $(e_{10},\dots,e_{40})$ lies on $V_{+}\cup V_{-}$, then $R_{1}(p)=7$, as for
  $p=\theta_{1}^{2}/2$. The same holds for $R_{2}$ with
  $(e_{01},\dots,e_{04})$. Each $V_{\pm}$ contains the surface of points
  $(w\lambda\pm\frac{3}{2\lambda},\,w\lambda^{2}\pm\frac12,\,w\lambda^{3},
  \,w\lambda^{4})$, $w\in\CC$, $\lambda\in\CC^{\times}$.
\item $r(\gamma)\ge80$, with equality exactly when $p$ is a constant, and
  $r(\gamma)$ takes exactly the fifteen values
  \[
    80,\ 84,\ 87,\ 88,\ 91,\ 92,\ 94,\ 95,\ 96,\ 104,\ 107,\ 108,\ 110,\ 111,\
    112 ,
  \]
  the values $87$, $91$, $94$, $95$, $107$, $110$ and $111$ only when $R_{1}$
  or $R_{2}$ is $7$. The value $100$ is not taken.
\item For $\omega=\sum_{\sigma}u_{\sigma}\alpha_{\sigma}$ with arbitrary
  nonzero $u_{\sigma}$, the same holds with $e_{ij}$ replaced by
  $\kappa_{1}^{i}\kappa_{2}^{j}e_{ij}$, where $\kappa_{t}^{4}=
  (u_{\sigma}u_{\bbar\sigma})^{-1}$ for the embeddings over $\tau_{t}$. The
  four choices of $\kappa_{t}$ fix $V_{+}\cup V_{-}$, and $\mu$, $\rho_{1}$,
  $\rho_{2}$ do not depend on them.
\end{enumerate}
\end{theorem}

\begin{proof}
(i) Let $S$ be the torus of the basis that fixes $\theta_{1}$, $\theta_{2}$
and every $\alpha_{\sigma}$. If $x_{\sigma,k}$ has weight $a_{\sigma,k}$, then
$y_{\sigma,k}$ has weight $-a_{\bbar\sigma,k}$ and $\partial_{\sigma,k}$ has
weight $-a_{\sigma,k}$, and the only relations are
$a_{\sigma,1}+a_{\sigma,2}=a_{\bbar\sigma,1}+a_{\bbar\sigma,2}$, one for each
real place; so $S$ has dimension six. Attach to $y_{\sigma,k}$ the label
$(\bbar\sigma,k)$ and to $\partial_{\sigma,k}$ the label $(\sigma,k)$. A basis
element of $HT^{2}$ is a product of two of these and has weight
$-(a_{P}+a_{Q})$ for its two labels $P,Q$, and two such weights agree exactly
when the pairs of labels coincide or are $\{(\sigma,1),(\sigma,2)\}$ and
$\{(\bbar\sigma,1),(\bbar\sigma,2)\}$. So $HT^{2}$ splits into $34$ weight
spaces: the eight lines spanned by $y_{\sigma,k}\otimes
\partial_{\bbar\sigma,k}$; sixteen spaces of dimension four, one for each
pair of labels over different real places; eight of dimension four, one for
each pair $\{(\sigma,k),(\bbar\sigma,l)\}$ over one place; and $W_{1}$,
$W_{2}$. The class $\gamma$ is invariant under $S$, so contraction carries
each weight space into the same weight of $H^{*}(A,\CC)$, and $r(\gamma)$ is
the sum of the ranks on the $34$ spaces.

Two rules compute these ranks. For a label $P=(\rho,k)$ with $\rho$ over
$\tau_{t}$ put $y_{P}=y_{\bbar\rho,k}$, the element of $H^{0,1}$ that carries
it. Contraction by $\partial_{\rho,k}$ is a derivation that maps $\theta_{t}$ to
$\pm y_{P}$ and the other $\theta_{t'}$ to $0$, so it maps $f(\theta_{1},
\theta_{2})$ to $\pm y_{P}\,(\partial_{t}f)(\theta_{1},\theta_{2})$. Hence a basis
element with labels $P,Q$ carries $p$ to $\pm y_{P}y_{Q}\,(Dp)(\theta_{1},
\theta_{2})$, where $D$ is $1$, $\partial_{t}$ or $\partial_{t}\partial_{t'}$
according as none, one or both labels are carried by elements of $T_{0}A$, and
$t,t'$ are their places. Multiplication by $y_{P}y_{Q}$ removes from
$\theta_{1}$ and $\theta_{2}$ the terms that contain $y_{P}$ or $y_{Q}$, and
the products $\theta_{1}^{a}\theta_{2}^{b}$ of what remains are linearly
independent as long as $a$ and $b$ do not exceed the numbers of remaining
terms. Every monomial of $y_{P}y_{Q}f(\theta_{1},\theta_{2})$ contains, with
each $x_{\sigma,i}$, its partner $y_{\bbar\sigma,i}$. The second rule is that
the images of $\omega$ below contain some $x_{\sigma,i}$ without its partner,
so they are independent of every image of $p$.

A line $y_{\sigma,k}\otimes\partial_{\bbar\sigma,k}$ carries $p$ to
$\pm y_{\sigma,k}^{2}(\cdots)=0$ and $\alpha_{\bbar\sigma}$ to
$\pm y_{\sigma,k}x_{\bbar\sigma,k'}y_{\bbar\sigma,1}y_{\bbar\sigma,2}\ne0$,
$k'\ne k$: rank one. On a space of the second kind, with labels
$P=(\rho,k)$ over $\tau_{1}$ and $Q=(\rho',l)$ over $\tau_{2}$, the four basis
elements are $y_{P}\wedge y_{Q}$, $y_{P}\otimes\partial_{\rho',l}$,
$y_{Q}\otimes\partial_{\rho,k}$ and $\partial_{\rho,k}\wedge\partial_{\rho',l}$.
They carry $\omega$ to $y_{P}y_{Q}(\alpha_{\rho}+\alpha_{\rho'})$,
$y_{P}\,(\partial_{\rho',l}\lrcorner\,\alpha_{\rho'})$,
$y_{Q}\,(\partial_{\rho,k}\lrcorner\,\alpha_{\rho})$ and $0$, the factors $y_{P}$
and $y_{Q}$ killing $\alpha_{\bbar\rho}$ and $\alpha_{\bbar\rho'}$; these
contain $x_{\rho,k'}$ or $x_{\rho',l'}$ with $k'\ne k$, $l'\ne l$, without its
partner, and the first three are independent, having two, one and one such
factors $x$ at different places. The bivector carries $\gamma$ to
$\pm y_{P}y_{Q}(\partial_{1}\partial_{2}p)(\theta_{1},\theta_{2})$, with three
terms left in each $\theta_{t}$, and $\partial_{1}\partial_{2}p$ has degree at
most three in each variable; so this image is zero exactly when
$\partial_{1}\partial_{2}p=0$: rank $3+\mu(p)$. On a space of the third kind,
with labels $P=(\sigma,k)$ and $Q=(\bbar\sigma,l)$ over $\tau_{1}$, the element
$y_{P}\wedge y_{Q}$ carries $\omega$ to $y_{P}y_{Q}(\alpha_{\sigma'}+
\alpha_{\bbar\sigma'})$, $\sigma'$ and $\bbar\sigma'$ the embeddings over
$\tau_{2}$, and the other three elements kill $\omega$: $y_{P}\otimes
\partial_{\bbar\sigma,l}$ and $y_{Q}\otimes\partial_{\sigma,k}$ carry
$\alpha_{\bbar\sigma}$ and $\alpha_{\sigma}$ into multiples of
$y_{\bbar\sigma,k}y_{\bbar\sigma,1}y_{\bbar\sigma,2}=0$ and
$y_{\sigma,l}y_{\sigma,1}y_{\sigma,2}=0$, and the bivector meets no
$\alpha$. They carry $p$ to $\pm y_{P}y_{Q}(\partial_{1}p)$, twice, and to
$\pm y_{P}y_{Q}(\partial_{1}^{2}p)$, with two terms left in $\theta_{1}$ and
four in $\theta_{2}$. So the rank is $1$ plus the dimension of the span of
$\partial_{1}p$ and $\partial_{1}^{2}p$ modulo $\theta_{1}^{3}$; in the basis
$\theta_{1}^{a}\theta_{2}^{j}/(a!\,j!)$, $a\le2$, their coefficient vectors are
$(e_{a+1,j})$ and $(e_{a+2,j})$, and the rank is $1+\rho_{1}(p)$; over
$\tau_{2}$ it is $1+\rho_{2}(p)$. The total is
$8+16(3+\mu)+4(1+\rho_{1})+4(1+\rho_{2})+R_{1}+R_{2}$.

(ii) Over $\tau_{1}$ write $x_{k},X_{k},y_{k},Y_{k}$ for $x_{\sigma,k},
x_{\bbar\sigma,k},y_{\sigma,k},y_{\bbar\sigma,k}$, so that
$\theta_{1}=\sum_{k}(x_{k}Y_{k}+X_{k}y_{k})$, $\alpha_{\sigma}=x_{1}x_{2}y_{1}y_{2}$
and $\alpha_{\bbar\sigma}=X_{1}X_{2}Y_{1}Y_{2}$, and put $c=y_{1}y_{2}$,
$C=Y_{1}Y_{2}$, $a_{k}=x_{k}y_{1}y_{2}Y_{k}$, $A_{k}=X_{k}y_{k}Y_{1}Y_{2}$,
$b=x_{1}x_{2}y_{1}y_{2}Y_{1}Y_{2}$ and $B=X_{1}X_{2}y_{1}y_{2}Y_{1}Y_{2}$. Let
$\alpha'$ be the part of $\omega$ over $\tau_{2}$, so that $\alpha'\theta_{2}=0$,
and put $\langle e_{i}\rangle=\sum_{j}e_{ij}\theta_{2}^{j}/j!$. Since $c$ kills
the terms $X_{k}y_{k}$ of $\theta_{1}$, $c\,\theta_{1}^{i}/i!$ is $c$,
$a_{1}+a_{2}$, $-b$ and $0$ for $i=0$, $1$, $2$ and $i\ge3$, so that
$c\gamma=B+c(\alpha'+\langle e_{0}\rangle)+(a_{1}+a_{2})\langle e_{1}\rangle
-b\langle e_{2}\rangle$; in the same way, with $\partial$ acting as a left
derivation, the other seven basis elements of $W_{1}$ carry $\gamma$ to
\begin{align*}
  C\gamma&=b+C(\alpha'+\langle e_{0}\rangle)+(A_{1}+A_{2})\langle e_{1}\rangle
    -B\langle e_{2}\rangle,\\
  \partial_{\sigma,1}\lrcorner\,\partial_{\sigma,2}\lrcorner\,\gamma&=
    -c+C\langle e_{2}\rangle+(A_{1}+A_{2})\langle e_{3}\rangle-B\langle e_{4}\rangle,\\
  \partial_{\bbar\sigma,1}\lrcorner\,\partial_{\bbar\sigma,2}\lrcorner\,\gamma&=
    -C+c\langle e_{2}\rangle+(a_{1}+a_{2})\langle e_{3}\rangle-b\langle e_{4}\rangle,\\
  \pm y_{k}\,\partial_{\bbar\sigma,k'}\lrcorner\,\gamma&=
    A_{k}+c\langle e_{1}\rangle+(a_{1}+a_{2})\langle e_{2}\rangle-b\langle e_{3}\rangle,\\
  \pm Y_{k}\,\partial_{\sigma,k'}\lrcorner\,\gamma&=
    a_{k}+C\langle e_{1}\rangle+(A_{1}+A_{2})\langle e_{2}\rangle-B\langle e_{3}\rangle,
\end{align*}
where $\{k,k'\}=\{1,2\}$ and the sign is $+$ for $k=1$. The differences of
the last two pairs are $A_{1}-A_{2}$ and $a_{1}-a_{2}$. In a basis containing
$A_{1}-A_{2}$, $a_{1}-a_{2}$, $c\alpha'$ and $C\alpha'$, these four
coordinates form the identity matrix on the two differences and on $c\gamma$,
$C\gamma$, and vanish on the other four images. So $R_{1}$ is $4$ plus the rank of the two
images of the bivectors and of the halves of the sums of the last two pairs,
$\tfrac12(A_{1}+A_{2})+c\langle e_{1}\rangle+(a_{1}+a_{2})\langle e_{2}\rangle
-b\langle e_{3}\rangle$ and $\tfrac12(a_{1}+a_{2})+C\langle e_{1}\rangle+
(A_{1}+A_{2})\langle e_{2}\rangle-B\langle e_{3}\rangle$. In the coordinates
along $c$, $a_{1}+a_{2}$, $b$, $C$, $-(A_{1}+A_{2})$, $B$ times
$\theta_{2}^{j}/j!$, and up to the signs of some columns, the matrix of these
four vectors is $K=\bigl(K^{0}\,\big|\,\mathrm{diag}(H,H)\bigr)$: the columns
with $j=0$ form, with $e_{k}=e_{k0}$,
\[
  K^{0}=\begin{pmatrix}
  e_{1}&e_{2}&-e_{3}&0&-\frac12&0\\
  e_{2}&e_{3}&-e_{4}&-1&0&0\\
  0&\frac12&0&e_{1}&-e_{2}&-e_{3}\\
  -1&0&0&e_{2}&-e_{3}&-e_{4}
  \end{pmatrix},
\]
and those with $j\ge1$ form $\mathrm{diag}(H,H)$, with $H$ the $2\times12$
matrix of the pairs $(e_{i,j},e_{i+1,j})$, $1\le i\le3$, $1\le j\le4$. So
$R_{1}=4+\operatorname{rk}K$. Let
$(a,b,c,d)$ annihilate the columns of $K^{0}$. The first two columns give
$d=ae_{1}+be_{2}$ and $c=-2(ae_{2}+be_{3})$, and the other four then say that
$(a,b)$ is orthogonal to the rows of
\[
  Q=\begin{pmatrix}
  e_{3}&e_{4}\\
  -e_{1}e_{2}&e_{2}^{2}-2e_{1}e_{3}-1\\
  2e_{2}^{2}-e_{1}e_{3}-\frac12&e_{2}e_{3}\\
  e_{1}e_{4}-2e_{2}e_{3}&e_{2}e_{4}-2e_{3}^{2}
  \end{pmatrix},
\]
so $\operatorname{rk}K^{0}=2+\operatorname{rk}Q$. If $Q$ were zero, then
$e_{3}=e_{4}=0$ by its first row, $e_{2}^{2}=\frac14$ by its third, and the
second row would have the entry $-\frac34$. So $\operatorname{rk}K\ge3$ and
$R_{1}\ge7$. If $R_{1}=7$, then $\operatorname{rk}Q=1$, and the $2\times2$
minors of $Q$ vanish exactly on $V_{+}\cup V_{-}$. Indeed, with
$P=2h_{1}+1$ and $m_{ij}$ the minor of rows $i$ and $j$, expanding gives
\begin{gather*}
  m_{12}=e_{2}h_{3}-e_{3}P,\quad 2m_{13}=e_{4}P-2e_{2}h_{2},\quad
  m_{14}=2e_{3}h_{2}-e_{4}h_{3},\\
  2m_{23}=e_{2}^{2}-P^{2},\quad m_{24}=Ph_{3}-e_{2}e_{3},\quad
  2m_{34}=e_{3}^{2}-h_{2}-4h_{1}h_{2}.
\end{gather*}
On $V_{\pm}$, $(P,h_{3},2h_{2})=\pm(e_{2},e_{3},e_{4})$, which makes the first
five vanish, and the last is $e_{3}^{2}-h_{2}-2(\pm e_{2}-1)h_{2}=
e_{3}^{2}+h_{2}-e_{2}e_{4}=0$. Conversely, let all six vanish. If
$e_{2}=e_{3}=e_{4}=0$ then $P=1$ and $m_{23}\ne0$; so
$(P,h_{3},2h_{2})=\lambda(e_{2},e_{3},e_{4})$ for some $\lambda$, by the first
three. If $e_{2}\ne0$, then $m_{23}=0$ gives $\lambda=\pm1$. If $e_{2}=0$, then
$P=0$, so $e_{1}e_{3}=-\frac12$, and $h_{3}=e_{1}e_{4}=\lambda e_{3}$ and
$2h_{2}=-2e_{3}^{2}=\lambda e_{4}$ give $e_{4}=-2\lambda e_{3}^{2}$ and then
$\lambda^{2}=1$. In both cases the point lies on $V_{\lambda}$. If moreover $\rho_{1}(p)\le1$, then
$h_{1}=h_{2}=h_{3}=0$, and the minors of the second and third rows and of the
first and third rows become $(e_{2}^{2}-1)/2$ and, once $e_{2}^{2}=1$,
$e_{4}/2$; so $e_{4}=0$, $e_{3}^{2}=e_{2}e_{4}=0$ and $e_{1}e_{3}=0\ne
e_{2}^{2}$, against $h_{1}=0$. Hence $\rho_{1}(p)=2$. If $\mu(p)=0$, then
$e_{ij}=0$ for $i,j\ge1$, so $H=0$ and $R_{1}=6+\operatorname{rk}Q$, which is
$7$ on $V_{\pm}$; the point $(0,1,0,0)$ of $p=\theta_{1}^{2}/2$ lies on $V_{-}$.
The parametrised points satisfy $h_{1}=\pm\frac12w\lambda^{2}-\frac14$,
$h_{2}=\pm\frac12w\lambda^{4}$ and $h_{3}=\pm w\lambda^{3}$, which are the
equations of $V_{\pm}$. The space $W_{2}$ is the same with the indices
exchanged.

(iii) A nonzero $e_{ij}$ with $i,j\ge1$ enters both Hankel matrices, so
$\mu=1$ forces $\rho_{1},\rho_{2}\ge1$, and by (ii) $R_{t}=7$ forces
$\rho_{t}=2$. The tuples $(\mu,\rho_{1},\rho_{2},R_{1},R_{2})$ allowed by
these two conditions give exactly the fifteen values, and $80$ only for
$(0,0,0,8,8)$; $\rho_{1}=\rho_{2}=0$ means that $e_{ij}=0$ unless $i=j=0$,
that is, $p$ is constant. Each value is taken. For $p=0$, $\theta_{1}$,
$\theta_{1}^{2}$, $\theta_{1}^{2}+\theta_{2}$, $\theta_{1}^{2}+\theta_{2}^{2}$,
$\theta_{1}\theta_{2}$, $\theta_{1}\theta_{2}^{2}$ and
$\theta_{1}^{2}\theta_{2}^{2}$ the points $(e_{10},\dots,e_{40})$ and
$(e_{01},\dots,e_{04})$ are $0$, $(1,0,0,0)$ or $(0,2,0,0)$, which lie on
neither $V_{\pm}$, as $2h_{1}+1$ is $1$, $1$ and $-7$ there; so $R_{1}=R_{2}=8$,
and $(\mu,\rho_{1},\rho_{2})$ is $(0,0,0)$, $(0,1,0)$, $(0,2,0)$, $(0,2,1)$,
$(0,2,2)$, $(1,1,1)$, $(1,1,2)$ and $(1,2,2)$, giving $80$, $84$, $88$, $92$,
$96$, $104$, $108$ and $112$. For $p=\theta_{1}^{2}/2+q$ with $q$ one of $0$,
$\theta_{2}$, $\theta_{2}^{2}$, $\theta_{2}^{2}/2$ the point $(0,1,0,0)$ lies on
$V_{-}$ and $\mu=0$, so $R_{1}=7$, and $R_{2}=7$ exactly for $q=\theta_{2}^{2}/2$;
this gives $87$, $91$, $95$ and $94$. For $p=\theta_{1}^{2}/2+\theta_{1}\theta_{2}
+q$ with $q$ one of $0$, $\theta_{2}^{2}$, $\theta_{2}^{2}/2$, the only nonzero
$e_{ij}$ with $i,j\ge1$ is $e_{11}=1$, so the columns of $\mathrm{diag}(H,H)$
are multiples of $(1,0,0,0)$ and $(0,0,1,0)$; at $e=(0,1,0,0)$ the columns of
$K^{0}$ span the space of $(1,0,0,0)$, $(0,0,1,0)$ and $(0,1,0,-1)$, so
$\operatorname{rk}K=3$ and $R_{1}=7$, and likewise $R_{2}=7$ for
$q=\theta_{2}^{2}/2$. With $\mu=1$ and $(\rho_{1},\rho_{2})=(2,1)$, $(2,2)$,
$(2,2)$ this gives $107$, $111$ and $110$.

(iv) The element of the torus of the proof of \Cref{thm:quarticp2prime}(ii)
that multiplies $x_{\sigma,k}$ by $s_{\sigma}$ and $y_{\sigma,k}$ by
$t_{\sigma}$, with $s_{\sigma}t_{\bbar\sigma}=\kappa_{t}$ for $\sigma$ over
$\tau_{t}$, multiplies $\theta_{t}$ by $\kappa_{t}$ and $\alpha_{\sigma}$ by
$(s_{\sigma}t_{\sigma})^{2}$, and these can be chosen to make every
$u_{\sigma}(s_{\sigma}t_{\sigma})^{2}$ equal to $1$ exactly when
$\kappa_{t}^{4}=(u_{\sigma}u_{\bbar\sigma})^{-1}$. It carries $\omega+p$ to the
normalised class with coefficients $\kappa_{1}^{i}\kappa_{2}^{j}e_{ij}$, and
contraction is equivariant. Replacing $\kappa_{t}$ by $-\kappa_{t}$ changes the
signs of $e_{1}$ and $e_{3}$, which fixes $V_{+}$ and $V_{-}$, and replacing
it by $\sqrt{-1}\,\kappa_{t}$ exchanges them; scaling the columns and the
second row of the Hankel matrices does not change their ranks.
\end{proof}

The count of (i), drawn in \Cref{fig:quarticrank}, is the analogue of the
dimension counts of \Cref{thm:p2primenumber} for a field of degree four: sixty-four of the rank
come from the Weil class alone, the Hankel terms measure how far $p$ is from
an exponential in each variable, and the two eight-dimensional spaces
$W_{t}$ carry the exceptional coefficients, the analogue of the exceptional
ratio of \Cref{rem:p2primeexceptional}. Indeed $\rho_{1}(p)\le1$ exactly when
$(e_{10},\dots,e_{40})$ lies on the cone over the twisted cubic, that is,
when the part of $p$ in $\theta_{1}$ alone is, up to a constant, a multiple
of $e^{\lambda\theta_{1}}-1$, of $\theta_{1}$ or of $\theta_{1}^{4}$.

\begin{figure}[tbp]
\centering
  \includegraphics[width=\textwidth]{fig_quarticrank}
\caption{The rank $r(\gamma)$ of \Cref{thm:quarticrank}, over the plane of
the Hankel part $16\mu+4\rho_{1}+4\rho_{2}$ and of $R_{1}+R_{2}$. The fifteen
points are the tuples allowed by the two conditions of the theorem, and each
carries its value. Filled points are reached with $\rho_{1}=\rho_{2}$, as the
character of a complex with rational coefficients is
(\Cref{cor:quarticleast}(ii)); red points need an exceptional coefficient
at one place at least. The value $80$ is the pure shape, excluded for every
CM field by \Cref{thm:cmpurefalse}; so $88$ is the least rank a complex
meeting the criterion can have. The level $r=100$ meets no point.}
\label{fig:quarticrank}
\end{figure}

\begin{corollary}[The least object the criterion asks for]
\label{cor:quarticleast}
Let $E$ be as in \Cref{thm:quarticp2prime}\textup{(v)}.
\begin{enumerate}[label=\textup{(\roman*)},leftmargin=2.6em,itemsep=2pt]
\item $\dim\Ext^{2}(E,E)\ge81$.
\item The coefficients of $P$ satisfy $e_{ji}=\iota(e_{ij})$ for the
  nontrivial automorphism $\iota$ of $F_{0}$, so $\rho_{1}=\rho_{2}$, and
  $r(\ch E)$ is one of $80$, $88$, $94$, $95$, $96$, $104$, $110$, $111$,
  $112$.
\item If $\dim\Ext^{2}(E,E)=r(\ch E)$, then $P$ is not constant and
  $\dim\Ext^{2}(E,E)$ is one of $88$, $94$, $95$, $96$, $104$, $110$, $111$,
  $112$. The value $88$ requires $\mu=0$ and $\rho_{1}=\rho_{2}=1$, and it is
  the rank of $N\omega_{1}+\eta_{s}$ for every $s\in F_{0}^{\times}$.
\end{enumerate}
\end{corollary}

\begin{proof}
(i) By \Cref{thm:quarticp2prime}(v), $\dim\Ext^{2}(E,E)\ge r(\ch E)\ge80$. If
$r(\ch E)=80$, then $P$ is a constant by \Cref{thm:quarticrank}(iii), so
$\ch(E)=c+N\omega_{1}$ and \Cref{thm:cmpurefalse} makes the inequality
strict. (ii) Write $\eta_{t}=\tau_{1}(t)\theta_{1}+\tau_{2}(t)\theta_{2}$
as before \Cref{prop:quarticquat}. A polynomial with rational coefficients in the
$\eta_{t}$ becomes, after exchanging $\theta_{1}$ and $\theta_{2}$, the
polynomial whose coefficients are the images under $\iota$ of the old ones,
so $e_{ji}=\iota(e_{ij})$; the Hankel matrix of $\rho_{2}$ is then the image
under $\iota$ of that of $\rho_{1}$, which has the same rank. The values are
those of \Cref{thm:quarticrank}(iii) with $\rho_{1}=\rho_{2}$, the scaling
of (iv) there changing neither $\mu$ nor the $\rho_{t}$. (iii) If $P$ were
constant, \Cref{thm:cmpurefalse} would give $\dim\Ext^{2}(E,E)>r(\ch E)$.
For $P=\eta_{s}$ the pairs are $(\tau_{t}(s),0)$, so $\rho_{1}=\rho_{2}=1$,
$\mu=0$, and $R_{t}=8$ by \Cref{thm:quarticrank}(ii).
\end{proof}

The Weil tori of \Cref{setup:cmweiltori} exclude more than the pure shape.
A polynomial in $\theta_{1},\theta_{2}$ can stay of Hodge type on all of
$S_{F}$, and then the argument of \Cref{thm:cmpurefalse} applies to it.

\begin{proposition}[The characters that the Weil tori exclude]
\label{prop:quarticweiltori}
Let $E$ be as in \Cref{thm:quarticp2prime}\textup{(v)}, and suppose that $P$
lies in the span of the $\theta_{1}^{i}\theta_{2}^{j}$ with
$i,j\in\{0,4\}$, as for $\ch(E)=c+N\omega_{1}+b\,\theta_{1}^{4}+
\iota(b)\,\theta_{2}^{4}$ with $b\in F_{0}$. Then $\sigma_{E}$ is not
injective and $\dim\Ext^{2}(E,E)>r(\ch E)$. Consequently, if
$\dim\Ext^{2}(E,E)=r(\ch E)=88$, then $P=c+f(\theta_{1})+f^{\iota}(\theta_{2})$
with $c\in\QQ$, where $f(\theta)$ is a nonzero multiple either of $\theta$ or
of $e^{\lambda\theta}-1$ for some $\lambda\ne0$, and $f^{\iota}$ is obtained
by applying $\iota$ to the coefficients of $f$.
\end{proposition}

\begin{proof}
In the basis before \Cref{thm:quarticrank}, $\theta_{1}$ is a sum of four
products of two classes, one from $V_{\sigma}$ and one from $V_{\bbar\sigma}$
for the two embeddings over $\tau_{1}$, in disjoint sets of basis vectors.
So $\theta_{1}^{4}$ is $4!$ times their product, a nonzero element of
$\bigwedge^{4}V_{\sigma}\otimes\bigwedge^{4}V_{\bbar\sigma}$, that is a
nonzero multiple of $\alpha_{\sigma}\alpha_{\bbar\sigma}$, and likewise for
$\theta_{2}^{4}$. Hence $\ch(E)=N\omega_{1}+P$ lies in the span of the
classes $\alpha_{\Sigma}$, which are of type $(p,p)$ on every member of
$S_{F}$ (\Cref{setup:cmweiltori}); being rational, it is a Hodge class on
every member. Its component of degree four is $N\omega_{1}\ne0$, since $P$
has none, and \Cref{thm:cmpurefalse} applies. If $\dim\Ext^{2}(E,E)=
r(\ch E)=88$, then $\mu=0$ and $\rho_{1}=\rho_{2}=1$ by
\Cref{cor:quarticleast}(iii), so $P$ has no mixed monomials, and by the
paragraph after \Cref{thm:quarticrank} its part in $\theta_{1}$ alone is a
multiple of $e^{\lambda\theta_{1}}-1$, of $\theta_{1}$ or of $\theta_{1}^{4}$;
its part in $\theta_{2}$ alone is obtained by $\iota$, by
\Cref{cor:quarticleast}(ii). In the third case $P$ lies in the span of
$1$, $\theta_{1}^{4}$ and $\theta_{2}^{4}$, which the first part excludes.
\end{proof}

The shape $c+b\,\theta_{1}^{4}+\iota(b)\,\theta_{2}^{4}$ has rank $88$, and
with a multiple of $\theta_{1}^{4}\theta_{2}^{4}$ added it has rank $104$,
by \Cref{thm:quarticrank}: $\rho_{1}=\rho_{2}=1$, $\mu$ is $0$ or $1$, and
$R_{1}=R_{2}=8$, since $(e_{10},\dots,e_{40})=(0,0,0,24b)$ with $b\ne0$ does
not meet $V_{+}\cup V_{-}$. \Cref{prop:quarticweiltori} removes
both. What remains at $88$ is a Chern character whose polynomial part is
exponential or linear in each $\theta_{t}$ separately, such as
$c+\eta_{s}+N\omega_{1}$. For these the classes $\theta_{t}$ are not of Hodge
type on the Weil tori, and the argument needs a family of complex tori along
which the character stays of Hodge type and whose members are not algebraic.
The next proposition, the analogue for a quartic field of
\Cref{prop:p2primelocus}(iii) and (iv), shows that there is none: the Hodge
locus of such a character is the polarised family itself.

\begin{proposition}[The Hodge locus of a character at a quartic field]
\label{prop:quarticlocus}
Keep the notation before \Cref{thm:quarticrank}, with $\gamma=\omega+p$ and
every $u_{\sigma}\ne0$, and for the embeddings $\sigma,\bbar\sigma$ over
$\tau_{t}$ put $\omega_{t}=u_{\sigma}\alpha_{\sigma}+u_{\bbar\sigma}
\alpha_{\bbar\sigma}$. Call $\tau_{t}$ a \emph{mid place} of $p$ if $p$ has a
monomial $\theta_{1}^{i}\theta_{2}^{j}$ with nonzero coefficient whose exponent
of $\theta_{t}$ is $1$, $2$ or $3$, and call a mid place \emph{exceptional} if
the only such monomial is $\theta_{t}^{2}$ itself and its coefficient $c_{t}$
in $p$ satisfies $c_{t}^{2}\theta_{t}^{4}=\frac34\,\omega_{t}^{2}$, the ratio
of \Cref{prop:p2primelocus}(iv).
\begin{enumerate}[label=\textup{(\roman*)},leftmargin=2.6em,itemsep=2pt]
\item The annihilator of $\gamma$ in $H^{1}(T_{A})$ is the space of $F$-linear
  $\xi\in\bigoplus_{\sigma}\Hom(V_{\sigma}^{1,0},V_{\sigma}^{0,1})$ with
  $\xi\lrcorner\theta_{t}=0$ at every mid place, plus, at each exceptional
  place, a line spanned by an element of $\Hom(V_{\sigma}^{1,0},
  V_{\bbar\sigma}^{0,1})\oplus\Hom(V_{\bbar\sigma}^{1,0},V_{\sigma}^{0,1})$
  with both components nonzero. Its dimension is $8+4a+b$, where $a$ is the
  number of places that are not mid places and $b$ the number of exceptional
  places.
\item The connected component through $A$ of the Hodge locus of $\gamma$, in
  the space of complex structures on the real torus underlying $A$, is
  $S_{F}$ when neither place is a mid place, that is when $p$ lies in the
  span of the $\theta_{1}^{i}\theta_{2}^{j}$ with $i,j\in\{0,4\}$; and it is
  $\cD_{F}$ when both places are mid places and neither is exceptional. In
  the second case all its members are abelian varieties polarised by
  $\theta_{1}+\theta_{2}$.
\item If both places are exceptional, then near $A$ the Hodge locus is the
  orbit of $A$ under $G_{1}\times G_{2}$, where $G_{t}\cong\mathrm{Spin}(4,3)$
  is the connected stabiliser of $X_{t}=\omega_{t}+c_{t}\theta_{t}^{2}$ in the
  real linear group of $V_{\sigma}\oplus V_{\bbar\sigma}$, and it has
  dimension $10$. A very general member $T$ of it has $\NS(T)=0$ and
  $\Hdg^{2}(T)\otimes\CC$ contained in the span of $X_{1}$ and $X_{2}$, and on
  every member, for every K\"ahler class $\kappa$, $\int X_{t}\,\kappa^{6}$ is
  nonzero with the sign of $c_{t}$, in the normalisation in which
  $\theta_{t}$ is a real class with $\theta_{t}^{2}$ positive.
\item If $\gamma$ is the complexified character of \Cref{thm:quarticp2prime}(v),
  then both places are mid places or neither is, and both are exceptional or
  neither is; and an exceptional $\gamma$ has $r(\gamma)=94$, or $110$ when
  $p$ contains $\theta_{1}^{4}\theta_{2}^{4}$.
\end{enumerate}
\end{proposition}

\begin{proof}
(i) Grade $\bigwedge^{*}H^{1}$ by the degrees in the four $V_{\sigma}$, and
write $\xi_{\sigma\sigma'}$ for the component of $\xi$ in
$\Hom(V_{\sigma'}^{1,0},V_{\sigma}^{0,1})$, which raises the degree in
$V_{\sigma}$ by one and lowers that in $V_{\sigma'}$ by one. The monomial
$\theta_{1}^{i}\theta_{2}^{j}$ has degree $i$ in the two embeddings over
$\tau_{1}$ and $j$ in those over $\tau_{2}$, and $\alpha_{\sigma}$ has degree
$4$ in $V_{\sigma}$. By the proof of \Cref{thm:quarticp2prime}(i),
$\xi_{\sigma\sigma}\lrcorner\omega=0$, and for $\sigma\ne\sigma'$ the only
term of $\xi_{\sigma\sigma'}\lrcorner\omega$ is
$u_{\sigma'}\xi_{\sigma\sigma'}\lrcorner\alpha_{\sigma'}$, of degree $3$ in
$V_{\sigma'}$ and $1$ in $V_{\sigma}$, which determines $\xi_{\sigma\sigma'}$.
A term $\xi_{\rho\rho'}\lrcorner\theta_{1}^{i}\theta_{2}^{j}$ has that degree
only if $\sigma,\sigma'$ lie over the same place $\tau_{t}$, $(i,j)$ is
$(2,0)$ or $(0,2)$ with the $2$ at $t$, and $(\rho,\rho')=(\sigma',\sigma)$.
So every component $\xi_{\sigma\sigma'}$ with $\sigma,\sigma'$ over different
places vanishes, and the equations split into the $F$-linear part and, at
each place, the pair $\xi'=\xi_{\bbar\sigma\sigma}$, $\xi''=\xi_{\sigma\bbar\sigma}$.

For $F$-linear $\xi$ put $D_{t}=\xi\lrcorner\theta_{t}\in V_{\sigma}^{0,1}
\otimes V_{\bbar\sigma}^{0,1}$; the map $\xi\mapsto(D_{1},D_{2})$ is onto, with
kernel $T_{[A]}\cD_{F}$. Then $\xi\lrcorner\theta_{1}^{i}\theta_{2}^{j}
=i\,\theta_{1}^{i-1}\theta_{2}^{j}D_{1}+j\,\theta_{1}^{i}\theta_{2}^{j-1}D_{2}$,
the monomials lie in different degrees, and the two terms have different
Hodge types at the two places, so each vanishes. Now $\theta_{t}^{3}D_{t}=0$
for every $D_{t}$, as a degree count shows, while multiplication by
$\theta_{t}^{k}$ is injective on $V_{\sigma}^{0,1}\otimes V_{\bbar\sigma}^{0,1}$
for $k\le2$: $\theta_{t}^{2}y_{\sigma,k}y_{\bbar\sigma,l}$ is $2$ times the
product of $y_{\sigma,k}y_{\bbar\sigma,l}$ with the two terms of $\theta_{t}$
that avoid $y_{\sigma,k}$ and $y_{\bbar\sigma,l}$, and different $(k,l)$ give
different monomials. Multiplication by $\theta_{t'}^{j}$, $j\le4$, acts on the
other factor of $\bigwedge^{*}H^{1}=\bigwedge^{*}(V_{\sigma}\oplus
V_{\bbar\sigma})\otimes\bigwedge^{*}(V_{\sigma'}\oplus V_{\bbar\sigma'})$ and is
nonzero there. So the conditions on the $F$-linear part are exactly $D_{t}=0$
at the mid places, a space of dimension $16-4(2-a)=8+4a$.

At a place $\tau_{t}$ let $c_{t}$ be the coefficient of $\theta_{t}^{2}$ in $p$.
The terms involving $\xi''$ are $u_{\bbar\sigma}\xi''\lrcorner
\alpha_{\bbar\sigma}+c_{t}\,\xi'\lrcorner\theta_{t}^{2}$, of degree $3$ in
$V_{\bbar\sigma}$ and $1$ in $V_{\sigma}$, the same with $\sigma$ and
$\bbar\sigma$ and with $\xi'$ and $\xi''$ exchanged, and
$i\,\theta_{t}^{i-1}\theta_{t'}^{j}(\xi'\lrcorner\theta_{t})$ and its analogue
for $\xi''$ for every other monomial with exponent $i\in\{1,2,3\}$ of
$\theta_{t}$. Here $\xi'\lrcorner\theta_{t}$ lies in the line
$\bigwedge^{2}V_{\bbar\sigma}^{0,1}$, on which multiplication by
$\theta_{t}^{i-1}\theta_{t'}^{j}$ is injective for $i\le3$ by the same count. If
such another monomial occurs, $\xi'\lrcorner\theta_{t}=\xi''\lrcorner
\theta_{t}=0$, and then $\xi'=\xi''=0$, the maps $\xi'\mapsto
\xi'\lrcorner\alpha_{\sigma}$ and $\xi''\mapsto\xi''\lrcorner\alpha_{\bbar\sigma}$
being injective. Otherwise $\xi''\lrcorner\theta_{t}^{2}=2\theta_{t}\,
(\xi''\lrcorner\theta_{t})$ is a multiple of $\theta_{t}\,y_{\sigma,1}
y_{\sigma,2}$, which is $\xi'_{0}\lrcorner\alpha_{\sigma}$ for one element
$\xi'_{0}$, and likewise with $\sigma$ and $\bbar\sigma$ exchanged; so $\xi'$
and $\xi''$ are multiples of $\xi'_{0}$ and $\xi''_{0}$. Explicitly,
$\xi'_{0}(x_{\sigma,1})=-y_{\bbar\sigma,2}$ and
$\xi'_{0}(x_{\sigma,2})=y_{\bbar\sigma,1}$ give
$\xi'_{0}\lrcorner\alpha_{\sigma}=\theta_{t}\,y_{\sigma,1}y_{\sigma,2}$ and
$\xi'_{0}\lrcorner\theta_{t}=2y_{\bbar\sigma,1}y_{\bbar\sigma,2}$, and
symmetrically for $\xi''_{0}$. Writing $\xi'=a\xi'_{0}$ and
$\xi''=b\xi''_{0}$, the two remaining equations are
$u_{\sigma}a+4c_{t}b=0$ and $4c_{t}a+u_{\bbar\sigma}b=0$, and their
determinant vanishes exactly when $16c_{t}^{2}=u_{\sigma}u_{\bbar\sigma}$. Since $\theta_{t}^{4}=24\,
\alpha_{\sigma}\alpha_{\bbar\sigma}$ and $\omega_{t}^{2}=2u_{\sigma}
u_{\bbar\sigma}\alpha_{\sigma}\alpha_{\bbar\sigma}$, that is the condition
$c_{t}^{2}\theta_{t}^{4}=\frac34\omega_{t}^{2}$, and the solutions then form a
line on which both components are nonzero.

(ii) By \Cref{prop:p2primelocus}(i), applied to each homogeneous component,
the Hodge locus of $\gamma$ in the chart at $A$ is the linear space
$\Ann_{H^{1}(T_{A})}(\gamma)$. If neither place is a mid place, that space is
the space of all $F$-linear $\xi$, which is also the part of $S_{F}$ in the
chart; $\gamma$ is of Hodge type on all of $S_{F}$, as in the proof of
\Cref{prop:quarticweiltori}, and at every member of $S_{F}$ its component of
degree four is $\omega$, whose annihilator in $H^{1}(T)$ is again the space of
$F$-linear deformations (\Cref{thm:quarticp2prime}(i), whose proof uses only
the decomposition into the $V_{\sigma}$). So $S_{F}$ is open in the Hodge locus;
it is closed in the space of complex structures, since commuting with $F$ is a
closed condition and the dimensions of the $H^{1,0}\cap V_{\sigma}$ are
locally constant on it; and it is connected. If both places are mid places
and neither is exceptional, the annihilator is $T_{[A]}\cD_{F}$ by (i), at
every member of $\cD_{F}$, since the conditions of (i) are stated in terms of
the flat classes $\theta_{t}$ and $\omega_{t}$; and the argument of
\Cref{prop:p2primelocus}(iii) shows that $\cD_{F}$ is open, closed and
connected in the Hodge locus.

(iii) At each place $X_{t}$ is, in the $8$-dimensional space
$V_{\sigma}\oplus V_{\bbar\sigma}$ with the form $\theta_{t}$, the class of
\Cref{rem:p2primeexceptional}, and the argument there applies word for
word, with the factor $\mathfrak{su}(2,2)$ at $\tau_{t}$ of the Lie algebra of $G_{F}$
in place of the Hodge group: the stabiliser is $\mathfrak{so}(4,3)$, with an
orbit of dimension $5$ in the Grassmannian. The other components of $\gamma$ are constants, multiples of
$\theta_{t}^{4}$, which is $24$ times the volume form of
$V_{\sigma}\oplus V_{\bbar\sigma}$, and of $\theta_{1}^{4}\theta_{2}^{4}$, all
fixed by $G_{1}\times G_{2}$. So the orbit of $A$ lies in the Hodge locus,
and it has dimension $10$, which by (i) is that of the locus near $A$. A
rational class of Hodge type at a very general member is of Hodge type on the
whole orbit, hence killed by the complex structures of all its members,
which lie in $\mathfrak{g}_{1}\oplus\mathfrak{g}_{2}$ and generate an ideal
projecting nontrivially to both simple factors, that is all of it. The
invariants of $\mathfrak{so}(7,\CC)$ in the exterior algebra of its spin
representation are $1$, $X_{t}$ and the volume form
(\Cref{rem:p2primeexceptional}), so in degree two there are none and in
degree four only $X_{1}$ and $X_{2}$. For the sign, a K\"ahler class is
$\kappa_{h}$ for a positive hermitian form $h$ on the tangent space; only the
terms of $\kappa_{h}^{6}$ that contain the whole volume of the other place pair
with $X_{t}$, and they give $\det(h_{t'})$ times $\kappa_{h'}^{2}$ for the
Schur complement $h'$ of $h_{t'}$ in $h$, a positive hermitian form on the
tangent space at $\tau_{t}$. So $\int X_{t}\kappa_{h}^{6}$ is a positive
multiple of the integral of \Cref{rem:p2primeexceptional} for $h'$, which has
the sign of $c_{t}$; and $G_{1}\times G_{2}$ fixes $X_{t}$ and carries K\"ahler
cones to K\"ahler cones.

(iv) An automorphism of $\bbar\QQ$ inducing $\iota$ on $F_{0}$ fixes the
rational class $\gamma$ and exchanges $\theta_{1}$ with $\theta_{2}$ and
$\omega_{1}$ with $\omega_{2}$, so it carries the conditions at $\tau_{1}$ to
those at $\tau_{2}$. An exceptional $p$ lies in the span of $1$,
$\theta_{t}^{2}$, $\theta_{t}^{4}$ and $\theta_{1}^{4}\theta_{2}^{4}$ with
$c_{t}\ne0$, so $\rho_{1}=\rho_{2}=2$, $\mu=1$ exactly when
$\theta_{1}^{4}\theta_{2}^{4}$ occurs, and $(e_{10},\dots,e_{40})$ lies on
$V_{+}\cup V_{-}$, so $R_{1}=R_{2}=7$ by \Cref{thm:quarticrank}(ii); the
formula of \Cref{thm:quarticrank}(i) gives $94$ or $110$.
\end{proof}

\begin{corollary}[The Weil tori at rank $88$]\label{cor:quarticlocus}
Let $E$ be as in \Cref{thm:quarticp2prime}\textup{(v)} with
$\dim\Ext^{2}(E,E)=r(\ch E)=88$, so that $\ch(E)=N\omega_{1}+c+f(\theta_{1})
+f^{\iota}(\theta_{2})$ as in \Cref{prop:quarticweiltori}. Then the connected
component through $A$ of the Hodge locus of $\ch(E)$ is $\cD_{F}$, and every
complex torus reached from $A$ by a deformation along which $\ch(E)$ stays of
Hodge type is an abelian variety of the family. So no argument that deforms
$E$ itself to a torus that is not algebraic can exclude it; the characters
that such an argument excludes at $88$ are exactly those of
\Cref{prop:quarticweiltori}.
\end{corollary}

\begin{proof}
The coefficient $e_{10}$ of $f$ is nonzero, both for a multiple of $\theta$
and for one of $e^{\lambda\theta}-1$, and $e_{01}=\iota(e_{10})$, so both
places are mid places, with the monomials $\theta_{1}$ and $\theta_{2}$, and
neither is exceptional. \Cref{prop:quarticlocus}(ii) applies.
\end{proof}

As after \Cref{prop:p2primelocus}, this concerns $E$ itself: an
autoequivalence of $D^{b}(A)$ changes the character, and the proposition then
applies to the new one. For the exceptional characters of
\Cref{prop:quarticlocus}(iii), of rank $94$ or $110$, the proof of
\Cref{thm:cmpurefalse} can be started, since the locus contains tori without
divisors, but it breaks down as in \Cref{rem:p2primeexceptional}: the classes
$X_{t}$ have integrals of constant sign against every K\"ahler class, so
Bando and Siu give no vanishing: they bound the signs of the coefficients of
$\ch_{2}$ of a torsion-free sheaf along $X_{1}$ and $X_{2}$, which a complex,
an alternating sum, need not respect. So at a
quartic field the Weil tori and their exceptional analogues exclude exactly
the characters of \Cref{prop:quarticweiltori}, and what the criterion still
asks for, at $88$ and above, is the construction of an object.

The natural candidates are Markman's secant objects \cite{Mar25c}, in the
twisted form in which he states his reduction. For them a $B$-field twist of
$\ch(E)$, the class $\kappa$ of \Cref{thm:quarticobstruction}, is flat, and
under a non-degeneracy condition its component in $H^{4}$ is a polynomial in
the $\eta_{t}$ plus a nonzero Weil class
\cite[Proposition~10.2.1, Theorem~1.1.2]{Mar25c}. But $\ch(E)$ itself is not
flat, and $\kappa$ is not of the form $N\omega+P(\eta)$ in higher degrees
(\Cref{rem:quarticsecantshape}), so they are candidates only for a twisted
form of \Cref{thm:quarticp2prime}(v). The dimension count that works for
imaginary quadratic fields (\Cref{thm:twistedcert}) can be tested on them
exactly. Let $X$ be a principally polarised abelian fourfold with real
multiplication by the ring of integers $\cO_{F_{0}}$, write
$\theta_{1},\theta_{2}\in H^{2}(X,\RR)$ for the components of the
polarisation at the two real places of $F_{0}$, so that
$\theta=\theta_{1}+\theta_{2}$ and $\theta_{j}^{3}=0$, and let $\mathcal{R}$ be
the $\QQ$-subalgebra of $H^{\mathrm{ev}}(X,\QQ)$ spanned by the rational
combinations of the $\theta_{1}^{a}\theta_{2}^{b}$, $0\le a,b\le2$. A class
$v=\sum_{a,b}V_{ab}\,\theta_{1}^{a}\theta_{2}^{b}/(a!\,b!)$ of $\mathcal{R}$ has a
$3\times3$ matrix $V$ with entries in $F_{0}\otimes\RR$.

\begin{proposition}[The secant space of a quartic field]\label{prop:quarticsecant}
Let $F=F_{0}(\sqrt{-q})$ with $q\in F_{0}$ totally positive, let $t\in F_{0}$,
$s_{j}=\sqrt{-1}\,\sqrt{\tau_{j}(q)}$, and for $\varepsilon\in\{\pm1\}^{2}$ put
$e_{\varepsilon}=\exp\bigl(\sum_{j}(\tau_{j}(t)+\varepsilon_{j}s_{j})\theta_{j}\bigr)$.
\begin{enumerate}[label=\textup{(\roman*)},leftmargin=2.6em,itemsep=2pt]
\item The span $S(t,q)$ of the four $e_{\varepsilon}$ is a rational subspace of
  $\mathcal{R}$ of dimension $4=2^{[F_{0}:\QQ]}$, secant to the spinor variety.
  For $t=0$ it is the secant space of \cite[\S11.1, (11.1.2)]{Mar25c} for the
  polarisation $\theta$ and $q$; replacing $\theta$ by $f\theta$,
  $f\in F_{0}^{\times}$, replaces $q$ by $qf^{2}$, and a rational $F_{0}$-$B$-field
  $t$ transports Markman's $F$-structure by the rational automorphism of
  $X\times\widehat X$ acting on $H^{*}(X)$ as $e^{\theta_{t}}$,
  $\theta_{t}=\sum_{j}\tau_{j}(t)\theta_{j}$, and replaces $S(0,q)$ by
  $S(t,q)=e^{\theta_{t}}S(0,q)$.
\item For $v=\sum_{\varepsilon}w_{\varepsilon}e_{\varepsilon}\in S(t,q)$ let
  $\rho=\rk V$ and $N_{w}=\#\{\varepsilon:w_{\varepsilon}\ne0\}$. The ranks of
  contraction from $HT^{k}(X)$ into $v$ are $1$, $8$ and $2\rho+4N_{w}$ for
  $k=0,1,2$, with $(\rho,N_{w})\in\{(1,4),(2,4),(2,2)\}$, and $N_{w}=2$ only for
  biquadratic $F$ and $v$ in the secant plane of one of its two imaginary
  quadratic subfields \cite[\S11.2.1]{Mar25c}.
\item For $v=e^{\theta_{t}}(a\,u+v_{f}+c\,w)$ in the basis
  $u=(1-\tfrac12q_{1}\theta_{1}^{2})(1-\tfrac12q_{2}\theta_{2}^{2})$,
  $v_{f}=\theta_{f}-\tfrac12\theta_{\bar fq}\theta_{1}\theta_{2}$,
  $w=\theta_{1}\theta_{2}$ of $e^{-\theta_{t}}S(t,q)$, one has
  $\chi(v,v)=\int_{X}v^{\vee}v=4\bigl(\Nm(q)a^{2}+c^{2}+\Tr(\bar q f^{2})\bigr)$,
  which is positive definite.
\end{enumerate}
\end{proposition}

\begin{proof}
(i) The Galois group permutes the four exponentials as it permutes the CM
types of $F$, so their span is rational; since $\theta_{j}^{3}=0$,
$e^{\pm s_{j}\theta_{j}}=(1-\tfrac12\tau_{j}(q)\theta_{j}^{2})\pm s_{j}\theta_{j}$, and
the four classes are independent. The conjugates in
$\mathfrak{so}\cap\End^{0}(X\times\widehat X)$ of the standard structure have
eigenvalues $t\pm s$, and the pure spinors of their eigenspaces are the
$e_{\varepsilon}$.

(ii) Over $\CC$, $H^{1}(X)=U_{1}\oplus U_{2}$ with $U_{j}$ the
eigenspaces of $F_{0}$, and $HT^{*}(X)$ and $H^{*}(X)$ are tensor products over
$j$, compatibly with contraction. By \Cref{lem:bfield}(i),
$\xi\lrcorner e^{\lambda\theta_{j}}=y^{\lambda}(\xi)e^{\lambda\theta_{j}}$, and the
blocks $HT_{1}^{a}\otimes HT_{2}^{b}$ have independent images because they
shift the two Hodge bidegrees by $(a,b)$ while $v$ is balanced. The block
$(1,1)$ contributes $4N_{w}$, by the invertibility of the Hadamard matrix of
the signs; the blocks $(2,0)$ and $(0,2)$ contribute $\rho$ each; the blocks
$(1,0)$ and $(0,1)$ contribute $4$ each in degree one. Reality gives
$w_{-\varepsilon}=\overline{w_{\varepsilon}}$, and when $F$ is not biquadratic
the Galois group acts transitively on the CM types, so $N_{w}=4$.

(iii) One has
$\int_{X}e^{c_{1}\theta_{1}+c_{2}\theta_{2}}=c_{1}^{2}c_{2}^{2}$, and $\chi$ is
invariant under $e^{\theta_{t}}$; the form is positive definite because $q$ is
totally positive.
\end{proof}

The integral classes of $\mathcal{R}$ are described by one congruence, and
the description does not depend on the field.

\begin{lemma}[The integral classes of $\mathcal{R}$]\label{lem:quarticlattice}
Let $F_{0}$ be a real quadratic field with ring of integers $\cO$ and
different $\mathfrak d$, and let $X$ be a principally polarised abelian
fourfold with real multiplication by $\cO$. A class $v\in\mathcal{R}$ with
matrix
\[
  V=\begin{pmatrix}r&\bar h&\bar g\\ h&c&\bar k\\ g&k&m\end{pmatrix},
  \qquad r,c,m\in\QQ,\quad h,g,k\in F_{0},
\]
lies in $H^{\mathrm{ev}}(X,\ZZ)$ if and only if $r,c,m\in\ZZ$,
$h,g,k\in\cO$ and $c\equiv g$ modulo $\mathfrak d$.
\end{lemma}

\begin{proof}
Integrality can be tested one prime $p$ at a time, and after a faithfully
flat extension of $\ZZ_{p}$. Let $M=H_{1}(X,\ZZ)$ and let $\langle\,,\rangle$
be the polarisation form on $M$, that is, $\theta$. The Rosati involution fixes
$\cO$, so $\langle\,,\rangle=\Tr_{F_{0}/\QQ}\psi$ for a unique $\cO$-bilinear
alternating form $\psi$ with values in $\mathfrak d^{-1}$, and $\psi$ is
perfect because $\theta$ is principal. Write $\cO=\ZZ[\nu]$ and
$\delta=\nu-\bar\nu$, a generator of $\mathfrak d$
with $\bar\delta=-\delta$, and put $\cO_{p}=\cO\otimes\ZZ_{p}$, a product of
at most two discrete valuation rings. The form $\delta\psi$ on
$M_{p}=M\otimes\ZZ_{p}$ is perfect with values in $\cO_{p}$, so $M_{p}$ has a
symplectic $\cO_{p}$-basis $e,f,e^{*},f^{*}$. Hence
$M_{p}=\cO_{p}\otimes_{\ZZ_{p}}W$, with $W$ the $\ZZ_{p}$-span of this basis,
and $\langle\,,\rangle=B\otimes\omega_{0}$, where $\omega_{0}$ is the standard
symplectic form of $W$ and $B(\alpha,\beta)=\Tr(\delta^{-1}\alpha\beta)$.

Let $\Lambda$ be the ring of integers of a finite extension of $\QQ_{p}$ over
which $F_{0}$ splits, and fix an embedding of $F_{0}$ in its field of
fractions. The map $\alpha\otimes\lambda\mapsto(\alpha\lambda,\bar\alpha\lambda)$
identifies $\cO_{p}\otimes_{\ZZ_{p}}\Lambda$ with the lattice of pairs
$(x,y)\in\Lambda^{2}$ with $x\equiv y$ modulo $\delta$, whose basis is
$(1,1),(\delta,0)$. Let $\varepsilon_{1},\varepsilon_{2}$ be the coordinates
and $u_{1},\dots,u_{4}$ the basis of $W^{\vee}$ dual to $e,f,e^{*},f^{*}$. The
dual lattice has basis $\varepsilon_{2}$, $(\varepsilon_{1}-\varepsilon_{2})/\delta$,
so $H^{1}(X,\Lambda)$ has the basis
\[
  Y_{i}=\varepsilon_{2}\otimes u_{i},\qquad
  Z_{i}=\delta^{-1}(X_{i}-Y_{i}),\qquad\text{where }X_{i}=\varepsilon_{1}\otimes u_{i},
\]
and $H^{*}(X,\Lambda)$ is the exterior algebra on the $Y_{i}$ and $Z_{i}$. In
these coordinates $B=\delta^{-1}(\varepsilon_{1}^{2}-\varepsilon_{2}^{2})$, and the
components of $\theta=B\otimes\omega_{0}$ in the two eigenspaces of $\cO$, the
first being the one on which $\cO$ acts through the fixed embedding, are
\[
  \theta_{1}=\delta^{-1}(X_{1}X_{3}+X_{2}X_{4}),\qquad
  \theta_{2}=-\delta^{-1}(Y_{1}Y_{3}+Y_{2}Y_{4}).
\]
Substitute $X_{i}=Y_{i}+\delta Z_{i}$ in $v=\sum_{a,b}V_{ab}\,\theta_{1}^{a}
\theta_{2}^{b}/(a!\,b!)$ and expand in the monomials in the $Y_{i}$ and $Z_{i}$.
Up to sign, every coefficient is one of
\[
  r,\ m,\ h,\ k,\ c,\ g,\ g-c,\ \delta h,\ \delta k,\ \delta g,\ \delta^{2}g,\
  \frac{g-c}{\delta},\ \frac{h-\bar h}{\delta},\ \frac{k-\bar k}{\delta},\
  \frac{2c-g-\bar g}{\delta^{2}},
\]
and each of $r$, $m$, $h$, $k$, $c$, $g$ and $(g-c)/\delta$ occurs, so an
integral class satisfies the conditions at $p$. Conversely, suppose that they
hold. Every $x\in\cO_{p}$ has $x-\bar x\in\delta\cO_{p}$, which disposes of
$h-\bar h$ and $k-\bar k$; and $g=c+\delta g'$ with $g'\in\cO_{p}$ gives
$2c-g-\bar g=-\delta(g'-\bar g')\in\delta^{2}\Lambda$. So every coefficient
lies in $\Lambda$. When $p$ is unramified in $F_{0}$, $\delta$ is a unit at
$p$ and the congruence is empty.
\end{proof}

\begin{proposition}[Integral secant classes of small Euler characteristic]
\label{prop:quarticother}
Let $F_{0}=\QQ(\sqrt D)$ with $D>1$ squarefree, let $X$ be a principally
polarised abelian fourfold with real multiplication by $\cO_{F_{0}}$, and let
$v$ be an integral class in a space $S(t,q)$ of \Cref{prop:quarticsecant},
with $\rho$ and $N_{w}$ as there.
\begin{enumerate}[label=\textup{(\roman*)},leftmargin=2.6em,itemsep=2pt]
\item If $\rho=1$, then $\chi(v,v)\ne4$.
\item If $\rho=2$ and $N_{w}=4$, then $\chi(v,v)\ne6$.
\end{enumerate}
\end{proposition}

\begin{proof}
By \Cref{lem:quarticlattice} the entries of $V$ lie in $\cO=\cO_{F_{0}}$, the
diagonal ones in $\ZZ$, and $c\equiv g$ modulo $\mathfrak d$; conjugation is
trivial on $\cO/\mathfrak d$.

(i) Write $V=\lambda\alpha\bar\alpha^{T}$ with $\alpha\in\cO^{3}$, and let
$\mathfrak a$ be the ideal generated by the entries of $\alpha$; replacing
$\alpha$ by a multiple, we may take $\mathfrak a$ integral and prime to
$\mathfrak d$. The entries $\lambda\alpha_{i}\bar\alpha_{j}$ lie in $\cO$ and
generate $\lambda\mathfrak a\bar{\mathfrak a}=\lambda\Nm(\mathfrak a)\cO$, so
$\mu=\lambda\Nm(\mathfrak a)\in\ZZ$. With $D'=\alpha_{1}^{2}-\alpha_{0}\alpha_{2}
\in\mathfrak a^{2}$ one has $\chi=4\lambda^{2}\Nm(D')$, and
$\Nm(D')\in\Nm(\mathfrak a)^{2}\ZZ$, so $\chi=4$ forces $\mu=\pm1$ and
$(D')=\mathfrak a^{2}$. Since conjugation is trivial modulo $\mathfrak d$,
$\alpha_{1}\bar\alpha_{1}-\alpha_{2}\bar\alpha_{0}\equiv D'$ modulo
$\mathfrak a\mathfrak d$, and the congruence $c\equiv g$ gives
$\lambda D'\in\mathfrak d\,\bar{\mathfrak a}^{-1}$. The ideal of $\lambda D'$ is
$\mathfrak a\bar{\mathfrak a}^{-1}$, so $\mathfrak a\subseteq\mathfrak d$, which
is impossible for $\mathfrak a$ prime to $\mathfrak d$ and $\mathfrak d\ne(1)$.

(ii) The matrix $\operatorname{adj}V$ is Hermitian of rank one, so
$\operatorname{adj}V=\mu\,\bar l\,l^{T}$ with $l\in\cO^{3}$, $l^{T}V=0$ and
$\mu\in\QQ^{\times}$. Put $\Delta=l_{1}^{2}-4l_{0}l_{2}$ and
$m=\mu^{2}\Nm(\Delta)$. Replacing $l$ by $\lambda l$, $\lambda\in F_{0}^{\times}$,
replaces $\mu$ by $\mu/\Nm(\lambda)$ and $\Delta$ by $\lambda^{2}\Delta$, so $m$
depends only on $V$. As the entries of $\operatorname{adj}V$ are integral,
$\mu$ is integral at every prime at which $l$ is primitive, and $l$ can be
made primitive at the primes above any given $p$, because
$\cO\otimes\ZZ_{(p)}$ is a principal ideal domain; so $m\in\ZZ$. Let $J$ be
the matrix with rows $(0,0,1)$, $(0,-2,0)$, $(1,0,0)$. Since
$\int_{X}\theta_{1}^{2}\theta_{2}^{2}=4$, one has $\chi(v,v)=\operatorname{tr}
(JVJ\bar V)$; and $\operatorname{adj}J$ has rows $(0,0,2)$, $(0,-1,0)$,
$(2,0,0)$, so that $l^{T}(\operatorname{adj}J)\,l=-\Delta$ and
\[
  \sigma_{2}(JVJ\bar V)=\operatorname{tr}\operatorname{adj}(JVJ\bar V)
  =\mu^{2}\,\Delta\bar\Delta=m .
\]
At a real place, $V=PWQ^{T}$ with $W=(w_{\varepsilon})$ and $P$, $Q$ the
$3\times2$ matrices with columns $(1,\pm s_{1},s_{1}^{2})$ and
$(1,\pm s_{2},s_{2}^{2})$, when $t=0$; multiplication by $e^{\theta_{t}}$
replaces $V$ by $UV\bar U^{T}$ with $U$ unipotent and $U^{T}JU=J$, which
conjugates $JVJ\bar V$, so we may take $t=0$. As
$(1,z,z^{2})\,J\,(1,z',z'^{2})^{T}=(z-z')^{2}$ and
$w_{-\varepsilon}=\overline{w_{\varepsilon}}$, the nonzero eigenvalues of
$JVJ\bar V$ are $16\Nm(q)\bigl(|w_{++}|\pm|w_{+-}|\bigr)^{2}$, so
\[
  \chi^{2}-4m=\bigl(64\Nm(q)\,|w_{++}|\,|w_{+-}|\bigr)^{2}\ge0,
\]
with equality exactly when $N_{w}=2$. Modulo $\mathfrak d$ one has
$\bar V\equiv V$, $\operatorname{tr}(JV)=g-2c+\bar g\equiv0$ and
$\sigma_{2}(JV)=\operatorname{tr}\bigl(\operatorname{adj}(V)\operatorname{adj}(J)\bigr)
=\mu\,l^{T}(\operatorname{adj}J)\,\bar l\equiv-\mu\Delta$, with $l$ primitive at
the primes dividing $\mathfrak d$. Hence
\[
  \chi\equiv\operatorname{tr}\bigl((JV)^{2}\bigr)
  =\operatorname{tr}(JV)^{2}-2\sigma_{2}(JV)\equiv2\mu\Delta
  \pmod{\mathfrak d}.
\]

Now let $\chi=6$ and $N_{w}=4$. Then $1\le m\le8$, because $\Delta$ is
totally negative, and $m$ satisfies three conditions.
\begin{enumerate}[label=\textup{(\alph*)},leftmargin=2.6em,itemsep=2pt]
\item $4m\equiv36$ modulo $\mathfrak d\cap\ZZ$, which is $D\ZZ$ if
  $D\equiv1$ modulo $4$ and $2D\ZZ$ otherwise. As $D$ is squarefree, this
  says that $m\equiv9$ modulo the odd part $D_{o}$ of $D$.
\item $m\equiv0$ or $1$ modulo $4$. Take $l$ primitive at $2$, so that
  $\mu\in\ZZ_{(2)}$. Since $\Nm(x+4y)\equiv\Nm(x)$ modulo $4$ for
  $x,y\in\cO\otimes\ZZ_{(2)}$, $\Nm(\Delta)\equiv\Nm(l_{1}^{2})=\Nm(l_{1})^{2}$,
  and $m\equiv\bigl(\mu\Nm(l_{1})\bigr)^{2}$ modulo $4$.
\item If $D\equiv2$ modulo $4$, then $m$ is odd. The prime $\mathfrak p$ above
  $2$ is ramified, $v_{\mathfrak p}(2)=2$ and $v_{\mathfrak p}(\sqrt D)=1$, so
  $v_{\mathfrak p}(\mathfrak d)=3$ while $v_{\mathfrak p}(6)=2$; the congruence
  $6\equiv2\mu\Delta$ modulo $\mathfrak p^{3}$ forces
  $v_{\mathfrak p}(\mu\Delta)=0$, so $\mu$ and $\Nm(\Delta)$ are odd.
\end{enumerate}
If $D\equiv1$ modulo $4$ and $D\ne5$, then $D\ge13$ and (a) has no solution
in $[1,8]$. If $D\equiv3$ modulo $4$, (a) gives $m\in\{3,6\}$ for $D=3$,
$m=2$ for $D=7$ and nothing for $D\ge11$, and (b) excludes these values. If
$D=2D'$, then (b) and (c) give $m\equiv1$ modulo $4$, which with (a) says
$m\equiv9$ modulo $4D'$, and this has no solution in $[1,8]$ unless $D'=1$.
There remain $D=5$, with $m=4$, and $D=2$, with $m\in\{1,5\}$.

For $\QQ(\sqrt5)$, $m=4$ means that $\Delta$ is $-2$ times a unit or minus the
square of a unit, since $2$ is inert; and $\Delta\equiv l_{1}^{2}$ modulo $4$,
while neither $-1$ nor twice a unit is a square modulo $4$ in $\ZZ[\varphi]$.
For $\QQ(\sqrt2)$, (c) makes $\mu$ odd, so $\mu=\pm1$ and
$\Nm(\Delta)\in\{1,5\}$; the norms of the elements of $\ZZ[\sqrt2]$ prime to
$2$ are $\pm1$ modulo $8$, so $\Nm(\Delta)=1$ and $\Delta=-(1+\sqrt2)^{2k}$,
which is $3$ or $1+2\sqrt2$ modulo $4$, and neither is a square modulo $4$ in
$\ZZ[\sqrt2]$.
\end{proof}

\begin{theorem}[No dimension-count certificate for quartic CM fields]
\label{thm:quarticobstruction}
Let $F=F_{0}(\sqrt{-q})$ be a quartic CM field with real quadratic subfield
$F_{0}$, let $X$ be a principally polarised abelian fourfold with real
multiplication by $\cO_{F_{0}}$, and give $A=X\times\widehat X$ the action of
$F$ of \cite[\S11.1]{Mar25c}, transported by the $B$-field $e^{\theta_{t}}$ for a
rational $t\in F_{0}$. Let $F_{1},F_{2}$
be perfect complexes on $X$ with $\Ext^{<0}(F_{i},F_{i})=0$, for instance
sheaves, whose Chern characters lie in spaces $S(t',q')$ of
\Cref{prop:quarticsecant}, and let $E=\Phi(F_{1}\boxtimes F_{2}^{\vee})$ and
$\kappa=\ch(E)e^{\ell/2}$. Put $v_{1}=\ch(F_{1})$ and $v_{2}=\ch(F_{2}^{\vee})$.
\begin{enumerate}[label=\textup{(\roman*)},leftmargin=2.6em,itemsep=2pt]
\item $r(\kappa)=r(\ch E)=r^{2}(v_{1})+64+r^{2}(v_{2})$, where
  $r^{2}(v_{i})\in\{12,18,20\}$ is given by \Cref{prop:quarticsecant}(ii); so
  $r(\kappa)\le104$, and $r(\kappa)\in\{100,102,104\}$ unless $F$ is
  biquadratic and some $v_{i}$ has $N_{w}=2$.
\item $\dim\Ext^{2}(E,E)>r(\kappa)$. So the injectivity of $\sigma_{E}$ cannot
  be obtained from a dimension count for any such $E$.
\item If $t=0$ and both Chern characters lie in $S(0,q)$, then $\kappa$ is
  invariant under $G_{F}$, and so is a Hodge class on every member of the
  $(F,2,1)$-family. This is the case $[F:\QQ]=4$ of
  \cite[Proposition~10.2.1(1)]{Mar25c} and of the inclusion stated before
  \cite[Lemma~10.1.3]{Mar25c}, with other conventions for $\Phi$ and for the
  sign of the twist.
\end{enumerate}
\end{theorem}

\begin{proof}
(i) By \Cref{lem:bfield}(i), multiplication by $e^{\ell/2}$ conjugates
contraction by an automorphism of $HT^{2}$, so $r(\kappa)=r(\ch E)$; and, as
in the proof of \Cref{thm:factor}(iv), $\Phi$ identifies $r(\ch E)$ with the
rank of contraction from $HT^{2}(X\times X)=\bigoplus_{a+b=2}HT^{a}(X)\otimes
HT^{b}(X)$ into $v_{1}\boxtimes v_{2}$. Write $r^{k}$ for the rank of
contraction from $HT^{k}(X)$. The block $HT^{1}\otimes HT^{1}$ maps into
$H^{\mathrm{odd}}\otimes H^{\mathrm{odd}}$ by the tensor product of the two
contractions from $HT^{1}$, of rank $r^{1}(v_{1})r^{1}(v_{2})=64$, and the other
two blocks map into $H^{\mathrm{ev}}\otimes H^{\mathrm{ev}}$. If
$(x_{1}\lrcorner v_{1})\otimes v_{2}+v_{1}\otimes(x_{2}\lrcorner v_{2})=0$, then
$x_{1}\lrcorner v_{1}=\lambda v_{1}$ for a scalar $\lambda$, and $\lambda=0$
because contraction from $HT^{2}$ raises $q-p$ by two on classes of type
$(p,q)$ while $v_{1}\ne0$ is of type $(p,p)$; so those two blocks contribute
$r^{2}(v_{1})+r^{2}(v_{2})$, and $r^{2}(v_{i})=2\rho+4N_{w}\in\{12,18,20\}$ by
\Cref{prop:quarticsecant}(ii), with $12$ only for biquadratic $F$.

(ii) By the K\"unneth formula, as in \Cref{prop:orlovfour},
$\dim\Ext^{2}(E,E)=e_{0}e_{2}'+e_{1}e_{1}'+e_{2}e_{0}'$ with
$e_{k}=\dim\Ext^{k}(F_{1},F_{1})\ge r^{k}(v_{1})$, by the commutative square
$\sigma_{F}\circ\operatorname{ev}_{F}=(\,\cdot\,)\lrcorner\ch(F)$, and likewise
for $F_{2}$. So $\dim\Ext^{2}(E,E)\ge r(\kappa)$, with equality only if each
$F_{i}$ is minimal, $\dim\Ext^{k}(F_{i},F_{i})=1,8,r^{2}(\ch F_{i})$ for
$k=0,1,2$. For such $F_{i}$, Serre duality gives $\chi(F_{i},F_{i})=r^{2}-14$,
that is $4$, $6$ or $-2$ for $(\rho,N_{w})=(1,4)$, $(2,4)$, $(2,2)$. The last
is impossible because $\chi$ is positive on $S(t',q')$. The Chern character
of a perfect complex on a complex torus is integral, since by the K\"unneth
theorem the topological $K$-theory of a real torus is generated by products of
classes pulled back from its circle factors, whose Chern characters are
integral. So the first two values are excluded by \Cref{prop:quarticother}.

(iii) Over $\RR$, $H^{1}(X,\RR)=U_{1}\oplus U_{2}$, with $U_{j}$ the eigenspace
of $F_{0}$ for $\tau_{j}$, and the $U_{j}$ are orthogonal for $\theta$. As in
the proof of \Cref{prop:flatall}, the correspondence $Z\mapsto\ch(\Phi(Z))
e^{\ell/2}$, the Poincar\'e class $\ell=\ell_{1}+\ell_{2}$, the classes $\eta$
and $\gamma$ and the exponentials $e_{\varepsilon}$ all decompose as tensor
products over $j$ of classes of even degree, so that the image of
$e_{\varepsilon}\boxtimes e_{\varepsilon'}^{\vee}$ is the tensor product over $j$
of the images of $e^{\varepsilon_{j}s_{j}\theta_{j}}\boxtimes
e^{-\varepsilon_{j}'s_{j}\theta_{j}}$ in the $j$-th factor. On that factor
$\sqrt{-q}$ acts as in the imaginary quadratic construction with $d$ replaced
by $\tau_{j}(q)$, and the identities of \Cref{prop:flatall} at $n=2$, which
are identities in the exterior algebra of a symplectic space with $d$ a
symbol, show that each image is a multiple of $e^{\pm(s_{j}/2\tau_{j}(q))\eta_{j}}$
or of a Weil class $(\gamma_{j}\mp s_{j}\ell_{j})^{2}$ of the $j$-th factor.
These are invariant under the special unitary group $SU_{j}$ of the $j$-th
factor, and $G_{F}(\RR)=SU_{1}\times SU_{2}$ acts factor by factor; so the
rational class $\kappa$ is invariant under $G_{F}(\RR)$, which is Zariski dense
in $G_{F}$, and the Hodge group of every member of the family lies in $G_{F}$.
\end{proof}

\begin{remark}[The shape of the secant characters]\label{rem:quarticsecantshape}
Let $t=0$ in \Cref{thm:quarticobstruction}(iii), and write
$\eta=\eta_{1}+\eta_{2}$, $\gamma=\gamma_{1}+\gamma_{2}$ and
$\ell=\ell_{1}+\ell_{2}$ for the decompositions over the real places used in
its proof: $\eta_{j}=\tau_{j}(q)\theta_{j}+\widehat\theta_{j}$ and
$\gamma_{j}=\tau_{j}(q)\theta_{j}-\widehat\theta_{j}$, with $\theta_{j}$ and
$\widehat\theta_{j}$ here the pull-backs to $A$ of the components of the
polarisations of $X$ and $\widehat X$. The classes $\eta_{t}$ of $A$ are the
combinations of $\eta_{1}$ and $\eta_{2}$, and the Weil classes at the two
embeddings over $\tau_{j}$ are the multiples of
$(\gamma_{j}\mp s_{j}\ell_{j})^{2}$. Let $\kappa_{(a,b)}$ be the component of
$\kappa$ of degree $a$ at $\tau_{1}$ and $b$ at $\tau_{2}$. Although $\kappa$
is flat, it does not lie in $\CC[\eta_{1},\eta_{2}]+W_{F}\otimes\CC$ once
$\kappa_{4}$ has a nonzero Weil part, unlike the classes of
\Cref{prop:flatall}. Indeed, by the proof of \Cref{thm:quarticobstruction}(iii),
$\kappa$ is a sum of products $X_{1}X_{2}$ in which each $X_{j}$ is a multiple
of $e^{\pm c_{j}\eta_{j}}$, $c_{j}=s_{j}/(2\tau_{j}(q))$, or of
$(\gamma_{j}\mp s_{j}\ell_{j})^{2}$. The components $\kappa_{(4,0)}$ and
$\kappa_{(4,8)}$ receive exactly the terms whose second factor is an
exponential, through its parts of degree $0$ and $8$, and
$c_{2}^{4}=1/(16\tau_{2}(q)^{2})$ for both signs; so
\[
  \kappa_{(4,8)}=\frac{\kappa_{(4,0)}\,\eta_{2}^{4}}{384\,\tau_{2}(q)^{2}},
\]
and symmetrically. The polynomials in $\eta_{1},\eta_{2}$ of this bidegree
are the multiples of $\eta_{1}^{2}\eta_{2}^{4}$ and $W_{F}$ has no component
there, while $\kappa_{(4,0)}$ has a nonzero Weil part as soon as $\kappa_{4}$
has one, since a nonzero rational Weil class has nonzero components at all
four embeddings. The class $\ch(E)=\kappa e^{-\ell/2}$ itself is not flat
when $\rk(E)\ne0$: for $Y$ in the Lie algebra of $G_{F}$,
$D_{Y}\ch(E)=-\tfrac12(D_{Y}\ell)\ch(E)$, whose part of degree two is
$-\tfrac12\rk(E)\,D_{Y}\ell$, and $\ell$ is not flat. So these objects, when
their rank is not zero, are candidates only for a twisted form of \Cref{thm:quarticp2prime}(v), with
hypothesis a flat twisted character $\kappa$ with $\kappa_{4}=N\omega+P_{4}$,
which needs the extension of \Cref{thm:BF} to twisted objects \cite{Per26}
and the extraction of untwisted cycles of \cite[Section~2]{Mar25b}, as after
\Cref{thm:twistedcert}. This is the form in which Markman states his
reduction: he uses the normalised class $\ch(E)\exp(-c_{1}(E)/\rk E)$, which is
$\kappa e^{-\kappa_{1}/\kappa_{0}}$ when $\rk E\ne0$, asserts the shape of a
polynomial plus a nonzero Weil class only for its component in $H^{4}$, and
takes as propagation hypothesis that this component stays algebraic
\cite[Proposition~10.2.1, Corollary~10.2.3]{Mar25c}.
\end{remark}

\begin{remark}[Negative Ext groups]\label{rem:quarticnegext}
The hypothesis $\Ext^{<0}(F_{i},F_{i})=0$ in
\Cref{thm:quarticobstruction}(ii) can be replaced by a condition on the
amplitude. Let $F_{1},F_{2}$ be any perfect complexes with characters in
spaces $S(t',q')$, let $e_{k}=\dim\Ext^{k}(F_{1},F_{1})$ and
$e_{k}'=\dim\Ext^{k}(F_{2},F_{2})$, and put
$\nu_{i}=\sum_{j\ge3}(-1)^{j}\dim\Ext^{-j}(F_{i},F_{i})$. Then still
$\dim\Ext^{2}(E,E)\ge r(\kappa)$, and equality forces, for each $i$,
$\dim\Ext^{k}(F_{i},F_{i})=1,8,r^{2}(v_{i})$ for $k=0,1,2$,
$\Ext^{-1}(F_{i},F_{i})=\Ext^{-2}(F_{i},F_{i})=0$ and
$\nu_{i}=\bigl(\chi(v_{i},v_{i})-r^{2}(v_{i})+14\bigr)/2$. This gives
$\nu_{i}\ge2$ for $(\rho,N_{w})=(1,4)$ and $(2,2)$, and $\nu_{i}\ge1$ for
$(2,4)$ unless $F_{0}=\QQ(\sqrt3)$ and $\chi(v_{i},v_{i})=4$, where
$\nu_{i}=-1$, a case that needs integral classes of that kind. So
equality needs, for both $i$, some $\Ext^{-j}(F_{i},F_{i})\ne0$ with $j\ge4$
even, or $j\ge3$ odd in the exceptional case, and it fails as soon as one
$F_{i}$ has its cohomology sheaves in an interval of length at most $3$, at
most $2$ in the exceptional case. These conditions are necessary, not
sufficient. The strict inequality $\dim\Ext^{2}(E',E')>r(\ch E')$ passes to
the images $E'$ of such an $E$ under derived equivalences, push-forward and
pull-back along isogenies, and twists by line bundles or $B$-fields.

Indeed, $\dim\Ext^{2}(E,E)=\sum_{a\in\ZZ}e_{a}e_{2-a}'$ by the K\"unneth
formula, and $e_{a}=e_{4-a}$ by Serre duality. The terms $a=0,1,2$ give at
least $r(\kappa)$, with equality only if $e_{k}$ and $e_{k}'$ are the ranks
$r^{k}$ for $k\le2$, as in the proof of \Cref{thm:quarticobstruction}(ii); the
terms $a=3,4,-1,-2$ are $e_{1}e_{-1}'$, $e_{0}e_{-2}'$, $e_{-1}e_{1}'$ and
$e_{-2}e_{0}'$, which must vanish. Then
$\chi(F_{i},F_{i})=2e_{0}-2e_{1}+e_{2}+2\sum_{j\ge1}(-1)^{j}e_{-j}$ gives the
value of $\nu_{i}$. The number $\chi(v,v)$ is a positive even integer:
positive by \Cref{prop:quarticsecant}(iii), and even because $\int_{X}x^{2}$ is
even for $x\in H^{4}(X,\ZZ)=\bigwedge^{4}H^{1}(X,\ZZ)$. For $(1,4)$,
$\chi\in4\ZZ$ by the proof of \Cref{prop:quarticother}(i), and $\chi\ne4$; for
$(2,2)$, $r^{2}=12$. For $(2,4)$, $r^{2}=20$, and in the proof of
\Cref{prop:quarticother}(ii) $\chi^{2}-4m>0$ with $m\ge1$, so $\chi\ge4$, and
$\chi\ne6$. If $\chi=4$, then $m\in\{1,2,3\}$, and conditions (a) to (c) of
that proof, with $36$ replaced by $16$, say that $m\equiv4$ modulo the odd
part of $D$, so that $D\in\{2,3,6\}$; that $m=1$; and, for $D\equiv2$ modulo
$4$, that $m$ is even, because $v_{\mathfrak p}(4)=4\ge3$ here gives
$v_{\mathfrak p}(\mu\Delta)\ge1$. Only $D=3$ remains. A complex with cohomology
sheaves in an interval of length $l$ has $\Ext^{-j}(F,F)=0$ for $j>l$.
Equivalences preserve Ext algebras and, through the induced isomorphisms of
$HT^{*}$ and $H^{*}$, ranks of contraction (\cite{CRV12}, as in the proof of
\Cref{thm:factor}(iv)); for an isogeny $g$,
$g^{*}g_{*}E=\bigoplus_{x\in\ker g}t_{x}^{*}E$, so $\Ext^{2}(g_{*}E,g_{*}E)$
contains $\Ext^{2}(E,E)$ while $r(\ch g_{*}E)=r(\ch E)$, and likewise for
$g^{*}$; twists are \Cref{lem:bfield}(i).
\end{remark}

What \Cref{thm:quarticobstruction} and \Cref{rem:quarticnegext} leave for
$\mathbf{W}(F,2,\delta)$ is one of three things: an Orlov product both of
whose factors have nonzero Ext groups in some even degree at most $-4$ (or odd
degree at most $-3$ when $F_{0}=\QQ(\sqrt3)$); a kernel on $X\times X$ that is
not an external product; or a proof of the injectivity of $\sigma_{E}$ for
some Orlov product that does not go through a dimension count, which is what
Markman's question \cite[Question~11.2.2]{Mar25c} asks for. The next
proposition shows that a kernel whose transform has a pure spinor as its
character, for instance a line bundle on the graph of a homomorphism, gives no
Weil class.

\begin{proposition}[Pure spinors give no Weil class]
\label{prop:quarticexclusions}
Keep the notation of \Cref{rem:quarticsecantshape} on $A=X\times\widehat X$
with $t=0$, let $V_{j}\subseteq H^{1}(A,\RR)$ be the part at $\tau_{j}$, and
call a class flat if it is invariant under $G_{F}(\RR)$. Let
$v=c\,e^{\beta}\Omega_{N}$ be a real pure spinor, with $c\in\RR^{\times}$,
$\beta$ a real $2$-form, $N\subseteq H^{1}(A,\RR)$ and $\Omega_{N}$ a
generator of $\det N$. If $v$ is flat, then $N$ is one of $0$, $V_{1}$,
$V_{2}$, $H^{1}(A,\RR)$, and $v\in\RR[\eta_{1},\eta_{2}]$. This applies to
the characters of line bundles on translates of abelian subvarieties, of
points and of simple semi-homogeneous bundles, and of their images under
derived equivalences and isogenies, twisted by any real $B$-field; for
instance to $\Phi$ of a line bundle on the graph of a homomorphism $X\to
X\times X$. The component in $W_{F}\otimes\RR$ of such a class vanishes in
every degree: no such object has a character, twisted as stated, of the form
required by \Cref{thm:quarticp2prime}(v) or by its twisted form, whatever
the quartic field.
\end{proposition}

\begin{proof}
Over $\RR$, $G_{F}(\RR)=SU_{1}\times SU_{2}$ with $SU_{j}$ acting on $V_{j}$
and trivially on the other part, and $V_{j}\otimes\CC$ is the sum of the
standard representation of $SU_{j}\otimes\CC\cong\SL_{4}$ and its dual, which
are not isomorphic (proof of \Cref{prop:quarticquat}(i)). So $V_{1}$ and
$V_{2}$ are irreducible and not isomorphic, and $0$, $V_{1}$, $V_{2}$,
$H^{1}(A,\RR)$ are the only stable subspaces. The flat classes are
$\mathrm{Inv}_{1}\otimes\mathrm{Inv}_{2}$, where by the same proof
$\mathrm{Inv}_{j}\subseteq\bigwedge^{*}V_{j}$ is spanned by the powers of
$\eta_{j}$ and the two Weil classes $(\gamma_{j}\mp s_{j}\ell_{j})^{2}$; in
degree two it is $\RR\eta_{j}$, and $\eta_{j}^{4}$ spans $\bigwedge^{8}V_{j}$.

If $v$ is real and pure over $\CC$, then $N=\{\xi:\xi\wedge v=0\}$ is
defined over $\RR$, and so are $c$ and the class of $\beta$ modulo
$N\wedge H^{1}$. For $g\in G_{F}(\RR)$, $gv=v$ gives $gN=N$, and
$g\Omega_{N}=\det(g|_{N})\Omega_{N}=\Omega_{N}$ because the connected
semisimple group $G_{F}(\RR)$ has no nontrivial character. Hence
$(e^{g\beta-\beta}-1)\Omega_{N}=0$, whose part of lowest degree gives
$g\beta-\beta\in N\wedge H^{1}$. So $N$ is one of the four stable subspaces,
and the class of $\beta$ in $\bigwedge^{2}(H^{1}/N)$ is invariant; by complete
reducibility it is the image of an invariant $2$-form
$a_{1}\eta_{1}+a_{2}\eta_{2}$, and since $e^{n\wedge x}\Omega_{N}=\Omega_{N}$
for $n\in N$, $v=c\,e^{a_{1}\eta_{1}+a_{2}\eta_{2}}\Omega_{N}$, with
$\Omega_{N}$ a multiple of $1$, $\eta_{1}^{4}$, $\eta_{2}^{4}$ or
$\eta_{1}^{4}\eta_{2}^{4}$. The characters listed are pure spinors, and
derived equivalences, isogenies and twists preserve purity; this is the
classical description recalled in \Cref{setup:lagquadric} and after
\Cref{cor:naturalsums}.
\end{proof}

\subsubsection*{Markman's candidate for a quartic field}

Markman gives one explicit pair of sheaves for a biquadratic field
\cite[Example~11.2.7, Lemma~11.2.8]{Mar25c} and asks whether the second can
be chosen so that its semiregularity map is injective on the image of the
evaluation map \cite[Question~11.2.2]{Mar25c}. For one object on the fourfold
this is the weakened criterion of \Cref{setup:criterion}, and the local
method of \Cref{thm:candidatefails} decides it for his pair: both sheaves
fail it, at points of the curves that the construction glues in. Two local
computations are needed, one for a sheaf that is locally free on a smooth
support and one for the ideal of a curve in a smooth divisor. As in
\Cref{lem:localproducts}, for a germ of vector field $\theta$ and a perfect
complex $F$ of modules over a local ring we write
$a_{\theta}\in\Ext^{1}(F,F)$ for the class of the jet extension along
$\theta$.

\begin{lemma}[Free germs on smooth supports]\label{lem:freegerm}
Let $\cO$ be the local ring of a smooth complex variety $Y$ at a point $p$,
or its completion, let $C\ni p$ be a smooth subvariety of codimension
$c\ge1$, put $N=T_{p}Y/T_{p}C$, and let $F$ be a perfect complex of
$\cO$-modules quasi-isomorphic to a shift of $\cO_{C,p}^{\oplus r}$, $r\ge1$.
Then $\Ext^{2}_{\cO}(F,F)$ is a free $\cO_{C,p}$-module isomorphic to
$M_{r}(\cO_{C,p})\otimes\bigwedge^{2}N$, and modulo the maximal ideal of
$\cO_{C,p}$ the product $a_{\theta}a_{\theta'}$ is, up to a sign that does
not depend on $\theta,\theta'$, the element $\id\otimes(\bar\theta\wedge\bar\theta')$,
where $\bar\theta,\bar\theta'$ are the images of $\theta(p),\theta'(p)$ in
$N$. In particular $a_{\theta}a_{\theta'}\ne0$ whenever $\theta(p)$,
$\theta'(p)$ and $T_{p}C$ span a space of dimension $\dim C+2$.
\end{lemma}

\begin{proof}
Choose formal coordinates $x_{1},\dots,x_{m}$ at $p$ with
$C=\{x_{1}=\dots=x_{c}=0\}$. Ext groups of finitely generated modules and jet
extensions commute with the flat base change from
$R=\CC[x_{1},\dots,x_{m}]$ to $\cO$, and a shift changes neither the Ext
algebra nor the jet classes, so we may take $F=(R/(x_{1},\dots,x_{c}))^{\oplus r}$
over $R$. Write $R=A_{1}\otimes\dots\otimes A_{c}\otimes B$ with
$A_{i}=\CC[x_{i}]$ and $B=\CC[x_{c+1},\dots,x_{m}]$; then
$F=k_{1}\otimes\dots\otimes k_{c}\otimes B^{\oplus r}$, with $k_{i}$ the
residue field of $A_{i}$ at the origin. On the resolution
$A_{i}\xrightarrow{x_{i}}A_{i}$ of $k_{i}$ the jet chain map of
$\partial_{i}$ is $-\partial_{i}(x_{i})=-1$, as in the proof of
\Cref{lem:localproducts}, so its class $\gamma_{i}$ spans
$\Ext^{1}_{A_{i}}(k_{i},k_{i})=k_{i}$, and $\Ext^{j}_{A_{i}}(k_{i},k_{i})=0$ for
$j\ge2$. The module $B^{\oplus r}$ is free, so its Ext algebra is
$M_{r}(B)$ in degree zero and its jet classes vanish. By
\eqref{eq:kunnethext}, applied factor by factor, $\Ext^{*}_{R}(F,F)$ is the
graded tensor product of the exterior algebra on $\gamma_{1},\dots,\gamma_{c}$
with $M_{r}(B)$, and $a_{\partial_{i}}=\gamma_{i}\otimes\id$ for $i\le c$,
while $a_{\partial_{i}}=0$ for $i>c$. Hence
$\Ext^{2}_{R}(F,F)=\bigwedge^{2}\langle\gamma_{1},\dots,\gamma_{c}\rangle\otimes M_{r}(B)$,
and $B=\cO_{C}$ in these coordinates. For $\theta=\sum_{i}f^{i}\partial_{i}$
and $\theta'=\sum_{i}f'^{i}\partial_{i}$ the linearity of \Cref{lem:edge}
gives
\[
  a_{\theta}a_{\theta'}=\sum_{i,j\le c}f^{i}f'^{j}\,\gamma_{i}\gamma_{j}\otimes\id
  =\sum_{i<j\le c}\bigl(f^{i}f'^{j}-f^{j}f'^{i}\bigr)\,\gamma_{i}\gamma_{j}\otimes\id,
\]
because the $\gamma_{i}$ anticommute and square to zero. The functions
$x_{1},\dots,x_{c}$ act by zero, so modulo the maximal ideal of $\cO_{C}$ the
coefficients are those of $\theta(p)\wedge\theta'(p)$ on
$\partial_{i}\wedge\partial_{j}$, $i<j\le c$, which are the coordinates of
$\bar\theta\wedge\bar\theta'$ in $\bigwedge^{2}N$. A section of a free module
with nonzero reduction is nonzero.
\end{proof}

\begin{lemma}[A curve in a divisor]\label{lem:divisorgerm}
Let $\cO$, $Y$ and $p$ be as in \Cref{lem:freegerm}, let $D\ni p$ be a
smooth divisor, let $C\subset D$ be a smooth subvariety of codimension two in
$D$ through $p$, and let $F$ be quasi-isomorphic to a shift of the ideal
$I_{C\subset D}$ of $C$ in $\cO_{D}$, regarded as an $\cO$-module. Put
$\nu=T_{p}Y/T_{p}D$ and $N'=T_{p}D/T_{p}C$. Then
$\Ext^{2}_{\cO}(F,F)\cong\cO_{C,p}\otimes\nu\otimes N'$, and modulo the maximal
ideal of $\cO_{C,p}$ the product $a_{\theta}a_{\theta'}$ is, up to a fixed
sign, the image of $\theta(p)\wedge\theta'(p)$ under
\[
  \textstyle\bigwedge^{2}T_{p}Y\longrightarrow\bigwedge^{2}(T_{p}Y/T_{p}C)
  \longrightarrow\bigwedge^{2}(T_{p}Y/T_{p}C)\big/\bigwedge^{2}N'\cong\nu\otimes N'.
\]
In particular $a_{\theta}a_{\theta'}\ne0$ whenever $\theta(p)\wedge\theta'(p)$
does not lie in $\bigwedge^{2}T_{p}D+T_{p}C\wedge T_{p}Y$.
\end{lemma}

\begin{proof}
With formal coordinates in which $D=\{x_{1}=0\}$ and $C=\{x_{1}=x_{2}=x_{3}=0\}$,
and the same reduction to the polynomial ring, write
$R=\CC[x_{1}]\otimes\CC[x_{2},x_{3}]\otimes B$ with $B=\CC[x_{4},\dots,x_{m}]$.
Then $I_{C\subset D}=k\otimes\mathfrak m\otimes B$, where $k$ is the residue
field of $\CC[x_{1}]$ and $\mathfrak m=(x_{2},x_{3})$. As in the proof of
\Cref{lem:localproducts}, $\mathfrak m$ has projective dimension one,
$\Ext^{1}(\mathfrak m,\mathfrak m)=\mathfrak m/\mathfrak m^{2}$ has the basis
$\alpha_{2}=a_{\partial_{2}}$, $\alpha_{3}=a_{\partial_{3}}$ and is killed by
$\mathfrak m$, and $\Ext^{2}(\mathfrak m,\mathfrak m)=0$; the factor $k$
contributes $\gamma_{1}$ in degree one, as in \Cref{lem:freegerm}, and $B$
contributes nothing in positive degree. By \eqref{eq:kunnethext}
\[
  \Ext^{2}_{R}(I_{C\subset D},I_{C\subset D})
  =\Ext^{1}(k,k)\otimes\Ext^{1}(\mathfrak m,\mathfrak m)\otimes B
  =\langle\gamma_{1}\rangle\otimes\langle\alpha_{2},\alpha_{3}\rangle\otimes B,
\]
the other summands vanishing by the projective dimensions, and
$a_{\partial_{1}}a_{\partial_{j}}=-a_{\partial_{j}}a_{\partial_{1}}=\pm\gamma_{1}\otimes\alpha_{j}$
for $j=2,3$, while $a_{\partial_{2}}a_{\partial_{3}}=1\otimes\alpha_{2}\alpha_{3}=0$,
$a_{\partial_{1}}^{2}=0$ and $a_{\partial_{i}}=0$ for $i\ge4$. As before,
$a_{\theta}a_{\theta'}$ reduces to the component of $\theta(p)\wedge\theta'(p)$
on $\partial_{1}\wedge\langle\partial_{2},\partial_{3}\rangle$. In these
coordinates $\partial_{1}$ spans $\nu$, the images of $\partial_{2},\partial_{3}$
span $N'$, and the component is the image in
$\bigwedge^{2}(T_{p}Y/T_{p}C)/\bigwedge^{2}N'$. Its kernel on
$\bigwedge^{2}T_{p}Y$ is spanned by $\partial_{2}\wedge\partial_{3}$ and the
$\partial_{i}\wedge\partial_{j}$ with $j\ge4$, which is
$\bigwedge^{2}T_{p}D+T_{p}C\wedge T_{p}Y$.
\end{proof}

Now keep the notation of \Cref{prop:quarticsecant}, with $X$ an abelian
fourfold with real multiplication by $F_{0}$ and $\theta=\theta_{1}+\theta_{2}$
a polarisation in $\bigwedge^{2}_{F_{0}}H^{1}(X,\QQ)$. Over $\CC$,
$H^{1}(X)=U_{1}\oplus U_{2}$ with $U_{j}$ the eigenspace of $F_{0}$ at
$\tau_{j}$, and the tangent space splits accordingly as $T=T_{1}\oplus T_{2}$,
with $T_{j}$ of dimension two, dual to $U_{j}\cap H^{1,0}$; translation
identifies every $T_{p}X$ with $T$. Let $\pi_{j}$ span $\bigwedge^{2}T_{j}$
and put
\begin{equation}\label{eq:quarticcompensated}
  x_{j}=\bigl(\tfrac12\tau_{j}(q)\,\pi_{j}\lrcorner\,\theta_{j}^{2},\;0,\;\pi_{j}\bigr)
  \in HT^{2}(X),\qquad j=1,2.
\end{equation}
The six classes $(0,v,0)$ with $v\in\Hom(H^{1,0},H^{0,1})$ preserving the
decomposition by places and $v\lrcorner\,\theta_{1}=v\lrcorner\,\theta_{2}=0$
are the first order deformations of $X$ with its real multiplication and
polarisation; we call them the polarised $F_{0}$-linear deformations.

\begin{theorem}[The local obstruction for a quartic secant character]
\label{thm:quarticlocal}
In this notation, let $q\in F_{0}$ be totally positive and $t\in F_{0}$.
\begin{enumerate}[label=\textup{(\roman*)},leftmargin=2.6em,itemsep=2pt]
\item $x_{j}\lrcorner\,w=0$ for every $w\in S(0,q)$, and
  $(e^{-\theta_{t}}x_{j})\lrcorner\,w=0$ for every $w\in S(t,q)$, $j=1,2$. If
  $w\in S(t,q)$ has $(\rho,N_{w})=(2,4)$, then $\Ann_{HT^{2}}(w)$ is spanned
  by $e^{-\theta_{t}}x_{1}$, $e^{-\theta_{t}}x_{2}$ and the six polarised
  $F_{0}$-linear deformations, and its bivector components span exactly
  $\langle\pi_{1},\pi_{2}\rangle$.
\item Let $F$ be a perfect complex on $X$ with $\ch(F)\in S(t,q)$, and let
  $\theta,\theta'$ be translation invariant vector fields with
  $\theta\wedge\theta'=\pi_{j}$. If $a_{\theta}a_{\theta'}\ne0$ in
  $\Ext^{2}(F_{p},F_{p})$ at some point $p$, then
  $\operatorname{ev}_{F}(e^{-\theta_{t}}x_{j})\ne0$ and $F$ does not satisfy the
  weakened criterion. Neither do $\Phi(F\boxtimes G)$ and $\Phi(G\boxtimes F)$
  on $X\times\widehat X$, for any nonzero perfect complex $G$ on $X$, with
  $\Phi$ as in \Cref{thm:factor}(iv).
\item The hypothesis of \textup{(ii)} holds at $p$, for $j=1$ or for $j=2$,
  whenever $F_{p}$ is a shift of a free module of positive rank over the
  local ring of a smooth curve through $p$, since $T_{p}C$ lies in at most one
  of $T_{1},T_{2}$; it holds for both $j$ when $F_{p}$ is a shift of a free
  module over the residue field. It holds for $j$ when $F_{p}$ is a shift of
  the ideal of a smooth curve $C$ in a smooth divisor $D$ with
  $T_{p}C\not\subseteq T_{j}$ and $T_{j}\not\subseteq T_{p}D$.
\end{enumerate}
\end{theorem}

\begin{proof}
(i) Let $t=0$. By the proof of \Cref{lem:bfield}(i), for a bivector $\pi$
of type $(2,0)$ in the tangent directions and a class $B$ of type $(1,1)$,
$\pi\lrcorner\,e^{B}=\tfrac12(\pi\lrcorner\,B^{2})\,e^{B}$. The interior
products by vectors of $T_{j}$ are odd derivations that vanish on
$\bigwedge^{*}U_{j'}$ for $j'\ne j$, so $\pi_{j}\lrcorner(ab)=(\pi_{j}\lrcorner\,a)\,b$
for $a$ in the subalgebra generated by $U_{j}$ and $b$ an even class
generated by $U_{j'}$. Applied to
$e_{\varepsilon}=e^{\varepsilon_{1}s_{1}\theta_{1}}e^{\varepsilon_{2}s_{2}\theta_{2}}$,
with $s_{j}^{2}=-\tau_{j}(q)$, this gives
$\pi_{j}\lrcorner\,e_{\varepsilon}=-\tfrac12\tau_{j}(q)(\pi_{j}\lrcorner\,\theta_{j}^{2})\,e_{\varepsilon}$
for all four signs, and so $x_{j}\lrcorner\,e_{\varepsilon}=0$. A class
$(0,v,0)$ with $v$ preserving the places and killing $\theta_{1}$ and
$\theta_{2}$ acts as a derivation that kills every polynomial in
$\theta_{1},\theta_{2}$, hence all of $S(0,q)$. For general $t$,
\Cref{lem:bfield}(i) gives $y\lrcorner(we^{\theta_{t}})=((e^{\theta_{t}}y)\lrcorner\,w)e^{\theta_{t}}$,
so $e^{-\theta_{t}}$ carries $\Ann_{HT^{2}}(w)$ onto
$\Ann_{HT^{2}}(we^{\theta_{t}})$; it fixes the six classes $(0,v,0)$,
because $v\lrcorner\,\theta_{t}=0$, and by \eqref{eq:bfield} it does not
change bivector components. The eight classes are independent: the six have
no bivector component and are independent, and the other two have the
independent components $\pi_{1},\pi_{2}$. When $(\rho,N_{w})=(2,4)$ the rank
of contraction from $HT^{2}(X)$, of dimension $28$, is $2\rho+4N_{w}=20$ by
\Cref{prop:quarticsecant}(ii), so the annihilator has dimension $8$ and the
eight classes span it.

(ii) By (i), $y=e^{-\theta_{t}}x_{j}=(z,v,\pi_{j})$ lies in
$\Ann_{HT^{2}}(\ch F)$, and by \eqref{eq:evgeneral}
$\operatorname{ev}_{F}(y)=z\cdot\id_{F}+v\lrcorner\,\mathrm{At}(F)
+\tfrac12\pi_{j}\lrcorner\,\mathrm{At}(F)^{2}$. The last term is
$A_{\theta}A_{\theta'}$, as in the proof of \Cref{thm:candidatefails}, and its
germ at $p$ is $a_{\theta}a_{\theta'}\ne0$ by \Cref{lem:edge}. The germs of the
first two terms vanish: $z\cdot\id_{F}$ by \Cref{lem:edge}, and
$v\lrcorner\,\mathrm{At}(F)$ because it is a composite
$F\to F\otimes\Omega^{1}_{X}[1]\to F\otimes\Omega^{1}_{X}\otimes T_{X}[2]\to F[2]$
in which the middle arrow is $\id\otimes v$ with
$v\in\Ext^{1}(\cO_{X},T_{X})=H^{1}(T_{X})$, and $v$ restricts to zero on
every affine open set. So $\operatorname{ev}_{F}(y)\ne0$ with
$y\lrcorner\,\ch(F)=0$, and $F$ fails the weakened criterion. For
$G\boxtimes F$ and $F\boxtimes G$ the class $y\otimes1$, resp. $1\otimes y$,
preserves the Chern character, and its evaluation
$\pm\operatorname{ev}_{F}(y)\otimes\id_{G}$ is nonzero in its K\"unneth
summand, as in the proof of \Cref{thm:factor}(iv); the transport under
$\Phi$ is the one given there.

(iii) Since $T_{j}$ is a plane, the image of $\pi_{j}$ in
$\bigwedge^{2}(T/T_{p}C)$ is nonzero exactly when $T_{j}\cap T_{p}C=0$. For a
curve this says $T_{p}C\not\subseteq T_{j}$, which fails for at most one $j$
because $T_{1}\cap T_{2}=0$; for a point it always holds. \Cref{lem:freegerm}
then gives $a_{\theta}a_{\theta'}\ne0$. In the last case the image of
$\pi_{j}$ in $\bigwedge^{2}(T/T_{p}C)$ is the nonzero decomposable vector of
the plane $(T_{j}+T_{p}C)/T_{p}C$, and it lies in $\bigwedge^{2}N'$ only if
that plane is $N'$, that is, only if $T_{j}+T_{p}C=T_{p}D$, which would put
$T_{j}$ inside $T_{p}D$. So \Cref{lem:divisorgerm} applies.
\end{proof}

At a point where $F$ is locally free both sides of the local identity vanish,
as at a smooth point of a surface in \Cref{prop:smoothlocal}; the test sees
only where the sheaf drops dimension or fails to be locally free on its
support. Markman's construction produces exactly such points.

\begin{proposition}[The classes of Markman's candidate]
\label{prop:markmanquarticclasses}
Let $X$ be the Jacobian of a curve of genus four with its principal
polarisation $\Theta$, with real multiplication by a real quadratic field
$F_{0}$ and an automorphism $g$ with $g^{*}=\hat\eta(f)$ on $H^{1}(X,\QQ)$,
where $f\in F_{0}$ has norm one and $f^{2}\ne1$; let $q>0$ be rational, so
that $F=F_{0}(\sqrt{-q})$ is biquadratic, and write $f_{j}=\tau_{j}(f)$. Let
$\alpha_{0}=\Theta-\frac q6\Theta^{3}$, the character of the sheaf $E_{0}=F_{q}$
of \cite[Example~8.2.4]{Mar25a}, and
$\beta'=g^{*}\Theta-\frac q6(g^{-1})^{*}\Theta^{3}$, a multiple of which is the
character of the sheaf $E'$ of \cite[Example~11.2.7]{Mar25c}. With
$A_{j}=1-\frac q2\theta_{j}^{2}$ and $\theta=\Theta$:
\begin{enumerate}[label=\textup{(\roman*)},leftmargin=2.6em,itemsep=2pt]
\item $\alpha_{0}=\theta_{1}A_{2}+\theta_{2}A_{1}$ and
  $\beta'=f_{1}^{2}\theta_{1}A_{2}+f_{2}^{2}\theta_{2}A_{1}$, both in $S(0,q)$.
  With $s=\sqrt{-q}$, the coefficients of $\beta'$ on the four
  $e_{\varepsilon}$ are $w_{++}=-w_{--}=(f_{1}^{2}+f_{2}^{2})/4s$ and
  $w_{+-}=-w_{-+}=(f_{1}^{2}-f_{2}^{2})/4s$, so $\beta'$ has
  $(\rho,N_{w})=(2,4)$, while $\alpha_{0}$ lies on the secant plane of
  $\QQ(\sqrt{-q})$ and has $(\rho,N_{w})=(2,2)$.
\item The ranks of contraction from $HT^{2}(X)$ are $12$ for $\alpha_{0}$ and
  $20$ for $\beta'$. The annihilator of $\beta'$ is that of
  \Cref{thm:quarticlocal}(i); that of $\alpha_{0}$ contains, besides it, the
  classes $(\frac q2\,\pi\lrcorner\,\Theta^{2},0,\pi)$ for every bivector
  $\pi$, by \Cref{thm:factor}(ii) with $t=\Theta$ and $d=q$.
\item $\int_{X}\alpha_{0}\beta'=-4q\Tr_{F_{0}/\QQ}(f^{2})\ne0$, which is
  Markman's formula for the rank of $\Phi(E_{0}\boxtimes E')$
  \cite[Lemma~11.2.8]{Mar25c}; $\chi(\beta',\beta')=4q\Tr_{F_{0}/\QQ}(f^{4})$ and
  $\chi(\alpha_{0},\alpha_{0})=8q$; and the rank of contraction from
  $HT^{2}(X\times X)$ into $\alpha_{0}\boxtimes\beta'$ is $12+64+20=96$, as
  \Cref{thm:quarticobstruction}(i) predicts.
\end{enumerate}
\end{proposition}

\begin{proof}
(i) The automorphism $g$ acts on $U_{j}$ by $f_{j}$, hence on $\theta_{j}$ by
$f_{j}^{2}$, and $f_{1}f_{2}=1$. Since $\theta_{j}^{3}=0$,
$\Theta^{3}=3\theta_{1}^{2}\theta_{2}+3\theta_{1}\theta_{2}^{2}$ and
$(g^{-1})^{*}\Theta^{3}=3f_{1}^{-4}f_{2}^{-2}\theta_{1}^{2}\theta_{2}
+3f_{1}^{-2}f_{2}^{-4}\theta_{1}\theta_{2}^{2}
=3f_{2}^{2}\theta_{1}^{2}\theta_{2}+3f_{1}^{2}\theta_{1}\theta_{2}^{2}$, which
gives the two expressions. With $t=0$ and
$e_{\varepsilon}=A_{1}A_{2}-q\varepsilon_{1}\varepsilon_{2}\theta_{1}\theta_{2}
+s(\varepsilon_{1}\theta_{1}A_{2}+\varepsilon_{2}\theta_{2}A_{1})$, the class
$\sum_{\varepsilon}w_{\varepsilon}e_{\varepsilon}$ equals $\beta'$ exactly when
$\sum w_{\varepsilon}=\sum\varepsilon_{1}\varepsilon_{2}w_{\varepsilon}=0$,
$s\sum\varepsilon_{1}w_{\varepsilon}=f_{1}^{2}$ and
$s\sum\varepsilon_{2}w_{\varepsilon}=f_{2}^{2}$, whose solution is the one
stated. All four coefficients are nonzero because $f_{1}^{2}=f_{2}^{2}$ with
$f_{1}f_{2}=1$ would force $f^{2}=1$. For $\alpha_{0}$ the same computation
with $f_{1}=f_{2}=1$ gives $w_{+-}=w_{-+}=0$. The matrix $V$ of
\Cref{prop:quarticsecant} has rank two in both cases.

(ii) The ranks are $2\rho+4N_{w}$, by \Cref{prop:quarticsecant}(ii);
$\alpha_{0}=v_{\Theta}$ in the notation of \eqref{eq:uv}, with $d=q$.

(iii) By $\int_{X}\theta_{1}^{2}\theta_{2}^{2}=4$,
$\int\alpha_{0}\beta'=-2q(f_{1}^{2}+f_{2}^{2}+f_{1}^{-2}+f_{2}^{-2})=-4q\Tr(f^{2})$,
which is negative; the Euler pairings are
\Cref{prop:quarticsecant}(iii) in the basis given there. The last number is
\Cref{thm:quarticobstruction}(i) with $r^{2}(\alpha_{0})=12$ and
$r^{2}(\beta')=20$.
\end{proof}

\begin{corollary}[Markman's candidate does not satisfy the weakened
criterion]\label{cor:markmanquarticfails}
In \Cref{prop:markmanquarticclasses}, let $E'$ be the sheaf of
\cite[Example~11.2.7]{Mar25c}, obtained by gluing $N$ general translates of
$g^{*}F'$ to a line bundle on a curve $C'$ at the points where $C'$ meets
them, and assume that $C'$ is reduced.
\begin{enumerate}[label=\textup{(\roman*)},leftmargin=2.6em,itemsep=2pt]
\item $E'$ does not satisfy the weakened criterion, so it does not answer
  \cite[Question~11.2.2]{Mar25c}.
\item $E_{0}$ does not satisfy the weakened criterion.
\item For every nonzero perfect complex $G$ on $X$, neither
  $\Phi(E_{0}\boxtimes G)$ nor $\Phi(G\boxtimes E')$ nor
  $\Phi(G\boxtimes E'^{\vee})$ satisfies it; in particular Markman's object
  $\Phi(E_{0}\boxtimes E')$ does not.
\end{enumerate}
\end{corollary}

\begin{proof}
(i) Let $p$ be a smooth point of $C'$ that lies on none of the translates.
Near $p$ the sheaf $E'$ is the line bundle on $C'$, so $E'_{p}\cong\cO_{C',p}$,
and $\ch(E')$ is a multiple of $\beta'\in S(0,q)$.
\Cref{thm:quarticlocal}(ii) and (iii) apply with the $j$ for which
$T_{p}C'\not\subseteq T_{j}$.

(ii) In \cite[Example~8.2.4]{Mar25a} the sheaf $E_{0}=F_{q}$ is the
push-forward, from the theta divisor $\Theta\subset X$, of the ideal twisted
by $\cO(\Theta)$ of the union of a translate $W$ of the surface $W_{2}$ and
$q$ translates $C_{1},\dots,C_{q}$ of the curve $W_{1}$, each meeting $W$ in
one point, and its character is $\alpha_{0}$. Let $p$ be a point of $C_{1}$ outside $W$, the other $C_{i}$
and the finite singular locus of $\Theta$. Near $p$, $E_{0}$ is the ideal of
the smooth curve $C_{1}$ in the smooth divisor $\Theta$. Choose translation
invariant fields with $\theta(p)\notin T_{p}\Theta$ and
$\theta'(p)\in T_{p}\Theta\setminus T_{p}C_{1}$. Then
$\theta(p)\wedge\theta'(p)$ has nonzero image in $\nu\otimes N'$, and
\Cref{lem:divisorgerm} gives $a_{\theta}a_{\theta'}\ne0$. By
\Cref{prop:markmanquarticclasses}(ii) the class
$x=(\frac q2\,\pi\lrcorner\,\Theta^{2},0,\pi)$ with $\pi=\theta\wedge\theta'$
preserves $\alpha_{0}$, and the argument of \Cref{thm:quarticlocal}(ii) gives
$\operatorname{ev}_{E_{0}}(x)\ne0$.

(iii) The first two cases are the last assertion of
\Cref{thm:quarticlocal}(ii), with the class of (ii) for $E_{0}$. For
$E'^{\vee}$, the germ at $p$ is the dual of $\cO_{C'}$, a shift of a line
bundle on $C'$ because $C'$ is a local complete intersection of codimension
three near $p$, and $\ch(E'^{\vee})=-\ch(E')$ lies in $S(0,q)$.
\end{proof}

\Cref{fig:quarticlocal} draws the support of Markman's sheaf near a point
of the glued curve, with the two eigenplanes of the real multiplication.

\begin{remark}[What is left of Markman's question]\label{rem:question1122}
The question \cite[Question~11.2.2]{Mar25c} asks for a sheaf $F_{2}$ on $X$
whose character lies in $S(0,q)$ and forms with $\alpha_{0}$ a pair as in
\cite[Lemma~11.2.1]{Mar25c}, and whose semiregularity map is injective on the
image of $HT^{2}(X)$ in $\Ext^{2}(F_{2},F_{2})$. \Cref{cor:markmanquarticfails}
does not settle it. It decides it for the sheaf $E'$ of the example, and for
every choice of $N$, of the translates and of the gluing, because the
obstruction sits at any smooth point of the glued curve. By \Cref{thm:quarticlocal}, a sheaf that answers the question, and
likewise a first factor with character on the secant plane that meets the
criterion, has no point at which it is locally free on a smooth curve or on
a point, no point at which it is the ideal of a curve in a smooth divisor in
the position of \Cref{thm:quarticlocal}(iii), and at a point where it is
locally free on a smooth surface $S$ the tangent plane $T_{p}S$ must meet
both $T_{1}$ and $T_{2}$, by \Cref{lem:freegerm} with $c=2$. Corrections of
the character by sheaves on curves, the device of \cite[Example~11.2.7]{Mar25c},
are therefore excluded. Sheaves that are locally free on smooth divisors, and
the global part of the identity
$\operatorname{ev}_{F}(e^{-\theta_{t}}x_{j})=0$, are not reached by this
test. Even a positive answer would give an object with
$\dim\Ext^{2}\ge97$, by \Cref{thm:quarticobstruction}(ii) and
\Cref{prop:markmanquarticclasses}(iii), so its semiregularity could not come
from a dimension count.
\end{remark}

\begin{figure}[tbp]
\centering
  \includegraphics[width=\textwidth]{fig_quarticlocal}
\caption{\Cref{thm:quarticlocal} and \Cref{cor:markmanquarticfails}. Left: the
support of Markman's sheaf $E'$ near a point $p$ of the glued curve $C'$,
between two of the translated supports that the curve crosses at its gluing
points. The tangent space at $p$ is the sum of the eigenplanes $T_{1}$ and
$T_{2}$ of the real multiplication; the tangent line of $C'$ lies in at most
one of them, here in $T_{1}$, so $\pi_{2}$ has nonzero image in
$\bigwedge^{2}N$, $N=T_{p}X/T_{p}C'$, which is the germ of
$\operatorname{ev}_{E'}(x_{2})$ at $p$. Right: for ten germs at the origin
of $\CC^{4}$, the pairs $\{k,l\}$ of coordinate directions whose jet classes
have nonzero product in the local $\Ext^{2}$, and the projective dimension;
rows (i), (ii) and (x) are \Cref{lem:freegerm}, row (iii) is
\Cref{lem:divisorgerm}, row (vi) is \Cref{lem:localproducts}.}
\label{fig:quarticlocal}
\end{figure}

\subsection{Sextic CM fields}\label{ssec:sextic}

The dimension count of \Cref{thm:quarticobstruction} does not extend to CM
fields of degree six, for a reason of parity (\Cref{rem:sexticparity}). Let
$F=F_{0}(\sqrt{-q})$ be a sextic CM field with totally
real cubic subfield $F_{0}$, real places $\tau_{1},\tau_{2},\tau_{3}$ and $q\in
F_{0}$ totally positive, and let $X$ be a principally polarised abelian sixfold
with real multiplication by $\cO_{F_{0}}$, so that
$H^{1}(X,\CC)=U_{1}\oplus U_{2}\oplus U_{3}$ with $U_{j}$ the eigenspace of
$\tau_{j}$, of dimension four and Hodge type $(2,2)$. Let
$\theta_{j}\in H^{1,1}(X)$ be the component of the polarisation on $U_{j}$, so
that $\theta=\sum_{j}\theta_{j}$, $\theta_{j}^{3}=0$ and
$\int_{X}\theta_{1}^{2}\theta_{2}^{2}\theta_{3}^{2}=8$. For $t\in F_{0}$ and
$\varepsilon\in\{\pm1\}^{3}$ put $s_{j}=\sqrt{-1}\,\sqrt{\tau_{j}(q)}$ and
$e_{\varepsilon}=\exp\bigl(\sum_{j}(\tau_{j}(t)+\varepsilon_{j}s_{j})\theta_{j}\bigr)$,
and let $S(t,q)$ be the span of the eight $e_{\varepsilon}$. As in
\Cref{prop:quarticsecant}(i) it is a rational subspace of dimension
$8=2^{[F_{0}:\QQ]}$ of the algebra generated by the $\theta_{j}$, the Galois
group permuting the $e_{\varepsilon}$ as it permutes the CM types of $F$; for
$t=0$ it is the secant space of Markman's construction for $F$ on
$X\times\widehat{X}$ \cite[\S11.1]{Mar25c}, and
$S(t,q)=e^{\sum_{j}\tau_{j}(t)\theta_{j}}S(0,q)$. A real class
$v=\sum_{\varepsilon}w_{\varepsilon}e_{\varepsilon}$ has
$w_{-\varepsilon}=\bbar{w}_{\varepsilon}$. Write $N_{w}$ for the number of
nonzero $w_{\varepsilon}$; $A_{jj'}$ for the number of sign pairs
$(\sigma,\sigma')$ such that $w_{\varepsilon}\ne0$ for some $\varepsilon$ with
$\varepsilon_{j}=\sigma$ and $\varepsilon_{j'}=\sigma'$; $\rho_{j}$ for the rank
of the $2\times4$ matrix $(w_{\varepsilon})$ whose rows are indexed by
$\varepsilon_{j}$ and whose columns by the two other signs; and, for
$\{j,j',j''\}=\{1,2,3\}$ and $\sigma'=\pm1$, $M^{(j;j')}_{\sigma'}$ for the
$2\times2$ matrix $(w_{\varepsilon})_{\varepsilon_{j},\varepsilon_{j''}}$ with
$\varepsilon_{j'}=\sigma'$ fixed. Let $r^{k}(v)$ be the rank of the contraction
$HT^{k}(X)\to H^{*}(X,\CC)$, $\xi\mapsto\xi\lrcorner\,v$.

\begin{proposition}[The count in degree six]\label{prop:sexticcount}
Let $v\in S(t,q)$ be a nonzero real class.
\begin{enumerate}[label=\textup{(\roman*)},leftmargin=2.6em,itemsep=2pt]
\item $r^{1}(v)=12$,
  \[
    r^{2}(v)=4\,(A_{12}+A_{13}+A_{23})+\rho_{1}+\rho_{2}+\rho_{3},
    \qquad
    r^{3}(v)=8N_{w}+2\sum_{j\ne j'}\bigl(\rk M^{(j;j')}_{+}+\rk M^{(j;j')}_{-}\bigr).
  \]
  So $r^{2}(v)\le54$ and $r^{3}(v)\le112$, with equality when every
  $w_{\varepsilon}$ is nonzero and every $M^{(j;j')}_{\sigma'}$ is invertible,
  which holds outside a proper Zariski closed subset of $S(t,q)$. When $F$
  contains an imaginary quadratic field $K$ and $v$ lies in the secant plane
  of $K$, so that $N_{w}=2$, one has $(r^{2},r^{3})=(30,40)$.
\item $\chi(v,v)=\int_{X}v^{\vee}v=-64\,\Nm_{F_{0}/\QQ}(q)\sum_{\varepsilon}
  |w_{\varepsilon}|^{2}<0$. For a CM field of degree $2m$, with $X$ of dimension
  $2m$ and $2^{m}$ exponentials, the same computation gives
  $(-4)^{m}\Nm(q)\sum_{\varepsilon}|w_{\varepsilon}|^{2}$: the Euler form on
  the secant space is definite, of sign $(-1)^{m}$.
\item Let $F_{1},F_{2}$ be perfect complexes on $X$ with
  $\Ext^{<0}(F_{i},F_{i})=0$ and characters $v_{i}$ in spaces
  $S(t_{i},q_{i})$, let $A=X\times\widehat{X}$ carry the action of $F$ of
  \cite[\S11.1]{Mar25c}, transported by the exponential of a rational
  $F_{0}$-$B$-field, and put
  $E=\Phi(F_{1}\boxtimes F_{2}^{\vee})$ and $\kappa=\ch(E)e^{\ell/2}$. Then
  $r(\kappa)=r(\ch E)=r^{2}(v_{1})+144+r^{2}(v_{2})\le252$, while every class
  that stays of Hodge type over $\cD_{F}$ has rank at most
  $\dim HT^{2}(A)-\dim\cD_{F}=276-12=264$. Moreover
  $\dim\Ext^{2}(E,E)\ge r(\kappa)$, with equality if and only if each $F_{i}$
  has
  \[
    \dim\Ext^{k}(F_{i},F_{i})=1,\ 12,\ r^{2}(v_{i}),\ e_{3},\ r^{2}(v_{i}),\ 12,\ 1
    \quad(k=0,\dots,6),
  \]
  where $e_{3}=2r^{2}(v_{i})-22-\chi(v_{i},v_{i})$,
  and this profile is compatible with the bound $e_{3}\ge r^{3}(v_{i})$ exactly
  when $\chi(F_{i},F_{i})\le-(r^{3}(v_{i})-2r^{2}(v_{i})+22)$, that is
  $\chi(F_{i},F_{i})\le-26$ for $v_{i}$ general in its secant space, and
  $\chi(F_{i},F_{i})\le-2$ when $N_{w}=2$.
\item If $t_{1}=t_{2}=0$ and $q_{1}=q_{2}=q$, then $\kappa$ is invariant
  under $G_{F}$, and so is a Hodge class on every member of the
  $(F,2,1)$-family.
\end{enumerate}
\end{proposition}

\begin{proof}
(i) Choose a basis $x_{j1},x_{j2}$ of $U_{j}^{1,0}$ and $y_{j1},y_{j2}$ of
$U_{j}^{0,1}$ with $\theta_{j}=x_{j1}y_{j1}+x_{j2}y_{j2}$, and let
$\omega_{j}=y_{j1}y_{j2}$ span $\bigwedge^{2}U_{j}^{0,1}$. Then
$H^{*}(X,\CC)$ and $HT^{*}(X)=\bigwedge^{*}(H^{0,1}\oplus T)$ are the tensor
products over $j$ of the corresponding algebras of the three factors,
$HT^{k}(X)=\bigoplus_{a_{1}+a_{2}+a_{3}=k}HT^{a_{1}}_{1}\otimes HT^{a_{2}}_{2}
\otimes HT^{a_{3}}_{3}$, and contraction is compatible with the factorisation;
$v=\sum_{\varepsilon}w_{\varepsilon}\bigotimes_{j}e^{\lambda_{j}(\varepsilon_{j})\theta_{j}}$
with $\lambda_{j}(\pm1)=\tau_{j}(t)\pm s_{j}$, two distinct numbers. The blocks
of $HT^{k}(X)$ shift the three Hodge bidegrees of $H^{*}(X)$ by different
amounts, so their images are independent and $r^{k}(v)$ is the sum of their
ranks. Inside a block the bookkeeping rests on two facts about one factor:
the two classes $e^{\lambda_{j}(\pm1)\theta_{j}}$ are linearly independent,
and so are the subspaces $U_{j}^{0,1}e^{\lambda_{j}(\pm1)\theta_{j}}$ of
$H^{*}(U_{j})$, since $u\,e^{\lambda\theta_{j}}=u+\lambda\,u\theta_{j}$ for
$u\in U_{j}^{0,1}$; whereas $\omega_{j}e^{\lambda\theta_{j}}=\omega_{j}$ for
every $\lambda$, so that the summands of a block through $HT^{2}_{j}$ share the
factor $\omega_{j}$ and are combined, not juxtaposed. On one factor,
contraction into $e^{\lambda\theta_{j}}$ is as follows. For $\xi\in HT^{1}_{j}$,
$\xi\lrcorner\,e^{\lambda\theta_{j}}=y^{\lambda}(\xi)\,e^{\lambda\theta_{j}}$ with
$y^{\lambda}(\xi)=\xi^{0,1}+\lambda\,\xi_{T}\lrcorner\,\theta_{j}\in
U_{j}^{0,1}$, as in \Cref{lem:bfield}(i); $y^{\lambda}$ is onto, its kernel
$K^{\lambda}_{j}$ has dimension two, and
$HT^{1}_{j}=K^{\lambda_{j}(+1)}_{j}\oplus K^{\lambda_{j}(-1)}_{j}$ because
$\lambda_{j}(+1)\ne\lambda_{j}(-1)$ and $\xi_{T}\lrcorner\,\theta_{j}$
determines $\xi_{T}$. For $\xi\in HT^{2}_{j}$, since $\omega_{j}\theta_{j}=0$,
one finds $\omega_{j}\lrcorner\,e^{\lambda\theta_{j}}=\omega_{j}$,
$(y_{ja}\otimes\partial_{x_{jb}})\lrcorner\,e^{\lambda\theta_{j}}=\pm\lambda\,
\omega_{j}$ for $a\ne b$ and $0$ for $a=b$, and
$(\partial_{x_{j1}}\wedge\partial_{x_{j2}})\lrcorner\,e^{\lambda\theta_{j}}
=\pm\lambda^{2}\omega_{j}$: the image of $HT^{2}_{j}$ is the line
$\CC\omega_{j}$, $\xi\lrcorner\,e^{\lambda\theta_{j}}=c(\xi,\lambda)\,\omega_{j}$,
and the pairs $\bigl(c(\xi,\lambda_{j}(+1)),c(\xi,\lambda_{j}(-1))\bigr)$ fill
$\CC^{2}$. For $k\ge3$, $HT^{k}_{j}$ kills $e^{\lambda\theta_{j}}$, because
every element of it contains $\omega_{j}$ or $\partial_{x_{j1}}\wedge
\partial_{x_{j2}}$ together with a further factor, and the further factor
kills $\omega_{j}$.

Now count. The block $HT^{1}_{j}$ acts on $v$ by
$\xi\mapsto\sum_{\varepsilon}w_{\varepsilon}\,y^{\lambda_{j}(\varepsilon_{j})}(\xi)
\bigotimes_{j'}e^{\lambda_{j'}(\varepsilon_{j'})\theta_{j'}}$; on the summand
$K^{\lambda_{j}(-\sigma)}_{j}$ only the $\varepsilon$ with
$\varepsilon_{j}=\sigma$ contribute, and they map it isomorphically onto
$U_{j}^{0,1}$ tensored with a nonzero class as soon as one such
$w_{\varepsilon}$ is nonzero. Both signs occur in the support of a nonzero real
$w$, since $w_{-\varepsilon}=\bbar{w}_{\varepsilon}$, so the block has rank $4$
and $r^{1}(v)=12$. The block $HT^{1}_{j}\otimes HT^{1}_{j'}$ is the sum of the
four summands $K^{\lambda_{j}(-\sigma)}_{j}\otimes K^{\lambda_{j'}(-\sigma')}_{j'}$,
each of dimension four; the summand indexed by $(\sigma,\sigma')$ maps
isomorphically onto $U_{j}^{0,1}\otimes U_{j'}^{0,1}$ tensored with the class
$\sum_{\sigma''}w_{\varepsilon}e^{\lambda_{j''}(\sigma'')\theta_{j''}}$, which
is nonzero exactly when the pair $(\sigma,\sigma')$ occurs in the support, and
to zero otherwise; so the block has rank $4A_{jj'}$. The block $HT^{2}_{j}$
maps $\xi$ to $\sum_{\sigma}c(\xi,\lambda_{j}(\sigma))\,z_{\sigma}$ with
$z_{\sigma}=\omega_{j}\otimes\sum_{\mu}w_{\sigma\mu}e_{\mu}$, where $\mu$ runs
over the four sign patterns of the other two places and the $e_{\mu}$ are
independent; as $(c(\xi,\lambda_{j}(+1)),c(\xi,\lambda_{j}(-1)))$ fills
$\CC^{2}$, the rank is the dimension of the span of $z_{+},z_{-}$, which is
$\rho_{j}$. This gives $r^{2}(v)$. In degree three the block
$HT^{1}_{1}\otimes HT^{1}_{2}\otimes HT^{1}_{3}$ is the sum of eight summands
$\bigotimes_{j}K^{\lambda_{j}(-\varepsilon_{j})}_{j}$, of dimension eight, on
each of which exactly one term of $v$ acts, isomorphically when
$w_{\varepsilon}\ne0$ and by zero otherwise: rank $8N_{w}$. The block
$HT^{2}_{j}\otimes HT^{1}_{j'}$ factors through $\CC^{2}\otimes HT^{1}_{j'}$
by $c\otimes\mathrm{id}$, which is onto, and on
$\CC^{2}\otimes K^{\lambda_{j'}(-\sigma')}_{j'}$ the map is
$M^{(j;j')}_{\sigma'}\otimes y^{\lambda_{j'}(\sigma')}$ followed by tensoring
with independent classes, of rank $2\rk M^{(j;j')}_{\sigma'}$; the two values
of $\sigma'$ give independent images. The blocks $HT^{3}_{j}$ contribute
nothing. This gives $r^{3}(v)$. The maximal values are $48+6=54$ and
$64+6\cdot8=112$; the conditions for them are the nonvanishing of the
$w_{\varepsilon}$ and of the determinants of the $M^{(j;j')}_{\sigma'}$, which
are polynomials in the coordinates of $v$, nonzero on $S(t,q)$. If $F\supseteq K$
and $v$ lies in the secant plane of $K$, the support of $w$ is a pair
$\{\varepsilon_{0},-\varepsilon_{0}\}$, the two CM types of $F$ induced from
$K$; then $A_{jj'}=2$, $\rho_{j}=2$, $N_{w}=2$ and every
$M^{(j;j')}_{\sigma'}$ has one nonzero entry, so $(r^{2},r^{3})=(24+6,16+24)$.

(ii) Since $\theta_{j}^{3}=0$, $\theta^{2m}=\frac{(2m)!}{2^{m}}\theta_{1}^{2}
\cdots\theta_{m}^{2}$, and $\int_{X}\theta^{2m}/(2m)!=1$ gives
$\int_{X}\theta_{1}^{2}\cdots\theta_{m}^{2}=2^{m}$ and
$\int_{X}\exp(\sum_{j}c_{j}\theta_{j})=\prod_{j}c_{j}^{2}$. The dual of
$e^{\lambda\theta_{j}}$ is $e^{-\lambda\theta_{j}}$, so
$\int_{X}e_{\varepsilon}^{\vee}e_{\varepsilon'}=\prod_{j}\bigl(\lambda_{j}(\varepsilon'_{j})
-\lambda_{j}(\varepsilon_{j})\bigr)^{2}$, which is $0$ unless
$\varepsilon'=-\varepsilon$, where it is $\prod_{j}(2s_{j})^{2}=(-4)^{m}\Nm(q)$.
Hence $\chi(v,v)=(-4)^{m}\Nm(q)\sum_{\varepsilon}w_{\varepsilon}w_{-\varepsilon}
=(-4)^{m}\Nm(q)\sum_{\varepsilon}|w_{\varepsilon}|^{2}$.

(iii) As in the proof of \Cref{thm:quarticobstruction}(i), $r(\kappa)=r(\ch E)$
by \Cref{lem:bfield}(i), and $\Phi$ identifies $r(\ch E)$ with the rank of
contraction from $HT^{2}(X\times X)=\bigoplus_{a+b=2}HT^{a}(X)\otimes HT^{b}(X)$
into $v_{1}\boxtimes v_{2}$; the block $HT^{1}\otimes HT^{1}$ contributes
$r^{1}(v_{1})r^{1}(v_{2})=144$ and the other two contribute
$r^{2}(v_{1})+r^{2}(v_{2})$, by the type argument given there. The tangent
space $T_{[A]}\cD_{F}\subseteq H^{1}(T_{A})$, of dimension $mn^{2}=12$, kills
every class that stays of Hodge type along $\cD_{F}$, as in
\Cref{prop:p2primelocus}(i), and $\dim HT^{2}(A)=66+144+66$. By the K\"unneth
formula, with $\Ext^{<0}=0$,
$\dim\Ext^{2}(E,E)=e_{0}e_{2}'+e_{1}e_{1}'+e_{2}e_{0}'$ where
$e_{k}=\dim\Ext^{k}(F_{1},F_{1})$ and $e_{k}'=\dim\Ext^{k}(F_{2},F_{2})$; the
commutative square $\sigma\circ\operatorname{ev}=(\,\cdot\,)\lrcorner\ch$ of
\Cref{setup:criterion}, valid in every degree, gives $e_{k}\ge r^{k}(v_{1})$
and $e_{k}'\ge r^{k}(v_{2})$, and $e_{0},e_{0}'\ge1$. So
$\dim\Ext^{2}(E,E)\ge r^{2}(v_{2})+144+r^{2}(v_{1})$, with equality if and
only if $e_{0}=e_{0}'=1$, $e_{1}=e_{1}'=12$, $e_{2}=r^{2}(v_{1})$ and
$e_{2}'=r^{2}(v_{2})$. Serre duality on the sixfold gives $e_{k}=e_{6-k}$, so
$\chi(F_{i},F_{i})=2(e_{0}-e_{1}+e_{2})-e_{3}=2r^{2}(v_{i})-22-e_{3}$, and by
Riemann--Roch on an abelian variety $\chi(F_{i},F_{i})=\chi(v_{i},v_{i})$,
which is (ii); this determines $e_{3}$, and $e_{3}\ge r^{3}(v_{i})$ is the
stated inequality. With the maximal ranks it reads
$\chi\le-(112-108+22)=-26$, and with $N_{w}=2$ it reads
$\chi\le-(40-60+22)=-2$.

(iv) The proof of \Cref{thm:quarticobstruction}(iii) applies with three
factors in place of two: for $t=0$ the classes $e_{\varepsilon}\boxtimes
e_{\varepsilon'}^{\vee}$, the correspondence $Z\mapsto\ch(\Phi(Z))e^{\ell/2}$
and the Poincar\'e class decompose as tensor products over $j$ of classes of
even degree, on each factor $\sqrt{-q}$ acts as in the imaginary quadratic
construction with $d$ replaced by $\tau_{j}(q)$, and the identities of
\Cref{prop:flatall} at $n=2$ show that each image is a multiple of an
exponential of $\eta_{j}$ or of a Weil class of the $j$-th factor, invariant
under the special unitary group of that factor; $G_{F}(\RR)$ is the product of
the three, and is Zariski dense in $G_{F}$. Statement (iv) is the case
$[F:\QQ]=6$ of \cite[Proposition~10.2.1(1)]{Mar25c}, with the inclusion
$\widetilde\phi(B\otimes B)\subset A$ of \cite[p.~33]{Mar25c}, which Markman
proves for every CM field in his conventions for $\Phi$ and for the twist; we
make no claim of priority.
\end{proof}

\begin{remark}[The special role of degree four]
\label{rem:sexticparity}
In \Cref{thm:quarticobstruction} the equality $\dim\Ext^{2}(E,E)=r(\kappa)$
determines the whole Ext profile of each factor, because on a fourfold
$\Ext^{2}$ is the middle degree and Serre duality supplies the rest; the Euler
characteristic of a minimal factor is then the small positive integer
$r^{2}-14$, the Euler form on the secant space being positive definite, and
the arithmetic of \Cref{prop:quarticother} excludes it. On a sixfold the same
equality fixes $\Ext^{0}$, $\Ext^{1}$ and $\Ext^{2}$ and leaves the middle
group $\Ext^{3}$ free, and the Euler form has the opposite sign:
$\chi(F_{i},F_{i})=-64\Nm(q)\sum|w_{\varepsilon}|^{2}$ is a negative integer
that $e_{3}$ absorbs. So no dimension count excludes a minimal Orlov product
over a sextic field, and none does over any CM field of degree $2m\ge6$,
where $\Ext^{3},\dots,\Ext^{m}$ are free and the Euler form has sign
$(-1)^{m}$. What \Cref{prop:sexticcount}(iii) leaves for such an object is the
profile displayed there, with $\chi(F_{i},F_{i})\le-26$ for a general secant
character, and the flat twisted character of (iv), which stays of Hodge type
over the whole twelve-dimensional period domain. Whether a sheaf with that
profile exists on a sixfold with real multiplication by a cubic field is not
known; nothing in \Cref{thm:quarticobstruction} bears on it. A certificate
of the kind of \Cref{thm:twistedcert} for a sextic field would be an Orlov
product of such sheaves, with $\dim\Ext^{2}(E,E)=r(\kappa)\le252$ on the
twelvefold, whose twisted character has a nonzero component in the Weil
classes $W_{F}(A)$; flatness does not give that component, and
\Cref{lem:sexticweil} computes it. In degree six as in degree four, the Orlov
products never reach the rank of a class of general shape: $r(\kappa)\le252$
against $264$, and $r(\kappa)\le104$ against $112$.
\end{remark}

The count of \Cref{prop:sexticcount}(iii) is a condition on the Euler
characteristic of the character alone, so it can be tested on the lattice of
integral characters, as \Cref{prop:quarticother} tests the quartic one. In
degree four the test excludes every minimal character. In degree six it can
exclude an integral character only if $-26<\chi(v,v)<0$
(\Cref{prop:sexticintegral}(i),(ii)), and it excludes neither of the integral
classes $\operatorname{Re}e^{\sqrt{-1}\,\theta}$ and
$\operatorname{Im}e^{\sqrt{-1}\,\theta}$ (\Cref{prop:sexticintegral}(iv)). For
$t=0$, the case in which
\Cref{prop:sexticcount}(iv) makes $\kappa$ flat, put
$B_{j}=1-\tfrac12\tau_{j}(q)\,\theta_{j}^{2}$, so that
$e^{\pm s_{j}\theta_{j}}=B_{j}\pm s_{j}\theta_{j}$, and for
$c_{0},c_{3}\in\QQ$ and $f,g\in F_{0}$ let
\[
  v(c_{0},f,g,c_{3})=c_{0}B_{1}B_{2}B_{3}
  +\sum_{j}\tau_{j}(f)\,\theta_{j}B_{k}B_{l}
  +\sum_{j}\tau_{j}(g)\,B_{j}\theta_{k}\theta_{l}
  +c_{3}\,\theta_{1}\theta_{2}\theta_{3},
\]
where $\{j,k,l\}=\{1,2,3\}$ in each sum.

\begin{proposition}[Integral flat characters in degree six]
\label{prop:sexticintegral}
Keep the notation of \Cref{prop:sexticcount}.
\begin{enumerate}[label=\textup{(\roman*)},leftmargin=2.6em,itemsep=2pt]
\item For every perfect complex $F$ on a complex torus of dimension six,
  $\chi(F,F)$ is even. So the condition of \Cref{prop:sexticcount}(iii)
  holds for every $F$ whose character has $N_{w}=2$.
\item For every nonzero real $v\in S(t,q)$, $r^{3}(v)-2r^{2}(v)\le4$, with
  equality only if $(r^{2}(v),r^{3}(v))=(54,112)$. So for every shape the
  condition of \Cref{prop:sexticcount}(iii) is implied by
  $\chi(F_{i},F_{i})\le-26$.
\item The rational points of $S(0,q)$ are the classes $v(c_{0},f,g,c_{3})$,
  and
  \[
    \chi(v,v)=-8\bigl(\Nm(q)\,c_{0}^{2}+\Tr\bigl(\Nm(q)\,q^{-1}f^{2}\bigr)
    +\Tr(q\,g^{2})+c_{3}^{2}\bigr).
  \]
\item For $q=1$, the classes $\operatorname{Re}e^{\sqrt{-1}\,\theta}=v(1,0,-1,0)$
  and $\operatorname{Im}e^{\sqrt{-1}\,\theta}=v(0,1,0,-1)$ are integral
  points of $S(0,1)$ with $\chi=-32$ and $N_{w}=2$, so they satisfy the
  condition of \Cref{prop:sexticcount}(iii).
\end{enumerate}
\end{proposition}

\begin{proof}
(i) The Chern character $v=\ch(F)$ is integral, as in the proof of
\Cref{thm:quarticobstruction}(ii), and $\chi(F,F)=\int_{X}v^{\vee}v$ by
Riemann--Roch. Write $v=\sum_{k}v_{2k}$, so that
$\int_{X}v^{\vee}v=\sum_{k}(-1)^{k}\int_{X}v_{2k}v_{12-2k}$. The terms $k$ and
$6-k$ are equal, since $(-1)^{k}=(-1)^{6-k}$ and classes of even degree
commute, so they contribute an even number when $k\ne3$. For $k=3$ write
$v_{6}=\sum_{I}a_{I}e_{I}$ in the basis of $\bigwedge^{6}H^{1}(X,\ZZ)$ given
by the products $e_{I}$ of six elements of a basis of $H^{1}(X,\ZZ)$; since
$e_{I}e_{I}=0$ and $e_{I}e_{J}=e_{J}e_{I}$, $v_{6}^{2}=2\sum_{I<J}a_{I}a_{J}
e_{I}e_{J}$ is even. For $N_{w}=2$ the condition of
\Cref{prop:sexticcount}(iii) is $\chi(F,F)\le-2$, and $\chi(F,F)$ is a
negative even integer by \Cref{prop:sexticcount}(ii).

(ii) For $j$ and $\sigma=\pm1$ let $W_{j,\sigma}$ be the $2\times2$ slice
$(w_{\varepsilon})_{\varepsilon_{j}=\sigma}$, with rows and columns indexed by
the signs at the two other places. The matrices $M^{(j';j)}_{\sigma}$ are
$W_{j,\sigma}$ and its transpose, so
$\sum_{j\ne j'}\bigl(\rk M^{(j;j')}_{+}+\rk M^{(j;j')}_{-}\bigr)
=2\sum_{j,\sigma}\rk W_{j,\sigma}$. Since $w_{-\varepsilon}=\bbar{w}_{\varepsilon}$,
$W_{j,-}$ is the complex conjugate of $W_{j,+}$ with its rows and its columns
reversed, so the two slices have the same rank $k_{j}$. Let $S$ be the support
of $w$ and $N=N_{w}=|S|$. By \Cref{prop:sexticcount}(i),
\begin{equation}\label{eq:sexticexcess}
  r^{3}(v)-2r^{2}(v)=8N-8\,(A_{12}+A_{13}+A_{23})
  +\sum_{j}\bigl(8k_{j}-2\rho_{j}\bigr).
\end{equation}
The set $S$ is stable under $\varepsilon\mapsto-\varepsilon$, which exchanges
the two slices along each place; so $N$ is even, each slice has $N/2$ nonzero
entries, $k_{j}\ge1$, and $A_{jj'}\in\{2,4\}$, the projection of $S$ to the
signs at $j$ and $j'$ being stable under negation. Also $1\le\rho_{j}\le2$ and
$k_{j}\le2$, and a slice whose nonzero entries lie in one row or one column
has rank one. If $\rho_{j}=1$, then $w_{\varepsilon}=a_{\varepsilon_{j}}
W(\varepsilon')$, with $\varepsilon'$ the signs at the two other places, so
every slice along a place $j'\ne j$ is $a\otimes(\text{a row of }W)$ and
$k_{j'}\le1$. Now take the four values of $N$ in turn.

If $N=8$, every $A_{jj'}$ is $4$ and \eqref{eq:sexticexcess} is
$-32+\sum_{j}(8k_{j}-2\rho_{j})$. When every $\rho_{j}$ is $2$ this is at most
$-32+3\cdot12=4$, with equality exactly when every $k_{j}$ is $2$, and then
$r^{2}=48+6=54$ and $r^{3}=64+48=112$. When some $\rho_{j}$ is $1$, the two
other $k$ are at most $1$ and \eqref{eq:sexticexcess} is at most
$-32+14+6+6=-6$.

If $N=6$, $S$ misses one pair $\{\pm\varepsilon_{0}\}$. As $\varepsilon_{0}$
and $-\varepsilon_{0}$ differ at every place, they never have the same signs
at two places, so every sign pair at two places occurs in $S$, every
$A_{jj'}$ is $4$, and \eqref{eq:sexticexcess} is at most $48-96+3\cdot14=-6$.

If $N=4$, $S=\{\pm\varepsilon,\pm\varepsilon'\}$, and after replacing
$\varepsilon'$ by $-\varepsilon'$ if necessary, $\varepsilon$ and
$\varepsilon'$ differ at exactly one place $j$. For the two other places $k$
and $l$ one has $A_{kl}=2$ and $A_{jk}=A_{jl}=4$. The two nonzero entries of
a slice along $k$ are $w_{\pm\varepsilon}$ and $w_{\pm\varepsilon'}$ with the
same sign, whose indices agree at $l$; they lie in one row or one column, so
$k_{k}=1$, and likewise $k_{l}=1$. Hence \eqref{eq:sexticexcess} is at most
$32-80+8\,(2+1+1)-6=-22$.

If $N=2$, $S=\{\pm\varepsilon_{0}\}$, every $A_{jj'}$ is $2$, every $k_{j}$ is
$1$, and every $\rho_{j}$ is $2$, the two slices along $j$ having their
nonzero entries at different positions; so \eqref{eq:sexticexcess} is
$16-48+24-12=-20$, the value $40-60$ of \Cref{prop:sexticcount}(i).

So $r^{3}-2r^{2}\le4$, with equality only at $(54,112)$, and the threshold
$r^{3}-2r^{2}+22$ of \Cref{prop:sexticcount}(iii) is at most $26$.

(iii) Since $\theta_{j}^{3}=0$ and $s_{j}^{2}=-\tau_{j}(q)$,
$e^{\pm s_{j}\theta_{j}}=B_{j}\pm s_{j}\theta_{j}$, so the eight exponentials
and the eight products $\prod_{j}\beta_{j}$ with $\beta_{j}\in\{B_{j},\theta_{j}\}$
span the same complex space, the change of basis being the tensor cube of an
invertible $2\times2$ matrix. The Galois group permutes the $\tau_{j}$, the
$\theta_{j}$ and the $B_{j}$ compatibly, so $\sum_{a}C_{a}\prod_{j}\beta_{j}$,
indexed by the subsets $a$ of $\{1,2,3\}$ with $\beta_{j}=\theta_{j}$ for
$j\in a$, is rational exactly when $C_{\pi(a)}=\sigma(C_{a})$ for every
$\sigma$ inducing $\pi$; the stabiliser of $\{j\}$, and of its complement,
is the stabiliser of $\tau_{j}$, which gives the form $v(c_{0},f,g,c_{3})$.
The involution $v\mapsto v^{\vee}$ and the integral are multiplicative over
the three factors, with $\int_{X}\theta_{1}^{2}\theta_{2}^{2}\theta_{3}^{2}=8$;
on one factor $B_{j}^{\vee}B_{j}=1-\tau_{j}(q)\theta_{j}^{2}$,
$\theta_{j}^{\vee}\theta_{j}=-\theta_{j}^{2}$ and the mixed products are odd
in $\theta_{j}$. So
$\chi(v,v)=-8\sum_{a}C_{a}^{2}\prod_{j\notin a}\tau_{j}(q)$, which is the
formula.

(iv) For $q=1$, $s_{j}=\sqrt{-1}$ and $e_{(1,1,1)}=\prod_{j}e^{\sqrt{-1}\,
\theta_{j}}=e^{\sqrt{-1}\,\theta}$. The polarisation is principal, so each
$\theta^{k}/k!$ is integral, and so are the real and imaginary parts
$\sum_{k}(-1)^{k}\theta^{2k}/(2k)!$ and $\sum_{k}(-1)^{k}\theta^{2k+1}/(2k+1)!$.
They are $\tfrac12(e_{(1,1,1)}+e_{(-1,-1,-1)})$ and
$\tfrac1{2\sqrt{-1}}(e_{(1,1,1)}-e_{(-1,-1,-1)})$, so $N_{w}=2$ and
$\chi=-64\,(\tfrac14+\tfrac14)=-32$ by \Cref{prop:sexticcount}(ii); expanding
$\prod_{j}(B_{j}+\sqrt{-1}\,\theta_{j})$ gives the coordinates.
\end{proof}

The flatness of \Cref{prop:sexticcount}(iv) makes $\kappa$ a Hodge class on
every member of the family, but in degree four the Hodge classes of a very
general member are the products of divisor classes and the Weil classes
$W_{F}(A)$ (\Cref{prop:cmweil}(iii)), and only a nonzero component of
$\kappa$ in $W_{F}(A)$ bears on $\mathbf{W}(F,2,1)$. We compute that
component. Keep the notation of \Cref{prop:sexticcount}, with $t_{1}=t_{2}=0$,
$q_{1}=q_{2}=q$ and trivial $B$-field, put $d_{j}=\tau_{j}(q)$, and let
$\Psi(Z)=\ch(\Phi(Z))e^{\ell/2}$ as in \Cref{prop:flatall}, so that
$\kappa(v_{1},v_{2})=\Psi(v_{1}\boxtimes v_{2}^{\vee})$ for the characters
$v_{i}$ of $F_{i}$. As in the proof of \Cref{prop:sexticcount}(iv),
$H^{1}(A,\CC)$ is the sum of three factors $U_{j}\oplus U_{j}^{\vee}$, on each
of which $\sqrt{-q}$ acts as in the imaginary quadratic construction with
$\beta=\theta_{j}$ and $d=d_{j}$. On the $j$-th factor write
$\widehat\theta_{j}$ and $\ell_{j}$ for the components of the dual
polarisation and of the Poincar\'e class, and put
$\gamma_{j}=d_{j}\theta_{j}-\widehat\theta_{j}$, as in \eqref{eq:gammadef},
and $\eta_{j}=d_{j}\theta_{j}+\widehat\theta_{j}$, a positive multiple of the
polarisation of the factor for which $\sqrt{-q}$ is anti-self-adjoint
(\Cref{prop:splitgeom}(i)). For $\sigma=\pm1$ let $\tau_{j,\sigma}$ be the
embedding of $F$ over $\tau_{j}$ with
$\tau_{j,\sigma}(\sqrt{-q})=\sigma\sqrt{-1}\,\sqrt{\tau_{j}(q)}=\sigma s_{j}$,
and $V_{j,\sigma}\subset H^{1}(A,\CC)$ its eigenspace, of dimension four. By
\Cref{thm:splitclosed}(ii) with $n=2$ and $d=d_{j}$, the line
$\bigwedge^{4}V_{j,\sigma}$ is spanned by
\begin{equation}\label{eq:sexticomega}
  \omega_{j,\sigma}=-\frac{1}{2d_{j}^{2}}\bigl(\gamma_{j}-\sigma s_{j}\ell_{j}\bigr)^{2},
\end{equation}
so the six classes $\omega_{j,\sigma}$ span $W_{F}(A)\otimes\CC$. Since
$\eta_{j}$ pairs $V_{j,+}$ with $V_{j,-}$, the products $\eta_{i}\eta_{k}$ lie
in summands of $H^{4}(A,\CC)=\bigwedge^{4}\bigl(\bigoplus_{j,\sigma}
V_{j,\sigma}\bigr)$ other than the six lines $\bigwedge^{4}V_{j,\sigma}$. Let
$\kappa_{j,\sigma}(v_{1},v_{2})$ be the coefficient of $\omega_{j,\sigma}$ in
the component of $\kappa(v_{1},v_{2})$ in $\bigwedge^{4}V_{j,\sigma}$; the
\emph{$F$-Weil part} of $\kappa$ is
$\sum_{j,\sigma}\kappa_{j,\sigma}\,\omega_{j,\sigma}$, its component in
$W_{F}(A)\otimes\CC$. For $\varepsilon\in\{\pm1\}^{3}$ let
$\varepsilon^{(j)}$ be the sign vector that agrees with $\varepsilon$ at $j$
and only there.

\begin{lemma}[The $F$-Weil part of an Orlov product]\label{lem:sexticweil}
Let $v_{i}=\sum_{\varepsilon}w^{i}_{\varepsilon}e_{\varepsilon}\in S(0,q)$ for
$i=1,2$.
\begin{enumerate}[label=\textup{(\roman*)},leftmargin=2.6em,itemsep=2pt]
\item The degree-four part of $\kappa(v_{1},v_{2})$ is a quadratic form in the
  $\eta_{j}$ plus $\sum_{j,\sigma}\kappa_{j,\sigma}\,\omega_{j,\sigma}$, with
  \[
    \kappa_{j,\sigma}(v_{1},v_{2})=-16\Nm(q)\,\tau_{j}(q)
    \sum_{\varepsilon_{j}=\sigma}w^{1}_{\varepsilon}\,w^{2}_{\varepsilon^{(j)}}.
  \]
\item Let $v=v(c_{0},f,g,c_{3})$ be rational, and for $\{j,k,l\}=\{1,2,3\}$
  write $f_{j}=\tau_{j}(f)$, $g_{j}=\tau_{j}(g)$ and $q_{j}=\tau_{j}(q)$. The
  numbers
  \begin{align*}
    R_{j}&=c_{0}^{2}+\frac{f_{k}^{2}}{q_{k}}+\frac{f_{l}^{2}}{q_{l}}
      +\frac{g_{j}^{2}}{q_{k}q_{l}}-\frac{1}{q_{j}}\Bigl(f_{j}^{2}
      +\frac{g_{l}^{2}}{q_{k}}+\frac{g_{k}^{2}}{q_{l}}
      +\frac{c_{3}^{2}}{q_{k}q_{l}}\Bigr),\\
    I_{j}&=c_{0}f_{j}+\frac{f_{k}g_{l}}{q_{k}}+\frac{f_{l}g_{k}}{q_{l}}
      +\frac{c_{3}g_{j}}{q_{k}q_{l}}
  \end{align*}
  are $\tau_{j}(R)$ and $\tau_{j}(I)$ for elements $R,I\in F_{0}$. Put
  $\Omega(v)=R+2I/\sqrt{-q}\in F$. Then
  \[
    \kappa_{j,\sigma}(v,v)=-\Nm(q)\,\tau_{j}(q)\,\tau_{j,\sigma}\bigl(\Omega(v)\bigr).
  \]
  So the $F$-Weil part of $\kappa(v,v)$ is zero if $R=I=0$, and otherwise
  all six $\kappa_{j,\sigma}(v,v)$ are nonzero.
\item If $v_{1}$ and $v_{2}$ are supported on one pair
  $\{\varepsilon_{0},-\varepsilon_{0}\}$, the $F$-Weil part of
  $\kappa(v_{1},v_{2})$ is zero. For $q=1$, $F$ contains
  $K=\QQ(\sqrt{-1})$, the classes $\operatorname{Re}e^{\sqrt{-1}\,\theta}$ and
  $\operatorname{Im}e^{\sqrt{-1}\,\theta}$ span the plane $P_{\theta}$ of
  \Cref{def:secantplane} with $d=1$, and for
  $v_{i}=a_{i}\operatorname{Re}e^{\sqrt{-1}\,\theta}
  +b_{i}\operatorname{Im}e^{\sqrt{-1}\,\theta}$ with $a_{i},b_{i}\in\QQ$ the
  class $\kappa(v_{1},v_{2})$ lies in $\QQ[\eta]\oplus W_{K}(A)$, where
  $\eta=\theta+\widehat\theta$ and $W_{K}(A)\subseteq H^{12}(A,\QQ)$ are the
  Weil classes of $K$ on the twelvefold (\Cref{prop:composite}); its part of
  degree at most four is
  \[
    -32(a_{1}a_{2}+b_{1}b_{2})+16(a_{1}b_{2}-a_{2}b_{1})\,\eta
    +4(a_{1}a_{2}+b_{1}b_{2})\,\eta^{2}.
  \]
\end{enumerate}
\end{lemma}

\begin{proof}
(i) The classes $e_{\varepsilon}\boxtimes e_{\varepsilon'}^{\vee}
=\bigotimes_{j}\bigl(e^{\varepsilon_{j}s_{j}\theta_{j}}\boxtimes
e^{-\varepsilon'_{j}s_{j}\theta_{j}}\bigr)$ and $\Psi$ decompose over the three
factors, as in the proof of \Cref{prop:sexticcount}(iv), and on the $j$-th
factor the two identities of \Cref{prop:flatall} at $n=2$ read, with their
constants,
\[
  \Psi_{j}\bigl(e^{\sigma s_{j}\theta_{j}}\boxtimes e^{\sigma s_{j}\theta_{j}}\bigr)
  =-4d_{j}\,e^{\sigma(s_{j}/2d_{j})\eta_{j}},
  \qquad
  \Psi_{j}\bigl(e^{\sigma s_{j}\theta_{j}}\boxtimes e^{-\sigma s_{j}\theta_{j}}\bigr)
  =\tfrac12\bigl(\gamma_{j}-\sigma s_{j}\ell_{j}\bigr)^{2}
  =-d_{j}^{2}\,\omega_{j,\sigma},
\]
since $s_{j}^{2}=-d_{j}$. So
$\Psi(e_{\varepsilon}\boxtimes e_{\varepsilon'}^{\vee})$ is the product over
$j$ of the first class where $\varepsilon'_{j}=-\varepsilon_{j}$ and of the
second where $\varepsilon'_{j}=\varepsilon_{j}$. In degree four a product with
two or three factors of the second kind vanishes, one with none is a
quadratic form in the $\eta_{j}$, and one with exactly one, at $j$, has
$\varepsilon'=\varepsilon^{(j)}$ and equals
$-d_{j}^{2}\prod_{i\ne j}(-4d_{i})\,\omega_{j,\varepsilon_{j}}
=-16\Nm(q)\,d_{j}\,\omega_{j,\varepsilon_{j}}$. Summing
$w^{1}_{\varepsilon}w^{2}_{\varepsilon'}\Psi(e_{\varepsilon}\boxtimes
e_{\varepsilon'}^{\vee})$ over $\varepsilon$ and $\varepsilon'$ gives (i).

(ii) An element $h$ of the Galois group that fixes the place $j$ permutes
the two others, so it fixes $R_{j}$ and $I_{j}$, which are symmetric in $k$
and $l$; hence $R_{j},I_{j}\in\tau_{j}(F_{0})$, and moving $j$ by the Galois
group shows that they are the conjugates of two elements of $F_{0}$. Write
$\varepsilon=(\sigma,\mu)$ with $\mu$ the signs at $k$ and $l$, so that
$\varepsilon^{(j)}=(\sigma,-\mu)$. Since $e^{\pm s_{i}\theta_{i}}=B_{i}\pm
s_{i}\theta_{i}$, the class $v=\sum_{a}C_{a}\prod_{i\in a}\theta_{i}
\prod_{i\notin a}B_{i}$, with $C_{\emptyset}=c_{0}$, $C_{\{i\}}=f_{i}$,
$C_{\{1,2,3\}\setminus\{i\}}=g_{i}$ and $C_{\{1,2,3\}}=c_{3}$, has
$w_{\varepsilon}=\frac18\sum_{a}C_{a}\prod_{i\in a}\varepsilon_{i}/s_{i}$.
Grouping $a=b$ and $a=b\cup\{j\}$ for $b\subseteq\{k,l\}$ gives
$w_{(\sigma,\mu)}=\frac18\sum_{b}D_{b}\prod_{i\in b}\mu_{i}/s_{i}$ with
$D_{b}=C_{b}+u\,C_{b\cup\{j\}}$ and $u=\sigma/s_{j}=\tau_{j,\sigma}(1/\sqrt{-q})$.
Since $\sum_{\mu}\prod_{i\in b}\mu_{i}\prod_{i\in b'}(-\mu_{i})
=4(-1)^{|b|}\delta_{bb'}$ and $s_{i}^{2}=-q_{i}$,
\[
  \sum_{\varepsilon_{j}=\sigma}w_{\varepsilon}w_{\varepsilon^{(j)}}
  =\frac1{16}\Bigl((c_{0}+uf_{j})^{2}+\frac{(f_{k}+ug_{l})^{2}}{q_{k}}
  +\frac{(f_{l}+ug_{k})^{2}}{q_{l}}+\frac{(g_{j}+uc_{3})^{2}}{q_{k}q_{l}}\Bigr),
\]
and expanding with $u^{2}=-1/q_{j}$ gives $\frac1{16}(R_{j}+2uI_{j})
=\frac1{16}\tau_{j,\sigma}(\Omega(v))$. With (i) this is the formula. The
embedding $\tau_{j,\sigma}$ is injective and $1$, $1/\sqrt{-q}$ are
independent over $F_{0}$, so one $\kappa_{j,\sigma}(v,v)$ vanishes if and only
if $R=I=0$, and then all do.

(iii) If $w^{1}$ and $w^{2}$ vanish off $\{\varepsilon_{0},-\varepsilon_{0}\}$,
no term of (i) survives, since $\varepsilon^{(j)}$ is neither $\varepsilon$
nor $-\varepsilon$. For $q=1$, $s_{j}=\sqrt{-1}$ and
$e_{(1,1,1)}=e^{\sqrt{-1}\,\theta}$, whose real and imaginary parts are
$u_{\theta}$ and $v_{\theta}$ of \eqref{eq:uv} with $d=1$, supported on
$\{\pm(1,1,1)\}$. On each factor $\sqrt{-1}=\sqrt{-q}$ acts as in the
imaginary quadratic construction with $\beta=\theta_{j}$ and $d=1$, so $K$
acts on $A$ by the action \eqref{eq:Mdef} determined by $\theta$ with $d=1$,
and \Cref{prop:flatall} with $n=6$ and $d=1$ gives
$\kappa\in\QQ[\eta]\oplus W_{K}(A)$. For the constants write
$v_{1}=w_{1}e^{\sqrt{-1}\,\theta}+\bbar w_{1}e^{-\sqrt{-1}\,\theta}$ and
$v_{2}^{\vee}=a_{2}\operatorname{Re}e^{\sqrt{-1}\,\theta}
-b_{2}\operatorname{Im}e^{\sqrt{-1}\,\theta}
=w_{2}e^{\sqrt{-1}\,\theta}+\bbar w_{2}e^{-\sqrt{-1}\,\theta}$, with
$w_{1}=\frac12(a_{1}-\sqrt{-1}\,b_{1})$ and $w_{2}=\frac12(a_{2}+\sqrt{-1}\,
b_{2})$. By the first identity in the proof of (i),
$\Psi(e^{\pm\sqrt{-1}\,\theta}\boxtimes e^{\pm\sqrt{-1}\,\theta})
=-64\,e^{\pm(\sqrt{-1}/2)\eta}$, while the two mixed terms have degree
twelve; so in degrees below twelve $\kappa$ is
$-128\operatorname{Re}\bigl(w_{1}w_{2}\,e^{(\sqrt{-1}/2)\eta}\bigr)$, whose
terms of degree $0$, $2$ and $4$ are the stated ones.
\end{proof}

Markman states his reduction for the normalised class
$\ch(E)\exp(-c_{1}(E)/\rk E)$, and it is the component of degree four of that
class, a polynomial in the invariant divisor classes plus a Weil class, whose
algebraicity on the deformations he needs \cite[Proposition~10.2.1,
Corollary~10.2.3]{Mar25c}. When $\rk E\ne0$ that class is
$\kappa\exp(-\kappa_{1}/\rk E)$, with $\kappa_{1}$ the degree-two part of
$\kappa$, which is a linear form in the $\eta_{j}$ by the proof of
\Cref{lem:sexticweil}(i); so it has the same $F$-Weil part as $\kappa$.

For $q=1$ the Orlov products of complexes with characters in
$\QQ\operatorname{Re}e^{\sqrt{-1}\,\theta}\oplus\QQ\operatorname{Im}
e^{\sqrt{-1}\,\theta}$ are therefore the template of
\Cref{prop:orlovequality} at $(n,d)=(6,1)$, seen on the twelvefold through the
restriction of the action of $F$ to $K$. Their twisted characters are flat
for the family of $K$ and have a nonzero $K$-Weil part in $H^{12}$ when
$F_{1}=F_{2}\ne0$ (\Cref{prop:flatall}), and they are not an input to
$\mathrm{P2}_{\mathrm{split}}$. By \Cref{prop:composite}, $\mathbf{W}(F,2,1)$
gives the $K$-Weil classes on the members of the $F$-family, and the converse
is not available; a consequence for $K$ through \Cref{prop:descent} would
still need a twisted form of the propagation of \Cref{thm:semiregclosure}.

\begin{remark}[Natural objects]\label{rem:sexticnatural}
Let $X$ be very general with real multiplication by $\cO_{F_{0}}$, so that
$\NS(X)_{\QQ}=\{\theta_{t}=\sum_{j}\tau_{j}(t)\theta_{j}:t\in F_{0}\}$. Call
\emph{natural} the line bundles, the simple semi-homogeneous bundles of
\cite{Muk78}, the skyscraper sheaves of points, and the objects obtained from
these by shifts, duals, twists by line bundles, pull-back by automorphisms
commuting with $\cO_{F_{0}}$, and $\Phi_{\cP}$ followed by $\widehat X\cong X$.
On the $j$-th factor write classes in $\langle1,\theta_{j},\theta_{j}^{2}/2
\rangle$ as $(r,c,s)$, and put $\nu(t)=e^{t\theta_{j}}=(1,t,t^{2})$ and
$\nu(\infty)=(0,0,1)$. The character of a natural object is a multiple of
$\bigotimes_{j}\nu(\tau_{j}(t))$ for one $t\in F_{0}\cup\{\infty\}$, its
slope: a simple semi-homogeneous bundle has character $r\,e^{\theta_{t}}$
\cite{Muk78}, a point has slope $\infty$, and the operations act on slopes by
$t\mapsto-t$, $t+t'$, $tu^{2}$ with $u\in\cO_{F_{0}}^{\times}$, and $-1/t$,
the last by \Cref{thm:fmtheta} applied factor by factor. Then
\begin{enumerate}[label=\textup{(\alph*)},leftmargin=2.6em,itemsep=2pt]
\item no natural character lies in $S(0,q)$, for any totally positive $q$;
\item if a real combination of natural characters of $k$ pairwise distinct
  slopes is a nonzero element of $S(0,q)$, then $k\ge6$, and if $k=6$ it has
  $N_{w}=2$.
\end{enumerate}
So no cone of a morphism between two natural objects has a nonzero character
in $S(0,q)$, and by \Cref{lem:sexticweil}(iii) a class whose Orlov square has
a nonzero $F$-Weil part needs at least seven slopes. Six are attained:
$\operatorname{Re}e^{\sqrt{-1}\,\theta}=\frac52-\frac56(e^{\theta}+e^{-\theta})
+\frac1{12}(e^{2\theta}+e^{-2\theta})-10[\mathrm{pt}]$.

For (a): $S(0,q)$ is the tensor product of the planes
$P_{j}=\langle B_{j},\theta_{j}\rangle=\{s=-d_{j}r\}$, the form $c^{2}-rs$, for
which $\nu(t)$ and $\nu(\infty)$ are isotropic, is positive definite on
$P_{j}$, and a decomposable tensor lies in $\bigotimes_{j}P_{j}$ only if each
factor does. For (b), let $T=\sum_{i}a_{i}\bigotimes_{j}\nu(\tau_{j}(t_{i}))$
and $\lambda_{1}=s+d_{1}r$, which is positive on $\nu(t)$ and $\nu(\infty)$.
Applying $\lambda_{1}$ to the first factor gives a relation with nonzero
coefficients among the images, under the embedding of
$\PP^{1}\times\PP^{1}$ by forms of bidegree $(2,2)$, of the $k$ points
$(\tau_{2}(t_{i}),\tau_{3}(t_{i}))$, no two of which share a coordinate since
each $\tau_{j}$ is injective. Any five such points are independent, since the
product of the two vertical and the two horizontal lines through four of them
misses the fifth; so $k\ge6$. If $k=6$, a form of bidegree $(2,2)$ vanishing
at five of the points vanishes at the sixth; applied to $CD$, with $C$ the
curve of bidegree $(1,1)$ through three of the points and $D$ running over
the pencil of such curves through two others, whose base locus is those two
points, this puts the sixth point on $C$. So the six points lie on the graph
of a real M\"obius transformation, and likewise for the other pairs of
factors. Moving the second and third factors by these transformations, which
act through $\mathrm{GL}_{2}$ on $\langle1,\theta_{j},\theta_{j}^{2}/2\rangle$,
turns $T$ into a binary sextic, embedded in the tensor cube by
$l^{6}\mapsto(l^{2})^{\otimes3}$, that is annihilated by three definite
binary quadrics, the transforms of the $\lambda_{j}$. Two non-proportional
definite binary quadrics generate every binary form of degree at least three,
and so would annihilate the sextic; as $T\ne0$, the three quadrics are
proportional to one quadric $\Lambda$, and the sextic lies in the kernel of
contraction with $\Lambda$, spanned by $l_{+}^{6}$ and $l_{-}^{6}$ with
$l_{\pm}$ the isotropic lines of $\Lambda$. Moving back, $T$
is a combination of two conjugate decomposable tensors in $S(0,q)$, which are
$e_{\varepsilon}$ and $e_{-\varepsilon}$ for one $\varepsilon$; so
$N_{w}=2$.
\end{remark}
